% Options for packages loaded elsewhere
\PassOptionsToPackage{unicode}{hyperref}
\PassOptionsToPackage{hyphens}{url}
\PassOptionsToPackage{dvipsnames,svgnames,x11names}{xcolor}
\documentclass[
  spanish,
  11pt,
]{report}
\usepackage{xcolor}
\usepackage[margin=2.5cm,top=2.5cm,bottom=2.5cm]{geometry}
\usepackage{amsmath,amssymb}
\setcounter{secnumdepth}{5}
\usepackage{iftex}
\ifPDFTeX
  \usepackage[T1]{fontenc}
  \usepackage[utf8]{inputenc}
  \usepackage{textcomp} % provide euro and other symbols
\else % if luatex or xetex
  \usepackage{unicode-math} % this also loads fontspec
  \defaultfontfeatures{Scale=MatchLowercase}
  \defaultfontfeatures[\rmfamily]{Ligatures=TeX,Scale=1}
\fi
\usepackage{lmodern}
\ifPDFTeX\else
  % xetex/luatex font selection
  \setmainfont[]{Calibri}
  \setmonofont[]{Consolas}
\fi
% Use upquote if available, for straight quotes in verbatim environments
\IfFileExists{upquote.sty}{\usepackage{upquote}}{}
\IfFileExists{microtype.sty}{% use microtype if available
  \usepackage[]{microtype}
  \UseMicrotypeSet[protrusion]{basicmath} % disable protrusion for tt fonts
}{}
\makeatletter
\@ifundefined{KOMAClassName}{% if non-KOMA class
  \IfFileExists{parskip.sty}{%
    \usepackage{parskip}
  }{% else
    \setlength{\parindent}{0pt}
    \setlength{\parskip}{6pt plus 2pt minus 1pt}}
}{% if KOMA class
  \KOMAoptions{parskip=half}}
\makeatother
\usepackage{color}
\usepackage{fancyvrb}
\newcommand{\VerbBar}{|}
\newcommand{\VERB}{\Verb[commandchars=\\\{\}]}
\DefineVerbatimEnvironment{Highlighting}{Verbatim}{commandchars=\\\{\}}
% Add ',fontsize=\small' for more characters per line
\newenvironment{Shaded}{}{}
\newcommand{\AlertTok}[1]{\textcolor[rgb]{1.00,0.00,0.00}{\textbf{#1}}}
\newcommand{\AnnotationTok}[1]{\textcolor[rgb]{0.38,0.63,0.69}{\textbf{\textit{#1}}}}
\newcommand{\AttributeTok}[1]{\textcolor[rgb]{0.49,0.56,0.16}{#1}}
\newcommand{\BaseNTok}[1]{\textcolor[rgb]{0.25,0.63,0.44}{#1}}
\newcommand{\BuiltInTok}[1]{\textcolor[rgb]{0.00,0.50,0.00}{#1}}
\newcommand{\CharTok}[1]{\textcolor[rgb]{0.25,0.44,0.63}{#1}}
\newcommand{\CommentTok}[1]{\textcolor[rgb]{0.38,0.63,0.69}{\textit{#1}}}
\newcommand{\CommentVarTok}[1]{\textcolor[rgb]{0.38,0.63,0.69}{\textbf{\textit{#1}}}}
\newcommand{\ConstantTok}[1]{\textcolor[rgb]{0.53,0.00,0.00}{#1}}
\newcommand{\ControlFlowTok}[1]{\textcolor[rgb]{0.00,0.44,0.13}{\textbf{#1}}}
\newcommand{\DataTypeTok}[1]{\textcolor[rgb]{0.56,0.13,0.00}{#1}}
\newcommand{\DecValTok}[1]{\textcolor[rgb]{0.25,0.63,0.44}{#1}}
\newcommand{\DocumentationTok}[1]{\textcolor[rgb]{0.73,0.13,0.13}{\textit{#1}}}
\newcommand{\ErrorTok}[1]{\textcolor[rgb]{1.00,0.00,0.00}{\textbf{#1}}}
\newcommand{\ExtensionTok}[1]{#1}
\newcommand{\FloatTok}[1]{\textcolor[rgb]{0.25,0.63,0.44}{#1}}
\newcommand{\FunctionTok}[1]{\textcolor[rgb]{0.02,0.16,0.49}{#1}}
\newcommand{\ImportTok}[1]{\textcolor[rgb]{0.00,0.50,0.00}{\textbf{#1}}}
\newcommand{\InformationTok}[1]{\textcolor[rgb]{0.38,0.63,0.69}{\textbf{\textit{#1}}}}
\newcommand{\KeywordTok}[1]{\textcolor[rgb]{0.00,0.44,0.13}{\textbf{#1}}}
\newcommand{\NormalTok}[1]{#1}
\newcommand{\OperatorTok}[1]{\textcolor[rgb]{0.40,0.40,0.40}{#1}}
\newcommand{\OtherTok}[1]{\textcolor[rgb]{0.00,0.44,0.13}{#1}}
\newcommand{\PreprocessorTok}[1]{\textcolor[rgb]{0.74,0.48,0.00}{#1}}
\newcommand{\RegionMarkerTok}[1]{#1}
\newcommand{\SpecialCharTok}[1]{\textcolor[rgb]{0.25,0.44,0.63}{#1}}
\newcommand{\SpecialStringTok}[1]{\textcolor[rgb]{0.73,0.40,0.53}{#1}}
\newcommand{\StringTok}[1]{\textcolor[rgb]{0.25,0.44,0.63}{#1}}
\newcommand{\VariableTok}[1]{\textcolor[rgb]{0.10,0.09,0.49}{#1}}
\newcommand{\VerbatimStringTok}[1]{\textcolor[rgb]{0.25,0.44,0.63}{#1}}
\newcommand{\WarningTok}[1]{\textcolor[rgb]{0.38,0.63,0.69}{\textbf{\textit{#1}}}}
\usepackage{longtable,booktabs,array}
\newcounter{none} % for unnumbered tables
\usepackage{calc} % for calculating minipage widths
% Correct order of tables after \paragraph or \subparagraph
\usepackage{etoolbox}
\makeatletter
\patchcmd\longtable{\par}{\if@noskipsec\mbox{}\fi\par}{}{}
\makeatother
% Allow footnotes in longtable head/foot
\IfFileExists{footnotehyper.sty}{\usepackage{footnotehyper}}{\usepackage{footnote}}
\makesavenoteenv{longtable}
\ifLuaTeX
\usepackage[bidi=basic,shorthands=off]{babel}
\else
\usepackage[bidi=default,shorthands=off]{babel}
\fi
\ifPDFTeX
\else
\babelfont{rm}[]{Calibri}
\fi
\ifLuaTeX
  \usepackage{selnolig} % disable illegal ligatures
\fi
\setlength{\emergencystretch}{3em} % prevent overfull lines
\providecommand{\tightlist}{%
  \setlength{\itemsep}{0pt}\setlength{\parskip}{0pt}}
\usepackage{bookmark}
\IfFileExists{xurl.sty}{\usepackage{xurl}}{} % add URL line breaks if available
\urlstyle{same}
\hypersetup{
  pdftitle={Proyecto Final - Smart Parking Albacete},
  pdfauthor={alonso.marcos@alu.uclm.es},
  pdflang={es-ES},
  colorlinks=true,
  linkcolor={blue},
  filecolor={Maroon},
  citecolor={Blue},
  urlcolor={Blue},
  pdfcreator={LaTeX via pandoc}}

\title{Proyecto Final - Smart Parking Albacete}
\usepackage{etoolbox}
\makeatletter
\providecommand{\subtitle}[1]{% add subtitle to \maketitle
  \apptocmd{\@title}{\par {\large #1 \par}}{}{}
}
\makeatother
\subtitle{Internet de las Cosas y sus Aplicaciones (cod. 311482)}
\author{alonso.marcos@alu.uclm.es}
\date{Mayo 2026}

\begin{document}
\maketitle

{
\setcounter{tocdepth}{3}
\tableofcontents
}
`\newpage'

\newpage

\chapter{00. Resumen ejecutivo}\label{resumen-ejecutivo}

\section{0.1 Identificación del
proyecto}\label{identificaciuxf3n-del-proyecto}

{\def\LTcaptype{none} % do not increment counter
\begin{longtable}[]{@{}
  >{\raggedright\arraybackslash}p{(\linewidth - 2\tabcolsep) * \real{0.5000}}
  >{\raggedright\arraybackslash}p{(\linewidth - 2\tabcolsep) * \real{0.5000}}@{}}
\toprule\noalign{}
\begin{minipage}[b]{\linewidth}\raggedright
Campo
\end{minipage} & \begin{minipage}[b]{\linewidth}\raggedright
Valor
\end{minipage} \\
\midrule\noalign{}
\endhead
\bottomrule\noalign{}
\endlastfoot
Titulación & Máster Universitario en Big Data y Computación en la Nube
(UCLM) \\
Asignatura & Internet de las Cosas y sus Aplicaciones (cód. 311482, 6
ECTS) \\
Entregable & Proyecto final TR1 + prototipo \\
Curso & 2025-2026 \\
Autor & alonso.marcos@alu.uclm.es \\
Cliente ficticio & TECO S.L. \\
Cliente real (licitador) & Ayuntamiento de Albacete \\
Zona piloto & Entorno universitario de Albacete (BBOX
\texttt{38.976059,\ -1.858728\ →\ 38.983215,\ -1.846111}) \\
\end{longtable}
}

\section{0.2 Visión del proyecto}\label{visiuxf3n-del-proyecto}

El Ayuntamiento de Albacete saca a licitación una solución de
\textbf{smart parking} que permita monitorizar en tiempo real las plazas
de aparcamiento del entorno universitario, exponer su disponibilidad a
sistemas externos (aplicaciones de movilidad, vehículos autónomos vía
C-V2X, paneles informativos) y servir como base para un escalado
posterior a toda la ciudad.

TECO S.L. responde con una solución end-to-end basada en cuatro capas:

\begin{enumerate}
\def\labelenumi{\arabic{enumi}.}
\tightlist
\item
  \textbf{Telemetría}: sensores magnéticos AMR en cada plaza
  (independencia visual, bajo consumo) y cámaras ANPR puntuales en los
  accesos.
\item
  \textbf{Conectividad}: NB-IoT como red principal por su cobertura
  municipal y autonomía multianual; LoRaWAN como alternativa donde haya
  infraestructura privada disponible; fibra/5G para el backhaul de los
  gateways y cámaras.
\item
  \textbf{Plataforma cloud}: AWS IoT Core como puerta de entrada
  MQTT/TLS, IoT Rules + Lambda como capa de procesamiento, DynamoDB como
  almacén operacional, API Gateway como fachada pública.
\item
  \textbf{Aplicaciones}: API REST documentada con OpenAPI 3.0 (incluido
  formato GeoJSON para clientes cartográficos) y dashboard Streamlit
  para el operador, junto con la integración futura con FIWARE Orion
  (NGSI-v2) para interoperabilidad con otras smart cities.
\end{enumerate}

\section{0.3 Alcance del prototipo
entregado}\label{alcance-del-prototipo-entregado}

El prototipo cubre el \textbf{flujo extremo a extremo} desde sensor
hasta dashboard:

\begin{itemize}
\tightlist
\item
  40 plazas simuladas, distribuidas en 4 sub-zonas operativas reales
  sobre coordenadas exactas dentro de la BBOX (Campus UCLM, Estadio
  Carlos Belmonte, Hospital Universitario y entorno residencial sur).
\item
  Despliegue real en \textbf{AWS Academy Learner Lab} con scripts boto3
  idempotentes.
\item
  Verificación funcional: \texttt{GET\ /spots},
  \texttt{GET\ /spots/\{id\}}, \texttt{GET\ /zones},
  \texttt{GET\ /zones/\{id\}/kpis}, formato GeoJSON, dashboard con mapa,
  KPIs por sub-zona y serie temporal de ocupación.
\end{itemize}

\section{0.4 Decisiones técnicas más
relevantes}\label{decisiones-tuxe9cnicas-muxe1s-relevantes}

{\def\LTcaptype{none} % do not increment counter
\begin{longtable}[]{@{}
  >{\raggedright\arraybackslash}p{(\linewidth - 2\tabcolsep) * \real{0.4000}}
  >{\raggedright\arraybackslash}p{(\linewidth - 2\tabcolsep) * \real{0.6000}}@{}}
\toprule\noalign{}
\begin{minipage}[b]{\linewidth}\raggedright
Decisión
\end{minipage} & \begin{minipage}[b]{\linewidth}\raggedright
Justificación
\end{minipage} \\
\midrule\noalign{}
\endhead
\bottomrule\noalign{}
\endlastfoot
Sensor magnético AMR (no cámara por plaza) & Mejor balance
precisión/consumo/coste/privacidad; \textgreater5 años de autonomía. \\
NB-IoT como red principal & Cobertura LTE municipal existente, sin
desplegar infraestructura propia, coste OPEX bajo (\textless1
€/SIM/mes). \\
AWS IoT Core + IoT Rule + Lambda & Modelo serverless; coste por evento;
sin servidores que mantener. \\
DynamoDB on-demand & Sin capacidad provisionada; absorbe picos de
eventos sin throttling. \\
Lambda agregador (no Flink en piloto) & Volumen del piloto manejable en
una función simple; Flink se reserva para escenario ciudad. \\
Streamlit como dashboard & Velocidad de desarrollo y suficiente para una
demo de defensa. \\
FIWARE solo a nivel de diseño & El enunciado exige AWS; FIWARE se
documenta como capa de interoperabilidad NGSI-v2. \\
\end{longtable}
}

\section{0.5 Resultados clave}\label{resultados-clave}

\begin{itemize}
\tightlist
\item
  Latencia medida end-to-end (sensor → DynamoDB → API):
  \textbf{\textless{} 2 segundos} en condiciones del piloto.
\item
  Throughput verificado: 40 sensores publicando \textasciitilde4
  eventos/min sostenidos sin throttling.
\item
  KPIs por sub-zona generados cada vez que cambia el estado de una
  plaza, con persistencia como serie temporal.
\item
  Coste estimado del piloto en AWS: \textbf{\textless{} 5 USD/mes} para
  el volumen actual; escalable linealmente.
\item
  Coste estimado escenario ciudad (10 000 plazas): aproximadamente
  600-900 USD/mes en AWS más el CAPEX de sensores (\textasciitilde75
  €/plaza).
\end{itemize}

\section{0.6 Mapeo a los resultados de
aprendizaje}\label{mapeo-a-los-resultados-de-aprendizaje}

{\def\LTcaptype{none} % do not increment counter
\begin{longtable}[]{@{}
  >{\raggedright\arraybackslash}p{(\linewidth - 2\tabcolsep) * \real{0.2821}}
  >{\raggedright\arraybackslash}p{(\linewidth - 2\tabcolsep) * \real{0.7179}}@{}}
\toprule\noalign{}
\begin{minipage}[b]{\linewidth}\raggedright
Resultado
\end{minipage} & \begin{minipage}[b]{\linewidth}\raggedright
Cobertura en este proyecto
\end{minipage} \\
\midrule\noalign{}
\endhead
\bottomrule\noalign{}
\endlastfoot
\textbf{CN02} (arquitecturas de tratamiento masivo) & Capítulos 5, 7, 9
(cloud, modelo de datos, escalabilidad). \\
\textbf{HA03} (orquestación ETL, data lakes) & Capítulo 7 (S3 raw como
data lake bronze; DynamoDB como zona silver; pipeline en Lambda). \\
\textbf{CP02} (IoT, edge, streams) & Capítulos 3, 4, 5 (telemetría,
conectividad, edge) y 6 (Apache Flink en diseño). \\
\end{longtable}
}

\section{0.7 Estructura del documento}\label{estructura-del-documento}

\begin{verbatim}
00 Resumen ejecutivo (este documento)
01 Descripción del problema
02 Requisitos funcionales y no funcionales
03 Capa de telemetría
04 Conectividad
05 Arquitectura cloud en AWS
06 Integración con FIWARE y Apache Flink (diseño)
07 Modelo maestro de datos y APIs
08 Seguridad
09 Escalabilidad (piloto y ciudad)
10 Análisis de costes
11 Prototipo: ejecución, verificación y capturas
12 Limitaciones y trabajo futuro
13 Preparación de la defensa
\end{verbatim}

\newpage

\chapter{01. Descripción del
problema}\label{descripciuxf3n-del-problema}

\section{1.1 Contexto y motivación}\label{contexto-y-motivaciuxf3n}

El proyecto se inscribe en el marco de las iniciativas de \textbf{smart
city} que numerosos ayuntamientos españoles están impulsando como
respuesta al crecimiento del parque automovilístico, a los compromisos
de descarbonización derivados del Pacto Verde Europeo y a la inminente
llegada de la movilidad autónoma. La búsqueda manual de aparcamiento es
una de las patologías típicas del tráfico urbano: estudios clásicos
(Shoup, 2005) cifran entre el 15\% y el 30\% del tráfico de un centro
urbano el porcentaje atribuible a vehículos que circulan buscando plaza.
La consecuencia es triple: aumento del tiempo de viaje, incremento de
las emisiones y degradación del confort ciudadano.

El Ayuntamiento de Albacete, ciudad de unos 175 000 habitantes con un
Campus Universitario consolidado y un complejo sanitario-deportivo de
alta afluencia, plantea como zona piloto el entorno comprendido entre el
\textbf{Campus UCLM}, el \textbf{Estadio Municipal Carlos Belmonte}, el
\textbf{Hospital Universitario} y la zona residencial al sur de la
\textbf{AB-20}. El objetivo es disponer en tiempo real de la ocupación
de plazas y exponer esa información tanto a sistemas externos (apps,
paneles de mensajería variable, vehículos autónomos conectados vía
C-V2X) como a la propia plataforma municipal de gestión de movilidad.

\section{1.2 Cliente, actores y objetivos de
negocio}\label{cliente-actores-y-objetivos-de-negocio}

{\def\LTcaptype{none} % do not increment counter
\begin{longtable}[]{@{}
  >{\raggedright\arraybackslash}p{(\linewidth - 2\tabcolsep) * \real{0.5833}}
  >{\raggedright\arraybackslash}p{(\linewidth - 2\tabcolsep) * \real{0.4167}}@{}}
\toprule\noalign{}
\begin{minipage}[b]{\linewidth}\raggedright
Actor
\end{minipage} & \begin{minipage}[b]{\linewidth}\raggedright
Rol
\end{minipage} \\
\midrule\noalign{}
\endhead
\bottomrule\noalign{}
\endlastfoot
Ayuntamiento de Albacete & Cliente final; promotor de la licitación. \\
TECO S.L. & Adjudicatario (rol del autor); diseña, despliega y opera la
solución. \\
Conductores y ciudadanos & Usuarios indirectos vía app/panel. \\
Vehículos autónomos conectados & Consumidores futuros de la API mediante
C-V2X. \\
Concejalía de Movilidad & Cliente interno; explota KPIs para
planificación. \\
Servicios de emergencia (Hospital) & Caso de uso prioritario: acceso
rápido a plazas de servicio. \\
Operador externo de cámaras ANPR & Proporciona vídeo en los accesos. \\
\end{longtable}
}

Objetivos de negocio:

\begin{enumerate}
\def\labelenumi{\arabic{enumi}.}
\tightlist
\item
  Reducir tiempo medio de búsqueda de aparcamiento en la zona piloto.
\item
  Mejorar el acceso al Hospital y a la Facultad de Medicina en horario
  crítico.
\item
  Generar evidencia (KPIs) para planificar zonas reguladas (verde/azul)
  y políticas de movilidad.
\item
  Disponer de una plataforma con capacidad de escalado al resto de la
  ciudad sin rediseño.
\item
  Cumplir requisitos de interoperabilidad (NGSI-v2 / FIWARE) que el
  ecosistema de smart cities europeas demanda cada vez con más
  frecuencia en los pliegos.
\end{enumerate}

\section{1.3 Caracterización de la zona
piloto}\label{caracterizaciuxf3n-de-la-zona-piloto}

La zona piloto se ha delimitado mediante una \textbf{bounding box} real
definida por el responsable del proyecto:

{\def\LTcaptype{none} % do not increment counter
\begin{longtable}[]{@{}lll@{}}
\toprule\noalign{}
Esquina & Latitud & Longitud \\
\midrule\noalign{}
\endhead
\bottomrule\noalign{}
\endlastfoot
Suroeste & 38.976059 & -1.858728 \\
Noreste & 38.983215 & -1.846111 \\
\end{longtable}
}

Esta BBOX cubre una superficie aproximada de \textbf{0,55 km² (≈ 800 ×
800 m)} y engloba los siguientes polos generadores de demanda de
aparcamiento:

\begin{itemize}
\tightlist
\item
  \textbf{Z1-CAMPUS} (Universidad e investigación): Escuela Superior de
  Ingeniería Informática UCLM, Pabellón Universitario, edificios del
  campus a lo largo de Calle Imperial y Calle de la Navaja. Patrón de
  uso: pico fuerte de llegadas entre 7:30 y 10:00, segundo pico menor
  entre 16:00 y 18:00, vacío en fin de semana excepto eventos.
\item
  \textbf{Z2-DEPORTIVO} (eventos y ocio): Estadio Municipal Carlos
  Belmonte (capacidad \textasciitilde17 500), Campos de Fútbol ``Alba
  Redondo'' y ``José Copete'', restaurante Le Première. Patrón
  fuertemente correlacionado con el calendario deportivo; picos extremos
  en partidos (saturación), demanda baja entre semana.
\item
  \textbf{Z3-SANITARIO} (Hospital y facultades): Hospital Universitario
  de Albacete, Facultad de Medicina, Facultad de Farmacia. Patrón
  continuo 24/7 con turnos del personal sanitario y rotación de visitas;
  demanda crítica en urgencias.
\item
  \textbf{Z4-RESIDENCIAL} (zona sur de la AB-20): viviendas y servicios
  de la Avenida de la Mancha y Avenida Olimpia. Patrón residencial
  clásico: alta ocupación nocturna, liberación diurna parcial.
\end{itemize}

Los ejes viarios principales que concentran plazas de aparcamiento en
línea (zona blanca y futura zona regulada) son: Calle Imperial, Calle de
la Navaja, Calle Sancho Panza, Calle Duque de Rivas, Calle de la
Historia, Avenida del Arte, Calle San Juan, Calle La Química, Avenida de
la Mancha, Avenida Olimpia, Calle Maratón.

\section{1.4 Supuestos realistas
adoptados}\label{supuestos-realistas-adoptados}

Asunciones explícitas que el agente toma como base del diseño y del
análisis (consistentes con el espíritu del enunciado: ``Asuma y detalle,
de forma realista, cuantos parámetros necesite''):

{\def\LTcaptype{none} % do not increment counter
\begin{longtable}[]{@{}
  >{\raggedright\arraybackslash}p{(\linewidth - 4\tabcolsep) * \real{0.2439}}
  >{\raggedright\arraybackslash}p{(\linewidth - 4\tabcolsep) * \real{0.3902}}
  >{\raggedright\arraybackslash}p{(\linewidth - 4\tabcolsep) * \real{0.3659}}@{}}
\toprule\noalign{}
\begin{minipage}[b]{\linewidth}\raggedright
Supuesto
\end{minipage} & \begin{minipage}[b]{\linewidth}\raggedright
Valor adoptado
\end{minipage} & \begin{minipage}[b]{\linewidth}\raggedright
Justificación
\end{minipage} \\
\midrule\noalign{}
\endhead
\bottomrule\noalign{}
\endlastfoot
Número de plazas piloto & \textasciitilde500 & Coherente con el tamaño
de la BBOX y el patrón de aparcamiento en cordón típico. \\
Distribución por sub-zona & 125 plazas medias por zona & Equilibrio
entre simplicidad y representatividad. \\
Rotación media diaria & 4-8 ciclos/plaza Z3, 3-5 Z1, 1-3 Z4, picos
extremos en Z2 & Datos de campo de otras ciudades españolas
similares. \\
Cobertura NB-IoT & 100 \% en la BBOX & Albacete dispone de cobertura
LTE/NB-IoT completa de los operadores nacionales. \\
Disponibilidad mínima del sistema & 99,5 \% & Razonable para piloto;
objetivo 99,9 \% en versión productiva. \\
Latencia end-to-end máxima & 5 segundos & Suficiente para la guía humana
y compatible con consumo C-V2X periódico. \\
Tasa de envío por sensor & 1 evento/cambio + 1 heartbeat cada 5 min &
Reduce tráfico sin perder observabilidad. \\
Tamaño medio de mensaje & 0,5-2 KB JSON & Holgado para datos de plaza;
comprimible si fuera necesario. \\
Vida útil del sensor (batería) & 5-7 años & Específicación habitual de
los sensores AMR comerciales. \\
\end{longtable}
}

\section{1.5 Restricciones y consideraciones
específicas}\label{restricciones-y-consideraciones-especuxedficas}

\begin{itemize}
\tightlist
\item
  \textbf{Privacidad}: imposible (legalmente y socialmente) cubrir cada
  plaza con cámaras matriculeras. Las cámaras se restringen a accesos
  perimetrales y se justifican como detección agregada y para servicios
  complementarios (control de acceso a zonas restringidas).
\item
  \textbf{Despliegue urbano}: la instalación de sensores en calzada
  requiere conformidad municipal, planificación nocturna y, en zona de
  uso peatonal, alternativas de instalación bajo asfalto.
\item
  \textbf{Variabilidad meteorológica}: Albacete tiene veranos secos
  calurosos (hasta 40 °C) e inviernos fríos (heladas, nevadas
  puntuales). Los sensores deben certificarse IP67/IP68 y soportar el
  rango térmico.
\item
  \textbf{Cohabitación con otros servicios IoT}: la red municipal puede
  dar cabida en paralelo a contenedores inteligentes, iluminación,
  calidad del aire, etc. La arquitectura debe asumir que las plazas son
  sólo un dominio más.
\item
  \textbf{Continuidad operativa}: en una versión productiva la solución
  debe contemplar redundancia multi-AZ; el piloto se despliega en una
  sola región (us-east-1) por las limitaciones del AWS Academy Learner
  Lab.
\end{itemize}

\section{1.6 Riesgos identificados}\label{riesgos-identificados}

{\def\LTcaptype{none} % do not increment counter
\begin{longtable}[]{@{}
  >{\raggedright\arraybackslash}p{(\linewidth - 6\tabcolsep) * \real{0.1860}}
  >{\raggedright\arraybackslash}p{(\linewidth - 6\tabcolsep) * \real{0.3256}}
  >{\raggedright\arraybackslash}p{(\linewidth - 6\tabcolsep) * \real{0.2093}}
  >{\raggedright\arraybackslash}p{(\linewidth - 6\tabcolsep) * \real{0.2791}}@{}}
\toprule\noalign{}
\begin{minipage}[b]{\linewidth}\raggedright
Riesgo
\end{minipage} & \begin{minipage}[b]{\linewidth}\raggedright
Probabilidad
\end{minipage} & \begin{minipage}[b]{\linewidth}\raggedright
Impacto
\end{minipage} & \begin{minipage}[b]{\linewidth}\raggedright
Mitigación
\end{minipage} \\
\midrule\noalign{}
\endhead
\bottomrule\noalign{}
\endlastfoot
Falsos positivos del sensor magnético por motos cercanas & Media & Bajo
& Algoritmo de debounce en edge y umbral de confianza. \\
Pérdida de cobertura NB-IoT puntual & Baja & Medio & Buffer local en el
sensor y reintento exponencial. \\
Sensores robados/vandalizados & Baja & Bajo & Caja blindada; alarma de
manipulación; bajo coste unitario. \\
Saturación de la API en eventos deportivos & Media & Medio & API Gateway
con throttling configurable; caché ante terceros. \\
Fuga de datos personales & Muy baja & Alto & Política estricta: no se
almacena matrícula ni dato personal. \\
Caducidad credenciales lab (3h) & Alta (académico) & Alto & Scripts
idempotentes, teardown rápido, plan de relanzamiento. \\
\end{longtable}
}

\newpage

\chapter{02. Requisitos funcionales y no
funcionales}\label{requisitos-funcionales-y-no-funcionales}

La elicitación de requisitos parte del enunciado del campus virtual, se
completa con la guía operativa interna del proyecto y con la guía
docente de la asignatura, y se formaliza en este capítulo distinguiendo
entre \textbf{requisitos funcionales} (RF), \textbf{no funcionales}
(RNF) y \textbf{restricciones} (RST). Cada requisito incluye criterios
de aceptación verificables.

\section{2.1 Requisitos funcionales}\label{requisitos-funcionales}

{\def\LTcaptype{none} % do not increment counter
\begin{longtable}[]{@{}
  >{\raggedright\arraybackslash}p{(\linewidth - 4\tabcolsep) * \real{0.1026}}
  >{\raggedright\arraybackslash}p{(\linewidth - 4\tabcolsep) * \real{0.2821}}
  >{\raggedright\arraybackslash}p{(\linewidth - 4\tabcolsep) * \real{0.6154}}@{}}
\toprule\noalign{}
\begin{minipage}[b]{\linewidth}\raggedright
ID
\end{minipage} & \begin{minipage}[b]{\linewidth}\raggedright
Requisito
\end{minipage} & \begin{minipage}[b]{\linewidth}\raggedright
Criterio de aceptación
\end{minipage} \\
\midrule\noalign{}
\endhead
\bottomrule\noalign{}
\endlastfoot
RF-01 & El sistema debe detectar y reportar el estado (libre / ocupada)
de cada plaza monitorizada. & Cada cambio de estado se refleja en el
almacén operacional en ≤ 5 s. \\
RF-02 & Cada evento debe quedar asociado a la plaza concreta y a la
sub-zona. & Cada evento incluye \texttt{spotId} y \texttt{zoneId}. \\
RF-03 & El sistema debe transmitir eventos desde los sensores o desde un
simulador. & Prototipo: el simulador publica MQTT/TLS contra AWS IoT
Core. \\
RF-04 & El sistema debe mantener el estado actual de ocupación de cada
plaza. & Tabla DynamoDB \texttt{parking-state} con PK
\texttt{spotId}. \\
RF-05 & El sistema debe ofrecer un dashboard con la disponibilidad en
tiempo real. & Streamlit con KPIs, mapa y serie temporal; refresco ≤ 10
s. \\
RF-06 & El sistema debe exponer una API REST a sistemas externos. & API
Gateway + OpenAPI 3.0; endpoints \texttt{/spots}, \texttt{/zones},
\texttt{/zones/\{id\}/kpis}. \\
RF-07 & La API debe permitir análisis por sub-zona y por plaza. &
Parámetro \texttt{?zone=} y endpoint \texttt{/zones}. \\
RF-08 & La salida debe poder consumirse por mapas de terceros sin
transformación. & \texttt{?format=geojson} devuelve FeatureCollection
RFC 7946. \\
RF-09 & El sistema debe calcular agregados por sub-zona (libres,
ocupadas, \% ocupación). & Lambda agregador actualiza \texttt{zone-kpis}
por cambio de estado. \\
RF-10 & El sistema debe registrar histórico operacional. & Serie
temporal en \texttt{zone-kpis} indexada por
\texttt{(zoneId,\ windowEnd)}. \\
RF-11 & El sistema debe ser capaz de detectar sensores caídos. &
Heartbeat configurable + auditoría de \texttt{lastUpdated} (trabajo
futuro: alerta automática). \\
\end{longtable}
}

\section{2.2 Requisitos no funcionales}\label{requisitos-no-funcionales}

\subsection{Rendimiento y latencia}\label{rendimiento-y-latencia}

{\def\LTcaptype{none} % do not increment counter
\begin{longtable}[]{@{}
  >{\raggedright\arraybackslash}p{(\linewidth - 4\tabcolsep) * \real{0.1026}}
  >{\raggedright\arraybackslash}p{(\linewidth - 4\tabcolsep) * \real{0.2821}}
  >{\raggedright\arraybackslash}p{(\linewidth - 4\tabcolsep) * \real{0.6154}}@{}}
\toprule\noalign{}
\begin{minipage}[b]{\linewidth}\raggedright
ID
\end{minipage} & \begin{minipage}[b]{\linewidth}\raggedright
Requisito
\end{minipage} & \begin{minipage}[b]{\linewidth}\raggedright
Criterio de aceptación
\end{minipage} \\
\midrule\noalign{}
\endhead
\bottomrule\noalign{}
\endlastfoot
RNF-01 & Latencia E2E (sensor → consulta en API) ≤ 5 s en condiciones
normales. & Medición real en el prototipo: ≤ 2 s. \\
RNF-02 & La API debe responder en p95 ≤ 1 s. & Verificado con
curl/Streamlit. \\
RNF-03 & El sistema debe soportar 10 eventos/s sin throttling en el
piloto. & DynamoDB on-demand absorbe el pico. \\
\end{longtable}
}

\subsection{Escalabilidad}\label{escalabilidad}

{\def\LTcaptype{none} % do not increment counter
\begin{longtable}[]{@{}
  >{\raggedright\arraybackslash}p{(\linewidth - 4\tabcolsep) * \real{0.1026}}
  >{\raggedright\arraybackslash}p{(\linewidth - 4\tabcolsep) * \real{0.2821}}
  >{\raggedright\arraybackslash}p{(\linewidth - 4\tabcolsep) * \real{0.6154}}@{}}
\toprule\noalign{}
\begin{minipage}[b]{\linewidth}\raggedright
ID
\end{minipage} & \begin{minipage}[b]{\linewidth}\raggedright
Requisito
\end{minipage} & \begin{minipage}[b]{\linewidth}\raggedright
Criterio de aceptación
\end{minipage} \\
\midrule\noalign{}
\endhead
\bottomrule\noalign{}
\endlastfoot
RNF-04 & La arquitectura debe escalar de 500 plazas a ≥ 10 000 sin
rediseño. & Plan documentado en el capítulo 9 (sharding por zona, GSI
por estado, Timestream/Kinesis). \\
RNF-05 & Los servicios cloud usados deben ser serverless o gestionados.
& IoT Core, Lambda, DynamoDB, API Gateway son todos gestionados. \\
\end{longtable}
}

\subsection{Fiabilidad y
disponibilidad}\label{fiabilidad-y-disponibilidad}

{\def\LTcaptype{none} % do not increment counter
\begin{longtable}[]{@{}
  >{\raggedright\arraybackslash}p{(\linewidth - 4\tabcolsep) * \real{0.1026}}
  >{\raggedright\arraybackslash}p{(\linewidth - 4\tabcolsep) * \real{0.2821}}
  >{\raggedright\arraybackslash}p{(\linewidth - 4\tabcolsep) * \real{0.6154}}@{}}
\toprule\noalign{}
\begin{minipage}[b]{\linewidth}\raggedright
ID
\end{minipage} & \begin{minipage}[b]{\linewidth}\raggedright
Requisito
\end{minipage} & \begin{minipage}[b]{\linewidth}\raggedright
Criterio de aceptación
\end{minipage} \\
\midrule\noalign{}
\endhead
\bottomrule\noalign{}
\endlastfoot
RNF-06 & Disponibilidad ≥ 99,5 \% en piloto. & SLA AWS IoT Core 99,9 \%;
SLA Lambda 99,95 \%. \\
RNF-07 & El sistema debe tolerar pérdida temporal de conectividad del
sensor. & Reintento con backoff en el cliente MQTT; buffer local. \\
RNF-08 & Cualquier componente debe poder reiniciarse sin perder estado.
& Estado en DynamoDB; Lambdas sin estado en memoria. \\
\end{longtable}
}

\subsection{Seguridad}\label{seguridad}

{\def\LTcaptype{none} % do not increment counter
\begin{longtable}[]{@{}
  >{\raggedright\arraybackslash}p{(\linewidth - 4\tabcolsep) * \real{0.1026}}
  >{\raggedright\arraybackslash}p{(\linewidth - 4\tabcolsep) * \real{0.2821}}
  >{\raggedright\arraybackslash}p{(\linewidth - 4\tabcolsep) * \real{0.6154}}@{}}
\toprule\noalign{}
\begin{minipage}[b]{\linewidth}\raggedright
ID
\end{minipage} & \begin{minipage}[b]{\linewidth}\raggedright
Requisito
\end{minipage} & \begin{minipage}[b]{\linewidth}\raggedright
Criterio de aceptación
\end{minipage} \\
\midrule\noalign{}
\endhead
\bottomrule\noalign{}
\endlastfoot
RNF-09 & Toda comunicación dispositivo-cloud debe ir cifrada. & MQTT
sobre TLS 1.2 con mTLS (cert X.509). \\
RNF-10 & Cada dispositivo debe tener identidad propia. & Cada Thing en
IoT Core con su certificado y policy. \\
RNF-11 & La API pública debe poder limitar el consumo de terceros. & API
Gateway permite throttling y API keys. \\
RNF-12 & No deben almacenarse datos personales identificables. & El
modelo no contiene matrículas ni identificadores del conductor. \\
RNF-13 & Auditoría y trazabilidad de eventos. & S3 raw + CloudTrail +
CloudWatch Logs. \\
\end{longtable}
}

\subsection{Mantenibilidad y
observabilidad}\label{mantenibilidad-y-observabilidad}

{\def\LTcaptype{none} % do not increment counter
\begin{longtable}[]{@{}
  >{\raggedright\arraybackslash}p{(\linewidth - 4\tabcolsep) * \real{0.1026}}
  >{\raggedright\arraybackslash}p{(\linewidth - 4\tabcolsep) * \real{0.2821}}
  >{\raggedright\arraybackslash}p{(\linewidth - 4\tabcolsep) * \real{0.6154}}@{}}
\toprule\noalign{}
\begin{minipage}[b]{\linewidth}\raggedright
ID
\end{minipage} & \begin{minipage}[b]{\linewidth}\raggedright
Requisito
\end{minipage} & \begin{minipage}[b]{\linewidth}\raggedright
Criterio de aceptación
\end{minipage} \\
\midrule\noalign{}
\endhead
\bottomrule\noalign{}
\endlastfoot
RNF-14 & El despliegue debe ser reproducible desde código. & Scripts
boto3 idempotentes; ningún paso manual en consola. \\
RNF-15 & Debe existir un procedimiento limpio de borrado. &
\texttt{99\_teardown.py}. \\
RNF-16 & Métricas y logs deben centralizarse. & CloudWatch para Lambda e
IoT; logs de regla MQTT. \\
\end{longtable}
}

\subsection{Coste}\label{coste}

{\def\LTcaptype{none} % do not increment counter
\begin{longtable}[]{@{}
  >{\raggedright\arraybackslash}p{(\linewidth - 4\tabcolsep) * \real{0.1026}}
  >{\raggedright\arraybackslash}p{(\linewidth - 4\tabcolsep) * \real{0.2821}}
  >{\raggedright\arraybackslash}p{(\linewidth - 4\tabcolsep) * \real{0.6154}}@{}}
\toprule\noalign{}
\begin{minipage}[b]{\linewidth}\raggedright
ID
\end{minipage} & \begin{minipage}[b]{\linewidth}\raggedright
Requisito
\end{minipage} & \begin{minipage}[b]{\linewidth}\raggedright
Criterio de aceptación
\end{minipage} \\
\midrule\noalign{}
\endhead
\bottomrule\noalign{}
\endlastfoot
RNF-17 & El coste cloud del piloto debe ser \textless{} 10 USD/mes. &
Estimación detallada en el capítulo 10. \\
\end{longtable}
}

\subsection{Interoperabilidad y
privacidad}\label{interoperabilidad-y-privacidad}

{\def\LTcaptype{none} % do not increment counter
\begin{longtable}[]{@{}
  >{\raggedright\arraybackslash}p{(\linewidth - 4\tabcolsep) * \real{0.1026}}
  >{\raggedright\arraybackslash}p{(\linewidth - 4\tabcolsep) * \real{0.2821}}
  >{\raggedright\arraybackslash}p{(\linewidth - 4\tabcolsep) * \real{0.6154}}@{}}
\toprule\noalign{}
\begin{minipage}[b]{\linewidth}\raggedright
ID
\end{minipage} & \begin{minipage}[b]{\linewidth}\raggedright
Requisito
\end{minipage} & \begin{minipage}[b]{\linewidth}\raggedright
Criterio de aceptación
\end{minipage} \\
\midrule\noalign{}
\endhead
\bottomrule\noalign{}
\endlastfoot
RNF-18 & La API debe ser autodescriptiva y estándar. & OpenAPI 3.0.1,
GeoJSON RFC 7946. \\
RNF-19 & Diseño compatible con NGSI-v2 (FIWARE). & Modelo de entidades
documentado; mapping en el capítulo 6. \\
RNF-20 & Cumplimiento RGPD. & Sin datos personales; documentado en el
capítulo 8. \\
\end{longtable}
}

\section{2.3 Restricciones}\label{restricciones}

{\def\LTcaptype{none} % do not increment counter
\begin{longtable}[]{@{}
  >{\raggedright\arraybackslash}p{(\linewidth - 4\tabcolsep) * \real{0.1600}}
  >{\raggedright\arraybackslash}p{(\linewidth - 4\tabcolsep) * \real{0.5200}}
  >{\raggedright\arraybackslash}p{(\linewidth - 4\tabcolsep) * \real{0.3200}}@{}}
\toprule\noalign{}
\begin{minipage}[b]{\linewidth}\raggedright
ID
\end{minipage} & \begin{minipage}[b]{\linewidth}\raggedright
Restricción
\end{minipage} & \begin{minipage}[b]{\linewidth}\raggedright
Origen
\end{minipage} \\
\midrule\noalign{}
\endhead
\bottomrule\noalign{}
\endlastfoot
RST-01 & La plataforma cloud debe ser Amazon AWS. & Enunciado. \\
RST-02 & Las credenciales del lab caducan a las 3 horas. & AWS Academy
Learner Lab. \\
RST-03 & Solo se dispone del \texttt{LabRole} predefinido; no se pueden
crear roles IAM nuevos. & AWS Academy. \\
RST-04 & El prototipo debe demostrarse en defensa oral. & Guía
docente. \\
RST-05 & La documentación se entrega en castellano. & Guía docente. \\
RST-06 & El prototipo debe ejecutarse desde un equipo Windows estándar.
& Entorno del autor. \\
\end{longtable}
}

\section{2.4 Matriz de trazabilidad}\label{matriz-de-trazabilidad}

{\def\LTcaptype{none} % do not increment counter
\begin{longtable}[]{@{}
  >{\raggedright\arraybackslash}p{(\linewidth - 4\tabcolsep) * \real{0.1746}}
  >{\raggedright\arraybackslash}p{(\linewidth - 4\tabcolsep) * \real{0.4127}}
  >{\raggedright\arraybackslash}p{(\linewidth - 4\tabcolsep) * \real{0.4127}}@{}}
\toprule\noalign{}
\begin{minipage}[b]{\linewidth}\raggedright
Requisito
\end{minipage} & \begin{minipage}[b]{\linewidth}\raggedright
Capítulo donde se aborda
\end{minipage} & \begin{minipage}[b]{\linewidth}\raggedright
Componente del prototipo
\end{minipage} \\
\midrule\noalign{}
\endhead
\bottomrule\noalign{}
\endlastfoot
RF-01..04 & 03, 05 & Simulador + IoT Core + Lambda ingest + DynamoDB
state \\
RF-05 & 11 & \texttt{dashboard/streamlit\_app.py} \\
RF-06..08 & 07 & Lambda API + API Gateway + OpenAPI \\
RF-09..10 & 05, 07 & Lambda aggregator + DynamoDB kpis \\
RF-11 & 03, 12 & Heartbeat en simulador; alerta como trabajo futuro \\
RNF-01..03 & 05, 11 & Mediciones reales del prototipo \\
RNF-04..05 & 09 & Plan de escalabilidad \\
RNF-06..08 & 08, 09 & Reintentos + AZ + estado externalizado \\
RNF-09..13 & 08 & mTLS + LabRole + S3 + CloudWatch \\
RNF-14..16 & 11 & Scripts \texttt{infra/*.py} + CloudWatch \\
RNF-17 & 10 & Estimación pricing \\
RNF-18..20 & 06, 07, 08 & OpenAPI + FIWARE NGSI-v2 + ausencia de PII \\
\end{longtable}
}

\newpage

\chapter{03. Capa de telemetría}\label{capa-de-telemetruxeda}

La capa de telemetría es la base del sistema: si el sensor no detecta
correctamente la ocupación o no es operable a coste razonable, toda la
capa de datos posterior pierde valor. Este capítulo justifica la
combinación de sensores elegida tras una comparación explícita de
alternativas, atendiendo a precisión, coste, consumo energético,
robustez al entorno, dificultad de instalación, mantenimiento y
privacidad.

\section{3.1 Catálogo de tecnologías
evaluadas}\label{catuxe1logo-de-tecnologuxedas-evaluadas}

{\def\LTcaptype{none} % do not increment counter
\begin{longtable}[]{@{}
  >{\raggedright\arraybackslash}p{(\linewidth - 4\tabcolsep) * \real{0.1846}}
  >{\raggedright\arraybackslash}p{(\linewidth - 4\tabcolsep) * \real{0.2769}}
  >{\raggedright\arraybackslash}p{(\linewidth - 4\tabcolsep) * \real{0.5385}}@{}}
\toprule\noalign{}
\begin{minipage}[b]{\linewidth}\raggedright
Tecnología
\end{minipage} & \begin{minipage}[b]{\linewidth}\raggedright
Principio físico
\end{minipage} & \begin{minipage}[b]{\linewidth}\raggedright
Producto comercial de referencia
\end{minipage} \\
\midrule\noalign{}
\endhead
\bottomrule\noalign{}
\endlastfoot
Magnético AMR pasivo & Distorsión del campo magnético terrestre por la
masa metálica del vehículo & Nedap SENSIT (1-2), Smart Parking PS101,
Bosch INA213 \\
Ultrasonido & Tiempo de vuelo de un eco entre el sensor cenital y el
suelo/coche & Libelium Smart Parking Ultrasonic, Worldsensing
Fastprk-U \\
Radar mmWave & Detección de presencia / movimiento por reflexión
electromagnética & TI IWR6843, Vayyar (multi-plaza) \\
LIDAR 2D/3D & Telemetría láser; ideal para cubrir varias plazas a la vez
& Velodyne Puck (estático), Quanergy M8 \\
Cámara visión artificial (ANPR + ocupación) & Procesamiento de imagen
sobre stream de vídeo & Cámaras IP + algoritmo YOLO/EfficientDet en edge
o cloud \\
Sensores ambientales (CO₂, ruido, temperatura) & Complemento de la
calidad del entorno & Libelium Smart Cities, Adeunis Comfort \\
\end{longtable}
}

\section{3.2 Análisis comparativo}\label{anuxe1lisis-comparativo}

{\def\LTcaptype{none} % do not increment counter
\begin{longtable}[]{@{}
  >{\raggedright\arraybackslash}p{(\linewidth - 10\tabcolsep) * \real{0.1639}}
  >{\raggedright\arraybackslash}p{(\linewidth - 10\tabcolsep) * \real{0.1803}}
  >{\raggedright\arraybackslash}p{(\linewidth - 10\tabcolsep) * \real{0.2131}}
  >{\raggedright\arraybackslash}p{(\linewidth - 10\tabcolsep) * \real{0.1148}}
  >{\raggedright\arraybackslash}p{(\linewidth - 10\tabcolsep) * \real{0.1148}}
  >{\raggedright\arraybackslash}p{(\linewidth - 10\tabcolsep) * \real{0.2131}}@{}}
\toprule\noalign{}
\begin{minipage}[b]{\linewidth}\raggedright
Criterio
\end{minipage} & \begin{minipage}[b]{\linewidth}\raggedright
Magnético
\end{minipage} & \begin{minipage}[b]{\linewidth}\raggedright
Ultrasónico
\end{minipage} & \begin{minipage}[b]{\linewidth}\raggedright
Radar
\end{minipage} & \begin{minipage}[b]{\linewidth}\raggedright
LIDAR
\end{minipage} & \begin{minipage}[b]{\linewidth}\raggedright
Cámara ANPR
\end{minipage} \\
\midrule\noalign{}
\endhead
\bottomrule\noalign{}
\endlastfoot
Precisión detección ocupación & Alta (\textgreater95 \%) si bien
calibrado & Alta en interior; baja en exterior por lluvia/viento & Muy
alta; tolera meteorología & Muy alta & Muy alta + matrícula \\
Identificación de vehículo & No & No & No & No (forma) & Sí
(matrícula) \\
Coste unitario & 60-120 € & 80-150 € & 200-500 € & 1500-5000 € &
500-1200 € por cámara (multi-plaza) \\
Consumo energético & 5-10 µA medio (\textgreater5 años con pila D) &
0,5-1 mA medio (1-2 años con batería) & 5-50 mA (necesita red eléctrica)
& Decenas de W (red) & Red eléctrica obligatoria \\
Robustez frente a clima & Excelente (sellado IP68 bajo asfalto) & Media
(lluvia distorsiona) & Excelente & Buena (suciedad cristal) & Media
(lluvia, niebla, deslumbramiento) \\
Privacidad & Muy alta (sin imagen, sin matrícula) & Muy alta & Alta &
Alta (no identifica personas) & \textbf{Baja} (procesa matrículas) \\
Despliegue & Bajo asfalto (corte nocturno) & Soporte en farola/pared &
Mástil/pared & Mástil elevado & Mástil/edificio + iluminación \\
Mantenimiento & Bajo (cambio batería \textgreater5 años) & Medio
(limpieza, recalibración) & Bajo & Alto (mecánico) & Medio (limpieza
óptica) \\
Adecuado para parking en cordón & \textbf{Sí} & Limitado (necesita
techo) & Sí & Sí (multi-plaza) & Sí, agregado \\
Adecuado para escenario ciudad & Sí & Limitado & Sí (mayor coste) & No
(coste/mantenimiento) & Solo en accesos \\
\end{longtable}
}

\section{3.3 Decisión razonada}\label{decisiuxf3n-razonada}

Se opta por una \textbf{arquitectura híbrida} con dos categorías de
sensores:

\subsection{3.3.1 Sensor principal: magnético AMR bajo
asfalto}\label{sensor-principal-magnuxe9tico-amr-bajo-asfalto}

\begin{itemize}
\tightlist
\item
  \textbf{Una unidad por plaza}, sellada IP68, anclada en el centro de
  la plaza bajo el pavimento (instalación nocturna en una sola
  operación).
\item
  Calibración inicial in situ (toma del campo magnético terrestre de
  referencia sin coche).
\item
  Detección de paso desde estado libre→ocupado y viceversa con umbral
  configurable; debounce típico 2-3 s para evitar falsos positivos por
  motos en plazas contiguas.
\item
  Reporta el evento por LPWAN (capítulo 4) y mantiene un heartbeat
  configurable (5 min por defecto en producción; 30 s en demo).
\item
  Vida útil 5-7 años con pila D (justifica el TCO frente al ultrasonido
  o al radar).
\item
  Modelo de coste asumido: \textbf{75 € por unidad} (precio medio en
  compras municipales de 2024-2025).
\end{itemize}

Justificación frente a las alternativas: - Frente a
\textbf{ultrasonido}, gana en exterior por robustez climática y vida
útil. - Frente a \textbf{radar mmWave}, gana en consumo y coste; el
radar es muy buena opción para zonas de alta facturación o donde haya
alimentación eléctrica disponible (p.~ej. plazas de carga eléctrica). -
Frente a \textbf{LIDAR}, el coste y el mantenimiento descartan su uso
por plaza; LIDAR podría plantearse para cubrir varias plazas desde un
solo poste en versiones futuras. - Frente a \textbf{cámara}, ganan
privacidad, coste y simplicidad operativa; las cámaras se reservan para
los accesos.

\subsection{3.3.2 Sensor complementario: cámaras ANPR en accesos a
zonas}\label{sensor-complementario-cuxe1maras-anpr-en-accesos-a-zonas}

\begin{itemize}
\tightlist
\item
  \textbf{6-8 cámaras ANPR estratégicas} en los accesos a la BBOX
  (entradas/salidas de la N-430, C. San Juan, C. Imperial, Av. de la
  Mancha) para conteo agregado de flujo y, opcionalmente, control de
  acceso a futuras zonas de bajas emisiones.
\item
  No se usan para detectar la ocupación de plazas individuales: se
  elimina así el principal riesgo RGPD.
\item
  Se procesan en gateway edge con un modelo ligero (YOLOv8-tiny o
  similar), enviando solo metadatos (timestamp, dirección, recuento) a
  la cloud.
\item
  Resultado: conteo agregado por puerta de acceso que el operador puede
  correlacionar con la ocupación detectada por los sensores.
\end{itemize}

\subsection{3.3.3 Sensores ambientales
puntuales}\label{sensores-ambientales-puntuales}

\begin{itemize}
\tightlist
\item
  1 nodo ambiental cada \textasciitilde150 plazas (NO₂, PM2.5, CO₂,
  ruido, temperatura, humedad) para correlacionar ocupación con calidad
  del aire y aportar valor añadido al expediente municipal de movilidad.
\item
  Estos nodos no influyen en la decisión ``libre/ocupada''; aportan
  analítica complementaria.
\end{itemize}

\section{3.4 Modelo del nodo IoT (sensor
magnético)}\label{modelo-del-nodo-iot-sensor-magnuxe9tico}

Cada sensor se modela como un \textbf{objeto IoT (Thing)} en AWS IoT
Core con los siguientes atributos persistentes:

\begin{Shaded}
\begin{Highlighting}[]
\FunctionTok{\{}
  \DataTypeTok{"spotId"}\FunctionTok{:} \StringTok{"ALB{-}Z1{-}001"}\FunctionTok{,}
  \DataTypeTok{"zoneId"}\FunctionTok{:} \StringTok{"Z1{-}CAMPUS"}\FunctionTok{,}
  \DataTypeTok{"street"}\FunctionTok{:} \StringTok{"C. Imperial"}\FunctionTok{,}
  \DataTypeTok{"lat"}\FunctionTok{:} \FloatTok{38.97765}\FunctionTok{,}
  \DataTypeTok{"lon"}\FunctionTok{:} \FloatTok{{-}1.85745}
\FunctionTok{\}}
\end{Highlighting}
\end{Shaded}

El payload de cada evento publicado por el sensor (topic
\texttt{parking/\{zoneId\}/spot/\{spotId\}/status}):

\begin{Shaded}
\begin{Highlighting}[]
\FunctionTok{\{}
  \DataTypeTok{"spotId"}\FunctionTok{:} \StringTok{"ALB{-}Z1{-}001"}\FunctionTok{,}
  \DataTypeTok{"zoneId"}\FunctionTok{:} \StringTok{"Z1{-}CAMPUS"}\FunctionTok{,}
  \DataTypeTok{"street"}\FunctionTok{:} \StringTok{"C. Imperial"}\FunctionTok{,}
  \DataTypeTok{"lat"}\FunctionTok{:} \FloatTok{38.97765}\FunctionTok{,}
  \DataTypeTok{"lon"}\FunctionTok{:} \FloatTok{{-}1.85745}\FunctionTok{,}
  \DataTypeTok{"status"}\FunctionTok{:} \StringTok{"occupied"}\FunctionTok{,}
  \DataTypeTok{"batteryLevel"}\FunctionTok{:} \FloatTok{92.2}\FunctionTok{,}
  \DataTypeTok{"confidence"}\FunctionTok{:} \FloatTok{0.94}\FunctionTok{,}
  \DataTypeTok{"sensorType"}\FunctionTok{:} \StringTok{"magnetic"}\FunctionTok{,}
  \DataTypeTok{"timestamp"}\FunctionTok{:} \DecValTok{1778929200072}
\FunctionTok{\}}
\end{Highlighting}
\end{Shaded}

\section{3.5 Política de envío}\label{poluxedtica-de-envuxedo}

{\def\LTcaptype{none} % do not increment counter
\begin{longtable}[]{@{}
  >{\raggedright\arraybackslash}p{(\linewidth - 4\tabcolsep) * \real{0.2667}}
  >{\raggedright\arraybackslash}p{(\linewidth - 4\tabcolsep) * \real{0.3667}}
  >{\raggedright\arraybackslash}p{(\linewidth - 4\tabcolsep) * \real{0.3667}}@{}}
\toprule\noalign{}
\begin{minipage}[b]{\linewidth}\raggedright
Evento
\end{minipage} & \begin{minipage}[b]{\linewidth}\raggedright
Condición
\end{minipage} & \begin{minipage}[b]{\linewidth}\raggedright
Frecuencia
\end{minipage} \\
\midrule\noalign{}
\endhead
\bottomrule\noalign{}
\endlastfoot
Cambio de estado & \texttt{status} actual ≠ \texttt{status} anterior,
tras debounce & inmediato \\
Heartbeat & Sin cambio en \texttt{T} minutos & cada 5 min
(configurable) \\
Auto-test diario & Una vez al día & 1/día \\
Alarma de manipulación & Acelerómetro detecta movimiento del sensor &
inmediato \\
Aviso de batería baja & \texttt{batteryLevel\ \textless{}\ 20\ \%} & una
vez al día hasta sustitución \\
\end{longtable}
}

Este patrón minimiza el tráfico LPWAN (factor dominante del coste OPEX)
sin perder observabilidad.

\section{3.6 Calibración, instalación y
mantenimiento}\label{calibraciuxf3n-instalaciuxf3n-y-mantenimiento}

\begin{itemize}
\tightlist
\item
  \textbf{Calibración inicial} en taller (linealización del
  magnetómetro) y final in situ (medida del campo terrestre sin
  vehículo). Tiempo medio por unidad: 5-10 min.
\item
  \textbf{Instalación} nocturna por cuadrillas de 2 operarios;
  rendimiento típico 30-40 plazas por noche con maquinaria menor.
\item
  \textbf{Mantenimiento}: revisión por muestreo cada 6 meses;
  sustitución de pila prevista a los 5 años; firmware actualizable OTA
  vía LPWAN o mediante visita técnica si la pila lo permite.
\end{itemize}

\section{3.7 Tabla resumen de la
elección}\label{tabla-resumen-de-la-elecciuxf3n}

{\def\LTcaptype{none} % do not increment counter
\begin{longtable}[]{@{}
  >{\raggedright\arraybackslash}p{(\linewidth - 4\tabcolsep) * \real{0.2821}}
  >{\raggedright\arraybackslash}p{(\linewidth - 4\tabcolsep) * \real{0.4103}}
  >{\raggedright\arraybackslash}p{(\linewidth - 4\tabcolsep) * \real{0.3077}}@{}}
\toprule\noalign{}
\begin{minipage}[b]{\linewidth}\raggedright
Parámetro
\end{minipage} & \begin{minipage}[b]{\linewidth}\raggedright
Valor adoptado
\end{minipage} & \begin{minipage}[b]{\linewidth}\raggedright
Comentario
\end{minipage} \\
\midrule\noalign{}
\endhead
\bottomrule\noalign{}
\endlastfoot
Tecnología principal & Magnético AMR bajo asfalto & Mejor
coste/consumo/privacidad \\
Tecnología complementaria & Cámaras ANPR puntuales (accesos) & Solo
conteo agregado, sin matrículas en cloud \\
Tecnología auxiliar & Nodos ambientales 1/150 plazas & Valor añadido (no
crítico) \\
Densidad & 1 sensor / plaza & Modelo más fiable \\
Cadencia & 1 evento/cambio + heartbeat 5 min & Mínimo tráfico, máxima
observabilidad \\
Vida útil sensor & 5-7 años & Habitual del segmento \\
Coste medio sensor & 75 € & Precio de mercado actual \\
\end{longtable}
}

\newpage

\chapter{04. Conectividad y recolección de
datos}\label{conectividad-y-recolecciuxf3n-de-datos}

La conectividad es probablemente la decisión que más condiciona el TCO
del proyecto: una vez instalado, el sensor genera tráfico durante años y
la elección equivocada de red puede multiplicar el coste operativo por
10. Este capítulo describe la estrategia de comunicación adoptada,
compara las tecnologías candidatas y presenta un plan de despliegue
preliminar.

\section{4.1 Criterios de evaluación}\label{criterios-de-evaluaciuxf3n}

{\def\LTcaptype{none} % do not increment counter
\begin{longtable}[]{@{}
  >{\raggedright\arraybackslash}p{(\linewidth - 2\tabcolsep) * \real{0.3571}}
  >{\raggedright\arraybackslash}p{(\linewidth - 2\tabcolsep) * \real{0.6429}}@{}}
\toprule\noalign{}
\begin{minipage}[b]{\linewidth}\raggedright
Criterio
\end{minipage} & \begin{minipage}[b]{\linewidth}\raggedright
Por qué importa
\end{minipage} \\
\midrule\noalign{}
\endhead
\bottomrule\noalign{}
\endlastfoot
Latencia & Determina el objetivo de RNF-01 (≤ 5 s end-to-end). \\
Cobertura & Define si hace falta desplegar infraestructura propia
(gateways). \\
Coste de despliegue (CAPEX) & Inversión inicial: gateways, antenas,
alimentación. \\
Coste operativo (OPEX) & Suscripción SIM, mantenimiento del backhaul. \\
Consumo energético del nodo & Determina la vida útil de la batería del
sensor. \\
Disponibilidad de red & Cobertura real medida en la zona; SLA del
operador. \\
Escalabilidad & Capacidad de añadir miles de nodos sin saturar. \\
Idoneidad para V2X & Soporte futuro al diálogo con vehículos
conectados. \\
\end{longtable}
}

\section{4.2 Tecnologías candidatas y
análisis}\label{tecnologuxedas-candidatas-y-anuxe1lisis}

\subsection{4.2.1 NB-IoT (3GPP LPWAN sobre
LTE)}\label{nb-iot-3gpp-lpwan-sobre-lte}

\begin{itemize}
\tightlist
\item
  \textbf{Latencia típica}: 1,6-10 s (modo CP, half-duplex).
\item
  \textbf{Cobertura}: usa la infraestructura LTE existente; Albacete
  está cubierta por los tres operadores nacionales (Movistar, Vodafone,
  Orange).
\item
  \textbf{CAPEX}: 0 € (red del operador).
\item
  \textbf{OPEX}: 0,30-1,00 €/SIM/mes en contratos M2M de gran volumen.
\item
  \textbf{Consumo}: muy bajo (modo PSM y eDRX); permite \textgreater5
  años con pila D.
\item
  \textbf{Pros}: cobertura nacional, sin gateways propios, soporte
  certificado.
\item
  \textbf{Contras}: dependencia del operador; cuotas mínimas por SIM.
\end{itemize}

\subsection{4.2.2 LoRaWAN}\label{lorawan}

\begin{itemize}
\tightlist
\item
  \textbf{Latencia típica}: 1-2 s en clase A (depende del downlink).
\item
  \textbf{Cobertura}: red privada (gateways propios) o The Things
  Network (cobertura parcial).
\item
  \textbf{CAPEX}: \textasciitilde600-1500 € por gateway exterior; con
  3-4 gateways se cubre la BBOX.
\item
  \textbf{OPEX}: \textasciitilde0 (red privada); o cuota TTN Industries.
\item
  \textbf{Consumo}: muy bajo, comparable a NB-IoT.
\item
  \textbf{Pros}: control total, sin dependencia de operador; ideal en
  zonas industriales o municipales con red propia.
\item
  \textbf{Contras}: capacidad de downlink limitada; gestión de la red
  municipal recae en TECO S.L. o el ayuntamiento.
\end{itemize}

\subsection{4.2.3 Sigfox}\label{sigfox}

\begin{itemize}
\tightlist
\item
  Cobertura europea madura pero \textbf{incertidumbre comercial} tras la
  quiebra de la matriz en 2022 y la reorganización posterior. Se
  descarta por riesgo de continuidad.
\end{itemize}

\subsection{4.2.4 Wi-Fi / Ethernet}\label{wi-fi-ethernet}

\begin{itemize}
\tightlist
\item
  \textbf{Latencia}: \textless100 ms.
\item
  \textbf{Cobertura}: limitada al edificio o radio del AP; no aplicable
  a plazas exteriores en cordón.
\item
  \textbf{Consumo}: alto; incompatible con el modelo de pila.
\item
  \textbf{Uso adecuado}: gateways edge, cámaras, paneles de mensajería
  variable, no sensores.
\end{itemize}

\subsection{4.2.5 5G NR}\label{g-nr}

\begin{itemize}
\tightlist
\item
  \textbf{Latencia}: \textless10 ms en URLLC; 30-50 ms en eMBB.
\item
  \textbf{Cobertura}: en expansión; en zona piloto disponible
  parcialmente.
\item
  \textbf{Consumo y coste por sensor}: aún elevados frente a NB-IoT.
\item
  \textbf{Uso adecuado}: backhaul de gateways edge, cámaras ANPR,
  conexión a vehículo conectado.
\end{itemize}

\subsection{4.2.6 C-V2X (Cellular V2X)}\label{c-v2x-cellular-v2x}

\begin{itemize}
\tightlist
\item
  Tecnología específica para comunicación vehículo-infraestructura sobre
  LTE/5G.
\item
  En este proyecto se considera \textbf{como interfaz de salida} hacia
  los vehículos autónomos, no como red de subida del sensor. El sistema
  expone los datos vía API REST y, si el operador del vehículo dispone
  de OBU C-V2X, puede consumir esa misma información a través de un
  broker MEC.
\end{itemize}

\section{4.3 Decisión razonada}\label{decisiuxf3n-razonada-1}

{\def\LTcaptype{none} % do not increment counter
\begin{longtable}[]{@{}
  >{\raggedright\arraybackslash}p{(\linewidth - 4\tabcolsep) * \real{0.2308}}
  >{\raggedright\arraybackslash}p{(\linewidth - 4\tabcolsep) * \real{0.4231}}
  >{\raggedright\arraybackslash}p{(\linewidth - 4\tabcolsep) * \real{0.3462}}@{}}
\toprule\noalign{}
\begin{minipage}[b]{\linewidth}\raggedright
Capa
\end{minipage} & \begin{minipage}[b]{\linewidth}\raggedright
Tecnología
\end{minipage} & \begin{minipage}[b]{\linewidth}\raggedright
Razones
\end{minipage} \\
\midrule\noalign{}
\endhead
\bottomrule\noalign{}
\endlastfoot
Sensor → Cloud & \textbf{NB-IoT} & Cobertura municipal sin desplegar
infraestructura; OPEX bajo; consumo compatible con vida útil
multianual. \\
Plan B / red privada & \textbf{LoRaWAN} & En zonas donde el ayuntamiento
ya disponga de gateways propios, se prefiere por independencia y coste
cero por mensaje. \\
Backhaul gateway / cámara & \textbf{Ethernet / fibra / 5G} & Necesario
para volumen y alimentación eléctrica. \\
Comunicación con vehículo autónomo & \textbf{API REST (HTTPS) + C-V2X
(futuro)} & Estándar y desacoplado; permite a cualquier OEM consumir. \\
\end{longtable}
}

Esta combinación es \textbf{defendible y realista}: refleja la práctica
habitual de los proyectos de smart parking municipales actualmente en
operación en España (Santander, Málaga, Pontevedra) donde el grueso de
la flota va sobre NB-IoT/LoRaWAN y los puntos críticos sobre fibra/5G.

\section{4.4 Plan de despliegue
preliminar}\label{plan-de-despliegue-preliminar}

\subsection{Fase 0 -- Diseño detallado y permisos (semanas
1-4)}\label{fase-0-diseuxf1o-detallado-y-permisos-semanas-1-4}

\begin{itemize}
\tightlist
\item
  Mapa exacto de plazas a sensorizar (≈ 500) con coordenadas,
  dimensiones y prioridades.
\item
  Coordinación con la Concejalía de Movilidad y la EMT para cortes
  nocturnos.
\item
  Acuerdo con operador NB-IoT (tarifa SIM M2M, ventana de provisión, APN
  privado).
\item
  Acuerdo con el responsable de la red WAN municipal para el backhaul de
  los gateways y cámaras.
\end{itemize}

\subsection{Fase 1 -- Piloto técnico (semanas
5-10)}\label{fase-1-piloto-tuxe9cnico-semanas-5-10}

\begin{itemize}
\tightlist
\item
  Despliegue de los \textbf{3 gateways edge} (Z1, Z2/Z3, Z4) sobre
  farolas o mobiliario urbano municipal, alimentados por la red
  eléctrica de alumbrado público.
\item
  Instalación de \textbf{20 sensores piloto} en Z1-CAMPUS para validar
  señal, consumo y patrones.
\item
  Despliegue del backend AWS (esta memoria).
\item
  Validación operativa durante 2-3 semanas.
\end{itemize}

\subsection{Fase 2 -- Despliegue completo del piloto (semanas
11-20)}\label{fase-2-despliegue-completo-del-piloto-semanas-11-20}

\begin{itemize}
\tightlist
\item
  Instalación nocturna de los \textasciitilde500 sensores restantes
  (cuadrillas de 2 personas, \textasciitilde35 sensores/noche).
\item
  6-8 cámaras ANPR en los accesos.
\item
  4 nodos ambientales puntuales.
\item
  Onboarding masivo en IoT Core mediante AWS IoT Provisioning Templates.
\end{itemize}

\subsection{Fase 3 -- Operación y métricas (mes
6-12)}\label{fase-3-operaciuxf3n-y-muxe9tricas-mes-6-12}

\begin{itemize}
\tightlist
\item
  KPIs de servicio: \% uptime, latencia E2E, tasa de falsos positivos,
  \% de baterías sustituidas.
\item
  Iteración del modelo de Lambda agregador en función de los patrones
  observados.
\item
  Integración con el sistema municipal de paneles de mensajería
  variable.
\end{itemize}

\subsection{Fase 4 -- Escalado a ciudad (año
2)}\label{fase-4-escalado-a-ciudad-auxf1o-2}

\begin{itemize}
\tightlist
\item
  Onboarding por barrios; despliegue de gateways adicionales si LoRaWAN
  es la red elegida.
\item
  Integración con C-V2X (RSU / MEC) si el ayuntamiento o un OEM lo
  solicitan.
\item
  Migración del histórico operacional a Amazon Timestream o S3 + Athena.
\end{itemize}

\section{4.5 Estimación de tráfico}\label{estimaciuxf3n-de-truxe1fico}

Asumiendo el escenario de operación normal:

\begin{itemize}
\tightlist
\item
  500 plazas × 5 cambios/día medios = 2500 eventos de cambio/día.
\item
  500 plazas × 288 heartbeats/día (cada 5 min) = 144 000 eventos/día.
\item
  Tamaño medio del payload: 0,5-2 KB.
\item
  Volumen diario: \textasciitilde150 000 mensajes / \textasciitilde150
  MB / día.
\item
  Volumen pico (evento deportivo): hasta 1000 cambios/min durante 15
  min.
\end{itemize}

Estos números son perfectamente asumibles tanto por NB-IoT (canal por
sector LTE) como por la combinación AWS IoT Core + Lambda + DynamoDB
on-demand, según se detalla en los capítulos 9 y 10.

\section{4.6 Diagrama de red}\label{diagrama-de-red}

\begin{Shaded}
\begin{Highlighting}[]
\NormalTok{flowchart LR}
\NormalTok{    subgraph Calle}
\NormalTok{        S1[Sensor magnetico]}
\NormalTok{        S2[Sensor magnetico]}
\NormalTok{        CAM[Camara ANPR acceso]}
\NormalTok{        ENV[Nodo ambiental]}
\NormalTok{    end}

\NormalTok{    subgraph Backhaul}
\NormalTok{        GW[Gateway edge {-} opcional LoRa o ANPR]}
\NormalTok{        SIM[Red NB{-}IoT operador]}
\NormalTok{    end}

\NormalTok{    subgraph Cloud}
\NormalTok{        IoT[AWS IoT Core]}
\NormalTok{    end}

\NormalTok{    subgraph Vehiculo}
\NormalTok{        V2X[OBU C{-}V2X / API]}
\NormalTok{    end}

\NormalTok{    S1 {-}{-}\textgreater{}|NB{-}IoT MQTT/TLS| SIM}
\NormalTok{    S2 {-}{-}\textgreater{}|NB{-}IoT MQTT/TLS| SIM}
\NormalTok{    CAM {-}{-}\textgreater{}|Ethernet| GW}
\NormalTok{    ENV {-}{-}\textgreater{}|NB{-}IoT| SIM}
\NormalTok{    SIM {-}{-}\textgreater{}|Internet TLS| IoT}
\NormalTok{    GW {-}{-}\textgreater{}|Fibra/5G TLS| IoT}
\NormalTok{    IoT {-}{-}\textgreater{}|REST HTTPS| V2X}
\end{Highlighting}
\end{Shaded}

\newpage

\chapter{05. Arquitectura cloud en AWS}\label{arquitectura-cloud-en-aws}

Este capítulo describe la arquitectura cloud completa desplegada en AWS
para el piloto. Se justifican los servicios elegidos, se enseña el flujo
de información extremo a extremo y se detalla qué responsabilidad asume
cada nivel (edge vs cloud).

\section{5.1 Servicios AWS utilizados}\label{servicios-aws-utilizados}

{\def\LTcaptype{none} % do not increment counter
\begin{longtable}[]{@{}
  >{\raggedright\arraybackslash}p{(\linewidth - 4\tabcolsep) * \real{0.2128}}
  >{\raggedright\arraybackslash}p{(\linewidth - 4\tabcolsep) * \real{0.4468}}
  >{\raggedright\arraybackslash}p{(\linewidth - 4\tabcolsep) * \real{0.3404}}@{}}
\toprule\noalign{}
\begin{minipage}[b]{\linewidth}\raggedright
Servicio
\end{minipage} & \begin{minipage}[b]{\linewidth}\raggedright
Rol en la solución
\end{minipage} & \begin{minipage}[b]{\linewidth}\raggedright
Justificación
\end{minipage} \\
\midrule\noalign{}
\endhead
\bottomrule\noalign{}
\endlastfoot
\textbf{AWS IoT Core} & Broker MQTT/TLS, registro de Things,
certificados, IoT Rules. & Servicio gestionado, integración nativa con
resto de AWS, mTLS estándar industrial. \\
\textbf{AWS IoT Rule} & Filtrado SQL del topic y enrutado a Lambda. &
Sin servidores; precio por evento; permite múltiples acciones. \\
\textbf{AWS Lambda} & Funciones de ingesta, agregación y API. &
Serverless, paga por uso, integración directa con IoT, DynamoDB y API
Gateway. \\
\textbf{Amazon DynamoDB} & Estado actual + serie temporal de KPIs. &
Latencia ms; on-demand absorbe picos; modelado clave-valor sencillo. \\
\textbf{Amazon API Gateway} & Fachada REST hacia terceros y dashboard. &
TLS, throttling, cuotas, autenticación, integración nativa con
Lambda. \\
\textbf{Amazon S3} (futuro/diseño) & Data lake bronze (raw events). &
Persistencia barata; integrable con Athena/Glue. \\
\textbf{AWS IoT Greengrass} (diseño) & Edge computing en gateway de
zona. & Para preprocesado y operación local sin conectividad. \\
\textbf{Amazon CloudWatch} & Logs y métricas de las Lambdas y reglas
IoT. & Incluido por defecto, observabilidad inmediata. \\
\textbf{IAM (LabRole)} & Identidad de las Lambdas. & Restricción del AWS
Academy: no se pueden crear roles propios. \\
\end{longtable}
}

Servicios \textbf{diseñados pero no desplegados} en el prototipo (por
restricciones del lab o por enfoque): IoT Greengrass, Timestream,
Kinesis Data Streams, Cognito, QuickSight.

\section{5.2 Diagrama de arquitectura
(general)}\label{diagrama-de-arquitectura-general}

\begin{Shaded}
\begin{Highlighting}[]
\NormalTok{flowchart LR}
\NormalTok{  subgraph Borde fisico}
\NormalTok{    SENSOR[Sensor magnetico AMR\textless{}br/\textgreater{}NB{-}IoT MQTT/TLS]}
\NormalTok{    GW[Gateway edge\textless{}br/\textgreater{}IoT Greengrass{-}diseno]}
\NormalTok{    CAM[Camara ANPR]}
\NormalTok{  end}

\NormalTok{  subgraph AWS\_Cloud}
\NormalTok{    IOTC[AWS IoT Core\textless{}br/\textgreater{}Message Broker + Registry]}
\NormalTok{    RULE\{IoT Rule\textless{}br/\textgreater{}SQL select * from parking/+/spot/+/status\}}
\NormalTok{    L1[Lambda ingest]}
\NormalTok{    L2[Lambda aggregator]}
\NormalTok{    L3[Lambda api]}
\NormalTok{    DDB1[(DynamoDB\textless{}br/\textgreater{}parking{-}state)]}
\NormalTok{    DDB2[(DynamoDB\textless{}br/\textgreater{}zone{-}kpis)]}
\NormalTok{    S3[(S3 raw {-} diseno)]}
\NormalTok{    APIGW[API Gateway REST]}
\NormalTok{    CW[CloudWatch Logs]}
\NormalTok{  end}

\NormalTok{  subgraph Consumidores}
\NormalTok{    DASH[Streamlit dashboard]}
\NormalTok{    TER[Vehiculos / apps / paneles MV]}
\NormalTok{  end}

\NormalTok{  SENSOR {-}{-}\textgreater{}|MQTT/TLS| IOTC}
\NormalTok{  GW {-}{-}\textgreater{} IOTC}
\NormalTok{  CAM {-}{-}\textgreater{} GW}
\NormalTok{  IOTC {-}{-}\textgreater{} RULE}
\NormalTok{  RULE {-}{-}\textgreater{} L1}
\NormalTok{  RULE {-}.{-}\textgreater{} S3}
\NormalTok{  L1 {-}{-}\textgreater{} DDB1}
\NormalTok{  L1 {-}{-} invoke {-}{-}\textgreater{} L2}
\NormalTok{  L2 {-}{-}\textgreater{} DDB2}
\NormalTok{  L1 {-}{-}\textgreater{} CW}
\NormalTok{  L2 {-}{-}\textgreater{} CW}
\NormalTok{  L3 {-}{-}\textgreater{} CW}
\NormalTok{  APIGW {-}{-}\textgreater{} L3}
\NormalTok{  L3 {-}{-}\textgreater{} DDB1}
\NormalTok{  L3 {-}{-}\textgreater{} DDB2}
\NormalTok{  DASH {-}{-}\textgreater{}|HTTPS| APIGW}
\NormalTok{  TER {-}{-}\textgreater{}|HTTPS| APIGW}
\end{Highlighting}
\end{Shaded}

\section{5.3 Flujo de información extremo a
extremo}\label{flujo-de-informaciuxf3n-extremo-a-extremo}

\begin{enumerate}
\def\labelenumi{\arabic{enumi}.}
\tightlist
\item
  \textbf{Sensor}: detecta cambio de estado, aplica debounce y publica
  un mensaje JSON en \texttt{parking/\{zoneId\}/spot/\{spotId\}/status}
  mediante MQTT/TLS contra el endpoint \texttt{iot:Data-ATS} de IoT
  Core.
\item
  \textbf{IoT Core}: autentica mediante mTLS (certificado X.509), valida
  la policy adjunta y entrega el mensaje al broker.
\item
  \textbf{IoT Rule}: la regla
  \texttt{smart\_parking\_albacete\_ingest\_rule} filtra todos los
  topics que casan con \texttt{parking/+/spot/+/status} mediante el SQL
  \texttt{SELECT\ *\ FROM\ \textquotesingle{}parking/+/spot/+/status\textquotesingle{}}
  y dispara una acción Lambda.
\item
  \textbf{Lambda ingest} (\texttt{smart-parking-albacete-ingest}):
  normaliza el payload, hace UPSERT en \texttt{parking-state}, y si el
  estado cambia respecto al anterior, invoca de forma asíncrona a la
  Lambda agregador con el \texttt{zoneId}.
\item
  \textbf{Lambda aggregator}
  (\texttt{smart-parking-albacete-aggregator}): hace un scan filtrado
  por zona en \texttt{parking-state}, calcula libres/ocupadas/ratio y
  escribe una fila en \texttt{zone-kpis} con la marca temporal
  \texttt{windowEnd}.
\item
  \textbf{API Gateway} expone los endpoints REST (\texttt{/spots},
  \texttt{/spots/\{id\}}, \texttt{/zones}, \texttt{/zones/\{id\}/kpis})
  y los integra con la Lambda \texttt{smart-parking-albacete-api}, que
  consulta DynamoDB y formatea la respuesta (JSON o GeoJSON).
\item
  \textbf{Cliente} (dashboard Streamlit o un tercero) consume la API por
  HTTPS.
\item
  \textbf{CloudWatch} recoge automáticamente logs y métricas de cada
  Lambda.
\end{enumerate}

\section{5.4 Modelado de los nodos en IoT
Core}\label{modelado-de-los-nodos-en-iot-core}

Cada plaza se modela como un \textbf{Thing} dentro de un \textbf{Thing
Type} común (\texttt{smart-parking-albacete-sensor-type}) y de un
\textbf{Thing Group} (\texttt{smart-parking-albacete-fleet}) que
facilita acciones masivas (despliegue OTA, búsqueda por atributos).

Cada Thing lleva como atributos buscables
(\texttt{searchableAttributes}):

\begin{itemize}
\tightlist
\item
  \texttt{zoneId}: sub-zona operacional (Z1-CAMPUS, Z2-DEPORTIVO,
  Z3-SANITARIO, Z4-RESIDENCIAL).
\item
  \texttt{street}: nombre de calle (sanitizado para cumplir las
  restricciones de IoT).
\item
  \texttt{lat}, \texttt{lon}: coordenadas geográficas.
\end{itemize}

Todos los Things del piloto comparten un único \textbf{certificado
X.509} y una única \textbf{policy}
(\texttt{smart-parking-albacete-sensor-policy}). En producción se
generaría un certificado por dispositivo para revocaciones granulares;
en el piloto se prioriza simplicidad operativa.

La policy autoriza únicamente las acciones MQTT necesarias y solo en los
topics del proyecto:

\begin{Shaded}
\begin{Highlighting}[]
\FunctionTok{\{}
  \DataTypeTok{"Version"}\FunctionTok{:} \StringTok{"2012{-}10{-}17"}\FunctionTok{,}
  \DataTypeTok{"Statement"}\FunctionTok{:} \OtherTok{[}
    \FunctionTok{\{}\DataTypeTok{"Effect"}\FunctionTok{:} \StringTok{"Allow"}\FunctionTok{,} \DataTypeTok{"Action"}\FunctionTok{:} \StringTok{"iot:Connect"}\FunctionTok{,} \DataTypeTok{"Resource"}\FunctionTok{:} \StringTok{"*"}\FunctionTok{\}}\OtherTok{,}
    \FunctionTok{\{}\DataTypeTok{"Effect"}\FunctionTok{:} \StringTok{"Allow"}\FunctionTok{,} \DataTypeTok{"Action"}\FunctionTok{:} \StringTok{"iot:Publish"}\FunctionTok{,}   \DataTypeTok{"Resource"}\FunctionTok{:} \StringTok{"arn:aws:iot:*:*:topic/parking/*"}\FunctionTok{\}}\OtherTok{,}
    \FunctionTok{\{}\DataTypeTok{"Effect"}\FunctionTok{:} \StringTok{"Allow"}\FunctionTok{,} \DataTypeTok{"Action"}\FunctionTok{:} \StringTok{"iot:Subscribe"}\FunctionTok{,} \DataTypeTok{"Resource"}\FunctionTok{:} \StringTok{"arn:aws:iot:*:*:topicfilter/parking/*"}\FunctionTok{\}}\OtherTok{,}
    \FunctionTok{\{}\DataTypeTok{"Effect"}\FunctionTok{:} \StringTok{"Allow"}\FunctionTok{,} \DataTypeTok{"Action"}\FunctionTok{:} \StringTok{"iot:Receive"}\FunctionTok{,}   \DataTypeTok{"Resource"}\FunctionTok{:} \StringTok{"arn:aws:iot:*:*:topic/parking/*"}\FunctionTok{\}}
  \OtherTok{]}
\FunctionTok{\}}
\end{Highlighting}
\end{Shaded}

\section{5.5 Capa edge}\label{capa-edge}

La guía del proyecto distingue claramente entre procesamiento \textbf{en
edge} y \textbf{en cloud}. En el diseño:

\subsection{En edge (Gateway con AWS IoT Greengrass v2 --
diseño)}\label{en-edge-gateway-con-aws-iot-greengrass-v2-diseuxf1o}

\begin{itemize}
\tightlist
\item
  Recogida de los sensores que no hablan NB-IoT directamente (LoRaWAN,
  Bluetooth de servicio).
\item
  Filtrado de ruido (descartar lecturas espurias del magnetómetro).
\item
  \textbf{Debounce y detección de cambios reales}: evita que un coche
  que apaga el motor y arranca rápido genere dos eventos.
\item
  \textbf{Caché temporal} si la conexión a AWS está caída: hasta 24
  horas de buffer con reenvío al recuperarse.
\item
  \textbf{Agregación simple por zona} para reducir tráfico (por ejemplo,
  enviar ``8 plazas libres en Z2-DEPORTIVO'' como evento agregado en vez
  de los 8 individuales en intervalos cortos).
\item
  \textbf{Procesamiento de cámaras ANPR}: ejecución local de un modelo
  ligero que solo envía a la cloud un evento con dirección y matrícula
  hash (sin texto plano).
\end{itemize}

\subsection{En cloud (lo implementado
realmente)}\label{en-cloud-lo-implementado-realmente}

\begin{itemize}
\tightlist
\item
  \textbf{Normalización} del payload (Lambda ingest).
\item
  \textbf{Persistencia del estado actual} (DynamoDB
  \texttt{parking-state}).
\item
  \textbf{Cálculo de agregados} (Lambda aggregator → DynamoDB
  \texttt{zone-kpis}).
\item
  \textbf{Histórico operacional} como serie temporal.
\item
  \textbf{APIs} internas y externas.
\item
  \textbf{Seguridad} (mTLS, IAM, RGPD).
\item
  \textbf{Observabilidad} (CloudWatch).
\item
  \textbf{Analítica avanzada} (S3 raw + Athena en escenario ciudad).
\end{itemize}

\section{5.6 Patrones de
comunicación}\label{patrones-de-comunicaciuxf3n}

\begin{itemize}
\tightlist
\item
  \textbf{Sensor → IoT Core}: pub MQTT QoS 1 sobre TLS 1.2.
\item
  \textbf{IoT Rule → Lambda}: invocación síncrona del runtime IoT
  (mantiene la semántica QoS 1).
\item
  \textbf{Lambda ingest → Lambda aggregator}: invocación asíncrona
  (\texttt{InvocationType=Event}) para no bloquear la ingesta.
\item
  \textbf{API Gateway → Lambda api}: integración AWS\_PROXY, sin
  transformaciones en el gateway.
\item
  \textbf{CORS habilitado} en la API para que el dashboard (servido en
  \texttt{localhost:8501}) pueda consumirla.
\end{itemize}

\section{5.7 Modelo serverless y
trade-offs}\label{modelo-serverless-y-trade-offs}

La elección de IoT Core + Lambda + DynamoDB + API Gateway responde a un
patrón \textbf{serverless puro}:

{\def\LTcaptype{none} % do not increment counter
\begin{longtable}[]{@{}
  >{\raggedright\arraybackslash}p{(\linewidth - 2\tabcolsep) * \real{0.3750}}
  >{\raggedright\arraybackslash}p{(\linewidth - 2\tabcolsep) * \real{0.6250}}@{}}
\toprule\noalign{}
\begin{minipage}[b]{\linewidth}\raggedright
Ventaja
\end{minipage} & \begin{minipage}[b]{\linewidth}\raggedright
Justificación
\end{minipage} \\
\midrule\noalign{}
\endhead
\bottomrule\noalign{}
\endlastfoot
Sin servidores que mantener & El equipo se centra en el problema de
negocio. \\
Coste proporcional al uso & Si no hay tráfico, no hay coste. \\
Escalado automático & Lambda y DynamoDB on-demand absorben picos sin
reconfigurar. \\
Alta disponibilidad nativa & Multi-AZ por defecto en cada servicio
gestionado. \\
\end{longtable}
}

{\def\LTcaptype{none} % do not increment counter
\begin{longtable}[]{@{}
  >{\raggedright\arraybackslash}p{(\linewidth - 2\tabcolsep) * \real{0.5217}}
  >{\raggedright\arraybackslash}p{(\linewidth - 2\tabcolsep) * \real{0.4783}}@{}}
\toprule\noalign{}
\begin{minipage}[b]{\linewidth}\raggedright
Limitación
\end{minipage} & \begin{minipage}[b]{\linewidth}\raggedright
Mitigación
\end{minipage} \\
\midrule\noalign{}
\endhead
\bottomrule\noalign{}
\endlastfoot
Cold-start de Lambda & Manejable en este caso (\textless{} 500 ms en
Python 3.12). Para producción crítica, Provisioned Concurrency. \\
Escaneos en DynamoDB & Aceptable para el volumen del piloto. En ciudad:
GSI o particionado. \\
Coste por evento puede crecer con el volumen & En ciudad, evaluar
Kinesis + Flink para batchear. \\
\end{longtable}
}

\section{5.8 Implementación efectiva en el prototipo
(resumen)}\label{implementaciuxf3n-efectiva-en-el-prototipo-resumen}

{\def\LTcaptype{none} % do not increment counter
\begin{longtable}[]{@{}
  >{\raggedright\arraybackslash}p{(\linewidth - 4\tabcolsep) * \real{0.3600}}
  >{\raggedright\arraybackslash}p{(\linewidth - 4\tabcolsep) * \real{0.3200}}
  >{\raggedright\arraybackslash}p{(\linewidth - 4\tabcolsep) * \real{0.3200}}@{}}
\toprule\noalign{}
\begin{minipage}[b]{\linewidth}\raggedright
Recurso
\end{minipage} & \begin{minipage}[b]{\linewidth}\raggedright
Nombre
\end{minipage} & \begin{minipage}[b]{\linewidth}\raggedright
Estado
\end{minipage} \\
\midrule\noalign{}
\endhead
\bottomrule\noalign{}
\endlastfoot
Thing Type & \texttt{smart-parking-albacete-sensor-type} & Creado \\
Thing Group & \texttt{smart-parking-albacete-fleet} & Creado \\
Things & \texttt{ALB-Z1-001} \ldots{} \texttt{ALB-Z4-010} (40 unidades)
& Creados \\
Certificado X.509 & Compartido & Creado \\
IoT Policy & \texttt{smart-parking-albacete-sensor-policy} & Creada \\
IoT Topic Rule & \texttt{smart\_parking\_albacete\_ingest\_rule} &
Creada \\
Tabla DynamoDB & \texttt{smart-parking-albacete-state} & Creada \\
Tabla DynamoDB & \texttt{smart-parking-albacete-zone-kpis} & Creada \\
Lambda & \texttt{smart-parking-albacete-ingest} & Desplegada \\
Lambda & \texttt{smart-parking-albacete-aggregator} & Desplegada \\
Lambda & \texttt{smart-parking-albacete-api} & Desplegada \\
API Gateway & \texttt{smart-parking-albacete-api} (id
\texttt{85fbp0svzc}) & Desplegada en stage \texttt{prod} \\
URL pública &
\texttt{https://85fbp0svzc.execute-api.us-east-1.amazonaws.com/prod} &
Operativa durante la sesión del lab \\
\end{longtable}
}

El cuaderno de despliegue, comandos exactos y verificación end-to-end
están en el \texttt{README\_GUIA.md} raíz y en el capítulo 11.

\newpage

\chapter{06. Integración con FIWARE y Apache Flink
(diseño)}\label{integraciuxf3n-con-fiware-y-apache-flink-diseuxf1o}

El enunciado del campus virtual exige una arquitectura cloud sobre AWS.
Sin embargo, los temas 4 y 5 de la asignatura introducen \textbf{FIWARE}
(con Orion Context Broker y la API NGSI-v2) y \textbf{Apache Flink} como
pieza de procesamiento de streams. Este capítulo razona cómo integrar
ambos componentes en la solución, qué ventajas aportan respecto a la
arquitectura serverless básica, y por qué se ha decidido
\textbf{documentarlos en el diseño sin incluirlos en la implementación
del prototipo}.

\section{6.1 Por qué considerar FIWARE en un proyecto sobre
AWS}\label{por-quuxe9-considerar-fiware-en-un-proyecto-sobre-aws}

FIWARE es un ecosistema \emph{open source} impulsado por la Unión
Europea y adoptado masivamente por proyectos de smart cities a nivel
europeo. Sus puntos fuertes son:

\begin{itemize}
\tightlist
\item
  \textbf{Modelo de datos compartido} (NGSI-v2 y, más recientemente,
  NGSI-LD) que facilita la \textbf{interoperabilidad} entre ciudades y
  entre dominios verticales (movilidad, calidad del aire,
  residuos\ldots).
\item
  \textbf{Catálogo de Smart Data Models} estandarizados (existen modelos
  específicos para \texttt{Parking}, \texttt{ParkingSpot},
  \texttt{OffStreetParking}, \texttt{OnStreetParking}).
\item
  \textbf{Orion Context Broker} como gestor de contexto en tiempo real
  con un mecanismo nativo de \textbf{suscripciones} (push) hacia
  consumidores arbitrarios.
\end{itemize}

Para un ayuntamiento, adoptar (al menos a nivel de contrato de API)
NGSI-v2 simplifica:

\begin{itemize}
\tightlist
\item
  La participación en redes europeas de smart cities (proyectos H2020 /
  Horizon Europe).
\item
  La sustitución del proveedor cloud sin que las aplicaciones
  consumidoras tengan que rehacerse.
\item
  La compatibilidad con proyectos del propio ayuntamiento ya basados en
  FIWARE.
\end{itemize}

\section{6.2 Arquitectura híbrida AWS +
FIWARE}\label{arquitectura-huxedbrida-aws-fiware}

La propuesta de TECO S.L. para fases posteriores del proyecto contempla
una \textbf{arquitectura híbrida} en la que AWS sigue siendo la espina
dorsal de ingesta y persistencia (decisión imposible de revertir tras la
inversión del piloto) y FIWARE Orion actúa como capa de \textbf{contexto
interoperable} expuesta a terceros y a otros dominios verticales
municipales.

\begin{Shaded}
\begin{Highlighting}[]
\NormalTok{flowchart LR}
\NormalTok{    SEN[Sensor magnetico] {-}{-}\textgreater{} GW[Gateway edge / IoT Greengrass]}
\NormalTok{    GW {-}{-}\textgreater{} IoT[AWS IoT Core]}
\NormalTok{    IoT {-}{-}\textgreater{} Rule\{IoT Rule\}}
\NormalTok{    Rule {-}{-}\textgreater{} Lam[Lambda normalizador]}
\NormalTok{    Lam {-}{-}\textgreater{} Orion[FIWARE Orion Context Broker]}
\NormalTok{    Orion {-}{-}\textgreater{} Mongo[(MongoDB)]}
\NormalTok{    Orion {-}{-} Suscripciones HTTP {-}{-}\textgreater{} Flink[Apache Flink]}
\NormalTok{    Orion {-}{-} NGSI{-}v2 {-}{-}\textgreater{} Apps[Aplicaciones / Vehiculos / OEM]}
\NormalTok{    Lam {-}{-}\textgreater{} DDB[(DynamoDB estado AWS)]}
\NormalTok{    Flink {-}{-}\textgreater{} DDB}
\NormalTok{    Flink {-}{-}\textgreater{} Dash[Dashboard]}
\end{Highlighting}
\end{Shaded}

Patrón de operación propuesto:

\begin{enumerate}
\def\labelenumi{\arabic{enumi}.}
\tightlist
\item
  El sensor publica MQTT como hasta ahora.
\item
  La IoT Rule sigue invocando la Lambda de ingesta, que:

  \begin{itemize}
  \tightlist
  \item
    Persiste el estado en DynamoDB (rápido, gestionado, sin downtime).
  \item
    Hace un POST a Orion (\texttt{/v2/entities/\{id\}/attrs}) con el
    nuevo estado, manteniendo el contexto NGSI-v2 sincronizado.
  \end{itemize}
\item
  Orion ofrece la API estándar a terceros que la prefieran a la REST
  nativa.
\item
  Orion notifica a Flink mediante una suscripción HTTP cada vez que
  cambia el atributo \texttt{status}.
\item
  Flink calcula agregados por ventanas y los devuelve a DynamoDB o a una
  nueva entidad NGSI-v2 (\texttt{Zone:UCLM:A} con atributos
  \texttt{freeSpots}, \texttt{occupiedSpots}, \texttt{occupancyRate}).
\end{enumerate}

\section{6.3 Modelo de entidades
NGSI-v2}\label{modelo-de-entidades-ngsi-v2}

Para alinearse con los Smart Data Models públicos:

\subsection{\texorpdfstring{Entidad
\texttt{ParkingSpot}}{Entidad ParkingSpot}}\label{entidad-parkingspot}

\begin{Shaded}
\begin{Highlighting}[]
\FunctionTok{\{}
  \DataTypeTok{"id"}\FunctionTok{:} \StringTok{"ParkingSpot:Albacete:ALB{-}Z1{-}001"}\FunctionTok{,}
  \DataTypeTok{"type"}\FunctionTok{:} \StringTok{"ParkingSpot"}\FunctionTok{,}
  \DataTypeTok{"status"}\FunctionTok{:} \FunctionTok{\{}
    \DataTypeTok{"type"}\FunctionTok{:} \StringTok{"Text"}\FunctionTok{,}
    \DataTypeTok{"value"}\FunctionTok{:} \StringTok{"free"}\FunctionTok{,}
    \DataTypeTok{"metadata"}\FunctionTok{:} \FunctionTok{\{}
      \DataTypeTok{"timestamp"}\FunctionTok{:} \FunctionTok{\{}
        \DataTypeTok{"type"}\FunctionTok{:} \StringTok{"DateTime"}\FunctionTok{,}
        \DataTypeTok{"value"}\FunctionTok{:} \StringTok{"2026{-}05{-}16T10:00:00Z"}
      \FunctionTok{\}}
    \FunctionTok{\}}
  \FunctionTok{\},}
  \DataTypeTok{"location"}\FunctionTok{:} \FunctionTok{\{}
    \DataTypeTok{"type"}\FunctionTok{:} \StringTok{"geo:json"}\FunctionTok{,}
    \DataTypeTok{"value"}\FunctionTok{:} \FunctionTok{\{}\DataTypeTok{"type"}\FunctionTok{:} \StringTok{"Point"}\FunctionTok{,} \DataTypeTok{"coordinates"}\FunctionTok{:} \OtherTok{[}\FloatTok{{-}1.85745}\OtherTok{,} \FloatTok{38.97765}\OtherTok{]}\FunctionTok{\}}
  \FunctionTok{\},}
  \DataTypeTok{"refParkingZone"}\FunctionTok{:} \FunctionTok{\{}\DataTypeTok{"type"}\FunctionTok{:} \StringTok{"Relationship"}\FunctionTok{,} \DataTypeTok{"value"}\FunctionTok{:} \StringTok{"ParkingZone:Albacete:Z1{-}CAMPUS"}\FunctionTok{\},}
  \DataTypeTok{"refParkingSensor"}\FunctionTok{:} \FunctionTok{\{}\DataTypeTok{"type"}\FunctionTok{:} \StringTok{"Relationship"}\FunctionTok{,} \DataTypeTok{"value"}\FunctionTok{:} \StringTok{"ParkingSensor:Albacete:ALB{-}Z1{-}001"}\FunctionTok{\}}
\FunctionTok{\}}
\end{Highlighting}
\end{Shaded}

\subsection{\texorpdfstring{Entidad
\texttt{ParkingSensor}}{Entidad ParkingSensor}}\label{entidad-parkingsensor}

\begin{Shaded}
\begin{Highlighting}[]
\FunctionTok{\{}
  \DataTypeTok{"id"}\FunctionTok{:} \StringTok{"ParkingSensor:Albacete:ALB{-}Z1{-}001"}\FunctionTok{,}
  \DataTypeTok{"type"}\FunctionTok{:} \StringTok{"ParkingSensor"}\FunctionTok{,}
  \DataTypeTok{"sensorType"}\FunctionTok{:} \FunctionTok{\{}\DataTypeTok{"type"}\FunctionTok{:} \StringTok{"Text"}\FunctionTok{,} \DataTypeTok{"value"}\FunctionTok{:} \StringTok{"magnetic"}\FunctionTok{\},}
  \DataTypeTok{"batteryLevel"}\FunctionTok{:} \FunctionTok{\{}\DataTypeTok{"type"}\FunctionTok{:} \StringTok{"Number"}\FunctionTok{,} \DataTypeTok{"value"}\FunctionTok{:} \DecValTok{92}\FunctionTok{\},}
  \DataTypeTok{"assignedSpot"}\FunctionTok{:} \FunctionTok{\{}\DataTypeTok{"type"}\FunctionTok{:} \StringTok{"Relationship"}\FunctionTok{,} \DataTypeTok{"value"}\FunctionTok{:} \StringTok{"ParkingSpot:Albacete:ALB{-}Z1{-}001"}\FunctionTok{\}}
\FunctionTok{\}}
\end{Highlighting}
\end{Shaded}

\subsection{\texorpdfstring{Entidad
\texttt{ParkingZone}}{Entidad ParkingZone}}\label{entidad-parkingzone}

\begin{Shaded}
\begin{Highlighting}[]
\FunctionTok{\{}
  \DataTypeTok{"id"}\FunctionTok{:} \StringTok{"ParkingZone:Albacete:Z1{-}CAMPUS"}\FunctionTok{,}
  \DataTypeTok{"type"}\FunctionTok{:} \StringTok{"ParkingZone"}\FunctionTok{,}
  \DataTypeTok{"totalSpots"}\FunctionTok{:} \FunctionTok{\{}\DataTypeTok{"type"}\FunctionTok{:} \StringTok{"Number"}\FunctionTok{,} \DataTypeTok{"value"}\FunctionTok{:} \DecValTok{125}\FunctionTok{\},}
  \DataTypeTok{"freeSpots"}\FunctionTok{:}  \FunctionTok{\{}\DataTypeTok{"type"}\FunctionTok{:} \StringTok{"Number"}\FunctionTok{,} \DataTypeTok{"value"}\FunctionTok{:} \DecValTok{56}\FunctionTok{\},}
  \DataTypeTok{"occupancyRate"}\FunctionTok{:} \FunctionTok{\{}\DataTypeTok{"type"}\FunctionTok{:} \StringTok{"Number"}\FunctionTok{,} \DataTypeTok{"value"}\FunctionTok{:} \FloatTok{0.55}\FunctionTok{\}}
\FunctionTok{\}}
\end{Highlighting}
\end{Shaded}

\section{6.4 Apache Flink: cuándo y para
qué}\label{apache-flink-cuuxe1ndo-y-para-quuxe9}

Para el \textbf{piloto} (≈500 plazas, \textasciitilde150 000
eventos/día), una Lambda agregador escrita en Python es más que
suficiente: latencia milisegundo, coste mensual irrelevante, complejidad
de mantenimiento nula.

En cambio, en el \textbf{escenario ciudad} (≥ 10 000 plazas, ≥ 3 000 000
eventos/día) Flink aporta ventajas que justifican introducirlo:

{\def\LTcaptype{none} % do not increment counter
\begin{longtable}[]{@{}
  >{\raggedright\arraybackslash}p{(\linewidth - 2\tabcolsep) * \real{0.4118}}
  >{\raggedright\arraybackslash}p{(\linewidth - 2\tabcolsep) * \real{0.5882}}@{}}
\toprule\noalign{}
\begin{minipage}[b]{\linewidth}\raggedright
Capacidad de Flink
\end{minipage} & \begin{minipage}[b]{\linewidth}\raggedright
Aplicación en smart parking
\end{minipage} \\
\midrule\noalign{}
\endhead
\bottomrule\noalign{}
\endlastfoot
\textbf{Procesamiento de streams} real (ventanas, watermarks) & Cálculo
continuo de ocupación con tolerancia a eventos fuera de orden. \\
\textbf{Ventanas tumbling} & KPIs por intervalo fijo (cada 1 min, 5 min,
1 h) sin recalcular toda la zona. \\
\textbf{Ventanas sliding} & Tendencias suavizadas (\% ocupación últimos
10 min refrescado cada 1 min) para paneles ciudadanos. \\
\textbf{Ventanas session} & Tiempo medio de estacionamiento real por
plaza, base para tarificación dinámica. \\
\textbf{CEP (Complex Event Processing)} & Detección de patrones:
sensores que reportan ocupación con intermitencia anómala (signo de
avería). \\
\textbf{Escalabilidad horizontal} & Paralelización por particiones del
topic Kafka/Kinesis. \\
\end{longtable}
}

Patrón típico recomendado:

\begin{Shaded}
\begin{Highlighting}[]
\NormalTok{flowchart LR}
\NormalTok{    IOT[AWS IoT Core] {-}{-}\textgreater{} KS[Kinesis Data Stream]}
\NormalTok{    KS {-}{-}\textgreater{} F[Flink job]}
\NormalTok{    F {-}{-}\textgreater{} KPIS[(DynamoDB zone{-}kpis)]}
\NormalTok{    F {-}{-}\textgreater{} ORION[FIWARE Orion]}
\NormalTok{    F {-}{-}\textgreater{} ALARM[SNS {-} alertas sensor caido]}
\end{Highlighting}
\end{Shaded}

Snippet conceptual de un job Flink que produce ocupación por zona cada
minuto:

\begin{Shaded}
\begin{Highlighting}[]
\NormalTok{DataStream}\OperatorTok{\textless{}}\NormalTok{ParkingEvent}\OperatorTok{\textgreater{}}\NormalTok{ events }\OperatorTok{=}\NormalTok{ env}
    \OperatorTok{.}\FunctionTok{addSource}\OperatorTok{(}\NormalTok{kinesisSource}\OperatorTok{)}
    \OperatorTok{.}\FunctionTok{assignTimestampsAndWatermarks}\OperatorTok{(}\NormalTok{WatermarkStrategy}
        \OperatorTok{.\textless{}}\NormalTok{ParkingEvent}\OperatorTok{\textgreater{}}\FunctionTok{forBoundedOutOfOrderness}\OperatorTok{(}\BuiltInTok{Duration}\OperatorTok{.}\FunctionTok{ofSeconds}\OperatorTok{(}\DecValTok{10}\OperatorTok{))}
        \OperatorTok{.}\FunctionTok{withTimestampAssigner}\OperatorTok{((}\NormalTok{e}\OperatorTok{,}\NormalTok{ t}\OperatorTok{)} \OperatorTok{{-}\textgreater{}}\NormalTok{ e}\OperatorTok{.}\FunctionTok{timestamp}\OperatorTok{));}

\NormalTok{events}
    \OperatorTok{.}\FunctionTok{keyBy}\OperatorTok{(}\NormalTok{ParkingEvent}\OperatorTok{::}\NormalTok{getZoneId}\OperatorTok{)}
    \OperatorTok{.}\FunctionTok{window}\OperatorTok{(}\NormalTok{TumblingEventTimeWindows}\OperatorTok{.}\FunctionTok{of}\OperatorTok{(}\BuiltInTok{Time}\OperatorTok{.}\FunctionTok{minutes}\OperatorTok{(}\DecValTok{1}\OperatorTok{)))}
    \OperatorTok{.}\FunctionTok{aggregate}\OperatorTok{(}\KeywordTok{new} \FunctionTok{ZoneOccupancyAggregator}\OperatorTok{())}
    \OperatorTok{.}\FunctionTok{addSink}\OperatorTok{(}\KeywordTok{new} \FunctionTok{DynamoDBSink}\OperatorTok{(}\NormalTok{kpisTable}\OperatorTok{));}
\end{Highlighting}
\end{Shaded}

\section{6.5 Decisión final: implementar solo en AWS para el
piloto}\label{decisiuxf3n-final-implementar-solo-en-aws-para-el-piloto}

Razones para \textbf{no} desplegar FIWARE/Flink en el prototipo:

\begin{enumerate}
\def\labelenumi{\arabic{enumi}.}
\tightlist
\item
  El AWS Academy Learner Lab tiene 3 h de vida y permisos restringidos
  (no se pueden crear roles IAM nuevos para servicios fuera de AWS).
\item
  Levantar Orion + MongoDB + Flink añade \textasciitilde3 contenedores
  Docker y, en una sesión académica, multiplica los puntos de fallo sin
  aportar valor académico extra (los conceptos quedan demostrados con la
  implementación AWS y este diseño).
\item
  La curva de coste/beneficio de Flink solo es positiva en el escenario
  ciudad. En el piloto, la Lambda agregador cumple los requisitos con un
  orden de magnitud menos de complejidad.
\item
  La interoperabilidad NGSI-v2 puede \textbf{simularse a nivel de
  contrato}: el cliente que quiera consumir vía NGSI-v2 verá la misma
  información que la API REST nativa, y la Lambda API podría exponer una
  segunda ruta que emule el formato NGSI-v2 a partir de DynamoDB
  (trabajo futuro de bajo coste).
\end{enumerate}

\section{6.6 Mapping AWS REST ↔ NGSI-v2 (referencia para
defensa)}\label{mapping-aws-rest-ngsi-v2-referencia-para-defensa}

{\def\LTcaptype{none} % do not increment counter
\begin{longtable}[]{@{}
  >{\raggedright\arraybackslash}p{(\linewidth - 2\tabcolsep) * \real{0.5082}}
  >{\raggedright\arraybackslash}p{(\linewidth - 2\tabcolsep) * \real{0.4918}}@{}}
\toprule\noalign{}
\begin{minipage}[b]{\linewidth}\raggedright
Endpoint REST (esta solución)
\end{minipage} & \begin{minipage}[b]{\linewidth}\raggedright
Endpoint NGSI-v2 equivalente
\end{minipage} \\
\midrule\noalign{}
\endhead
\bottomrule\noalign{}
\endlastfoot
\texttt{GET\ /spots} & \texttt{GET\ /v2/entities?type=ParkingSpot} \\
\texttt{GET\ /spots/\{id\}} & \texttt{GET\ /v2/entities/\{id\}} \\
\texttt{GET\ /zones} & \texttt{GET\ /v2/entities?type=ParkingZone} \\
\texttt{GET\ /zones/\{id\}/kpis} &
\texttt{GET\ /v2/entities/\{id\}?attrs=freeSpots,occupiedSpots,occupancyRate}
(más histórico en STH-Comet o Cygnus) \\
Notificación push a tercero & \texttt{POST\ /v2/subscriptions} con
\texttt{notification.http.url} \\
\end{longtable}
}

\newpage

\chapter{07. Modelo maestro de datos y
APIs}\label{modelo-maestro-de-datos-y-apis}

Este capítulo formaliza el \textbf{modelo de datos} que circula y se
persiste en el sistema, así como la \textbf{API REST} que se expone al
dashboard, a las aplicaciones del operador y a sistemas de terceros
(vehículos autónomos, plataformas municipales de movilidad).

\section{7.1 Modelo conceptual}\label{modelo-conceptual}

Las entidades principales del dominio son:

\begin{Shaded}
\begin{Highlighting}[]
\NormalTok{classDiagram}
\NormalTok{    class ParkingSpot \{}
\NormalTok{        +String spotId}
\NormalTok{        +String zoneId}
\NormalTok{        +String street}
\NormalTok{        +float lat}
\NormalTok{        +float lon}
\NormalTok{        +String status  "free|occupied|unknown"}
\NormalTok{        +int batteryLevel}
\NormalTok{        +float confidence}
\NormalTok{        +String sensorType}
\NormalTok{        +long lastUpdated}
\NormalTok{    \}}
\NormalTok{    class ParkingZone \{}
\NormalTok{        +String zoneId}
\NormalTok{        +int totalSpots}
\NormalTok{        +int freeSpots}
\NormalTok{        +int occupiedSpots}
\NormalTok{        +int unknownSpots}
\NormalTok{        +float occupancyRate}
\NormalTok{        +DateTime windowEnd}
\NormalTok{    \}}
\NormalTok{    class ParkingSensor \{}
\NormalTok{        +String sensorId}
\NormalTok{        +String technology}
\NormalTok{        +int batteryLevel}
\NormalTok{        +String firmwareVersion}
\NormalTok{    \}}
\NormalTok{    ParkingSpot "1" {-}{-} "1" ParkingSensor : equipado por}
\NormalTok{    ParkingSpot "*" {-}{-} "1" ParkingZone : pertenece a}
\end{Highlighting}
\end{Shaded}

\section{7.2 Modelado físico en
DynamoDB}\label{modelado-fuxedsico-en-dynamodb}

\subsection{\texorpdfstring{Tabla
\texttt{smart-parking-albacete-state}}{Tabla smart-parking-albacete-state}}\label{tabla-smart-parking-albacete-state}

{\def\LTcaptype{none} % do not increment counter
\begin{longtable}[]{@{}lll@{}}
\toprule\noalign{}
Atributo & Tipo & Rol \\
\midrule\noalign{}
\endhead
\bottomrule\noalign{}
\endlastfoot
\texttt{spotId} & String & Clave de partición (PK) \\
\texttt{zoneId} & String & Sub-zona \\
\texttt{street} & String & Calle \\
\texttt{lat}, \texttt{lon} & Number & Coordenadas \\
\texttt{status} & String & \texttt{free} / \texttt{occupied} /
\texttt{unknown} \\
\texttt{batteryLevel} & Number & \% batería \\
\texttt{confidence} & Number & Confianza de la última medición \\
\texttt{sensorType} & String & \texttt{magnetic} por defecto \\
\texttt{lastUpdated} & Number & Timestamp ms epoch \\
\end{longtable}
}

Modo de facturación: \textbf{PAY\_PER\_REQUEST} (on-demand): elimina la
necesidad de provisionar capacidad y absorbe picos sin throttling.

Operaciones más frecuentes:

\begin{itemize}
\tightlist
\item
  \texttt{PutItem} (UPSERT) desde Lambda ingest.
\item
  \texttt{Scan} filtrado por \texttt{zoneId} desde Lambda aggregator y
  Lambda api.
\end{itemize}

En escenario ciudad se introduciría un \textbf{GSI} con \texttt{zoneId}
como PK para sustituir el \texttt{Scan\ +\ Filter} por \texttt{Query}.

\subsection{\texorpdfstring{Tabla
\texttt{smart-parking-albacete-zone-kpis}}{Tabla smart-parking-albacete-zone-kpis}}\label{tabla-smart-parking-albacete-zone-kpis}

{\def\LTcaptype{none} % do not increment counter
\begin{longtable}[]{@{}lll@{}}
\toprule\noalign{}
Atributo & Tipo & Rol \\
\midrule\noalign{}
\endhead
\bottomrule\noalign{}
\endlastfoot
\texttt{zoneId} & String & Clave de partición (PK) \\
\texttt{windowEnd} & String & Clave de ordenación (SK), ISO 8601 \\
\texttt{totalSpots} & Number & Plazas totales en la zona \\
\texttt{freeSpots} & Number & Plazas libres \\
\texttt{occupiedSpots} & Number & Plazas ocupadas \\
\texttt{unknownSpots} & Number & Plazas sin datos \\
\texttt{occupancyRate} & Number & Ratio 0..1 \\
\texttt{computedAtMs} & Number & Timestamp ms del cálculo \\
\end{longtable}
}

Esta tabla actúa como \textbf{serie temporal} indexada por zona. Permite
consultas como ``últimos N KPIs de Z1-CAMPUS'' con \texttt{Query}
natural (descendiente por \texttt{windowEnd}).

En el escenario ciudad esta tabla podría migrarse a \textbf{Amazon
Timestream} para aprovechar funciones temporales nativas y agregaciones
eficientes.

\section{7.3 Esquema del payload MQTT (sensor → IoT
Core)}\label{esquema-del-payload-mqtt-sensor-iot-core}

\begin{Shaded}
\begin{Highlighting}[]
\FunctionTok{\{}
  \DataTypeTok{"spotId"}\FunctionTok{:} \StringTok{"ALB{-}Z1{-}001"}\FunctionTok{,}
  \DataTypeTok{"zoneId"}\FunctionTok{:} \StringTok{"Z1{-}CAMPUS"}\FunctionTok{,}
  \DataTypeTok{"street"}\FunctionTok{:} \StringTok{"C. Imperial"}\FunctionTok{,}
  \DataTypeTok{"lat"}\FunctionTok{:} \FloatTok{38.97765}\FunctionTok{,}
  \DataTypeTok{"lon"}\FunctionTok{:} \FloatTok{{-}1.85745}\FunctionTok{,}
  \DataTypeTok{"status"}\FunctionTok{:} \StringTok{"occupied"}\FunctionTok{,}
  \DataTypeTok{"batteryLevel"}\FunctionTok{:} \FloatTok{92.2}\FunctionTok{,}
  \DataTypeTok{"confidence"}\FunctionTok{:} \FloatTok{0.94}\FunctionTok{,}
  \DataTypeTok{"sensorType"}\FunctionTok{:} \StringTok{"magnetic"}\FunctionTok{,}
  \DataTypeTok{"timestamp"}\FunctionTok{:} \DecValTok{1778929200072}
\FunctionTok{\}}
\end{Highlighting}
\end{Shaded}

Topic: \texttt{parking/\{zoneId\}/spot/\{spotId\}/status}

\section{7.4 API REST: contrato OpenAPI
3.0}\label{api-rest-contrato-openapi-3.0}

El archivo \texttt{prototipo/api/openapi.yaml} contiene la
especificación completa. Resumen de endpoints:

{\def\LTcaptype{none} % do not increment counter
\begin{longtable}[]{@{}
  >{\raggedright\arraybackslash}p{(\linewidth - 6\tabcolsep) * \real{0.2051}}
  >{\raggedright\arraybackslash}p{(\linewidth - 6\tabcolsep) * \real{0.1538}}
  >{\raggedright\arraybackslash}p{(\linewidth - 6\tabcolsep) * \real{0.3333}}
  >{\raggedright\arraybackslash}p{(\linewidth - 6\tabcolsep) * \real{0.3077}}@{}}
\toprule\noalign{}
\begin{minipage}[b]{\linewidth}\raggedright
Método
\end{minipage} & \begin{minipage}[b]{\linewidth}\raggedright
Ruta
\end{minipage} & \begin{minipage}[b]{\linewidth}\raggedright
Descripción
\end{minipage} & \begin{minipage}[b]{\linewidth}\raggedright
Parámetros
\end{minipage} \\
\midrule\noalign{}
\endhead
\bottomrule\noalign{}
\endlastfoot
GET & \texttt{/spots} & Lista plazas (estado actual) & \texttt{?zone=}
(filtra), \texttt{?format=geojson} (RFC 7946) \\
GET & \texttt{/spots/\{spotId\}} & Detalle de una plaza & --- \\
GET & \texttt{/zones} & Lista de sub-zonas con KPIs vivos & --- \\
GET & \texttt{/zones/\{zoneId\}/kpis} & Serie temporal de KPIs &
\texttt{?limit=} (default 100) \\
\end{longtable}
}

Todos los endpoints devuelven
\texttt{Content-Type:\ application/json;\ charset=utf-8}. CORS
habilitado para el dashboard local.

\subsection{Ejemplo de petición}\label{ejemplo-de-peticiuxf3n}

\begin{Shaded}
\begin{Highlighting}[]
\ExtensionTok{curl}\NormalTok{ https://85fbp0svzc.execute{-}api.us{-}east{-}1.amazonaws.com/prod/zones}
\end{Highlighting}
\end{Shaded}

\subsection{Ejemplo de respuesta}\label{ejemplo-de-respuesta}

\begin{Shaded}
\begin{Highlighting}[]
\FunctionTok{\{}
  \DataTypeTok{"count"}\FunctionTok{:} \DecValTok{4}\FunctionTok{,}
  \DataTypeTok{"items"}\FunctionTok{:} \OtherTok{[}
    \FunctionTok{\{}\DataTypeTok{"zoneId"}\FunctionTok{:} \StringTok{"Z1{-}CAMPUS"}\FunctionTok{,}    \DataTypeTok{"total"}\FunctionTok{:} \DecValTok{10}\FunctionTok{,} \DataTypeTok{"free"}\FunctionTok{:} \DecValTok{6}\FunctionTok{,} \DataTypeTok{"occupied"}\FunctionTok{:} \DecValTok{4}\FunctionTok{,} \DataTypeTok{"unknown"}\FunctionTok{:} \DecValTok{0}\FunctionTok{,} \DataTypeTok{"occupancyRate"}\FunctionTok{:} \FloatTok{0.4}\FunctionTok{\}}\OtherTok{,}
    \FunctionTok{\{}\DataTypeTok{"zoneId"}\FunctionTok{:} \StringTok{"Z2{-}DEPORTIVO"}\FunctionTok{,} \DataTypeTok{"total"}\FunctionTok{:} \DecValTok{10}\FunctionTok{,} \DataTypeTok{"free"}\FunctionTok{:} \DecValTok{4}\FunctionTok{,} \DataTypeTok{"occupied"}\FunctionTok{:} \DecValTok{6}\FunctionTok{,} \DataTypeTok{"unknown"}\FunctionTok{:} \DecValTok{0}\FunctionTok{,} \DataTypeTok{"occupancyRate"}\FunctionTok{:} \FloatTok{0.6}\FunctionTok{\}}\OtherTok{,}
    \FunctionTok{\{}\DataTypeTok{"zoneId"}\FunctionTok{:} \StringTok{"Z3{-}SANITARIO"}\FunctionTok{,} \DataTypeTok{"total"}\FunctionTok{:} \DecValTok{10}\FunctionTok{,} \DataTypeTok{"free"}\FunctionTok{:} \DecValTok{6}\FunctionTok{,} \DataTypeTok{"occupied"}\FunctionTok{:} \DecValTok{4}\FunctionTok{,} \DataTypeTok{"unknown"}\FunctionTok{:} \DecValTok{0}\FunctionTok{,} \DataTypeTok{"occupancyRate"}\FunctionTok{:} \FloatTok{0.4}\FunctionTok{\}}\OtherTok{,}
    \FunctionTok{\{}\DataTypeTok{"zoneId"}\FunctionTok{:} \StringTok{"Z4{-}RESIDENCIAL"}\FunctionTok{,}\DataTypeTok{"total"}\FunctionTok{:} \DecValTok{9}\FunctionTok{,}  \DataTypeTok{"free"}\FunctionTok{:} \DecValTok{7}\FunctionTok{,} \DataTypeTok{"occupied"}\FunctionTok{:} \DecValTok{2}\FunctionTok{,} \DataTypeTok{"unknown"}\FunctionTok{:} \DecValTok{0}\FunctionTok{,} \DataTypeTok{"occupancyRate"}\FunctionTok{:} \FloatTok{0.222}\FunctionTok{\}}
  \OtherTok{]}
\FunctionTok{\}}
\end{Highlighting}
\end{Shaded}

\subsection{Ejemplo GeoJSON (consumido por mapas /
OEM)}\label{ejemplo-geojson-consumido-por-mapas-oem}

\begin{Shaded}
\begin{Highlighting}[]
\ExtensionTok{curl} \StringTok{\textquotesingle{}https://.../prod/spots?zone=Z1{-}CAMPUS\&format=geojson\textquotesingle{}}
\end{Highlighting}
\end{Shaded}

\begin{Shaded}
\begin{Highlighting}[]
\FunctionTok{\{}
  \DataTypeTok{"type"}\FunctionTok{:} \StringTok{"FeatureCollection"}\FunctionTok{,}
  \DataTypeTok{"features"}\FunctionTok{:} \OtherTok{[}
    \FunctionTok{\{}
      \DataTypeTok{"type"}\FunctionTok{:} \StringTok{"Feature"}\FunctionTok{,}
      \DataTypeTok{"geometry"}\FunctionTok{:} \FunctionTok{\{}\DataTypeTok{"type"}\FunctionTok{:} \StringTok{"Point"}\FunctionTok{,} \DataTypeTok{"coordinates"}\FunctionTok{:} \OtherTok{[}\FloatTok{{-}1.8572}\OtherTok{,} \FloatTok{38.9781}\OtherTok{]}\FunctionTok{\},}
      \DataTypeTok{"properties"}\FunctionTok{:} \FunctionTok{\{}
        \DataTypeTok{"spotId"}\FunctionTok{:} \StringTok{"ALB{-}Z1{-}002"}\FunctionTok{,}
        \DataTypeTok{"zoneId"}\FunctionTok{:} \StringTok{"Z1{-}CAMPUS"}\FunctionTok{,}
        \DataTypeTok{"status"}\FunctionTok{:} \StringTok{"occupied"}\FunctionTok{,}
        \DataTypeTok{"color"}\FunctionTok{:} \StringTok{"\#e74c3c"}\FunctionTok{,}
        \DataTypeTok{"batteryLevel"}\FunctionTok{:} \FloatTok{92.2}\FunctionTok{,}
        \DataTypeTok{"lastUpdated"}\FunctionTok{:} \DecValTok{1778929200072}
      \FunctionTok{\}}
    \FunctionTok{\}}
  \OtherTok{]}
\FunctionTok{\}}
\end{Highlighting}
\end{Shaded}

\section{7.5 APIs internas vs externas}\label{apis-internas-vs-externas}

La especificación distingue dos capas:

{\def\LTcaptype{none} % do not increment counter
\begin{longtable}[]{@{}
  >{\raggedright\arraybackslash}p{(\linewidth - 4\tabcolsep) * \real{0.1818}}
  >{\raggedright\arraybackslash}p{(\linewidth - 4\tabcolsep) * \real{0.3333}}
  >{\raggedright\arraybackslash}p{(\linewidth - 4\tabcolsep) * \real{0.4848}}@{}}
\toprule\noalign{}
\begin{minipage}[b]{\linewidth}\raggedright
Capa
\end{minipage} & \begin{minipage}[b]{\linewidth}\raggedright
Audiencia
\end{minipage} & \begin{minipage}[b]{\linewidth}\raggedright
Características
\end{minipage} \\
\midrule\noalign{}
\endhead
\bottomrule\noalign{}
\endlastfoot
\textbf{API pública} & Vehículos autónomos, apps ciudadanas, OEM,
ayuntamiento & Solo lectura (GET); rate-limit; cuotas; eventualmente API
keys o Cognito. \\
\textbf{API interna} & Backoffice operador / mantenimiento &
Lectura/escritura sobre \texttt{parking-state}; permite forzar
\texttt{unknown} para plazas en obras; protegida con IAM/Cognito y red
privada. \\
\end{longtable}
}

En el prototipo solo se ha desplegado la API pública (lectura). La API
interna se documenta como ampliación.

\section{7.6 Idempotencia y
consistencia}\label{idempotencia-y-consistencia}

\begin{itemize}
\tightlist
\item
  \textbf{Idempotencia de ingesta}: la Lambda ingest hace
  \texttt{PutItem} (UPSERT). Si el mismo evento llega dos veces (QoS 1
  puede entregar duplicados), el resultado es el mismo.
\item
  \textbf{Consistencia eventual} en DynamoDB: aceptable; la latencia es
  de pocos milisegundos.
\item
  \textbf{Orden de eventos}: el sistema confía en el \texttt{timestamp}
  del sensor. Para escenarios donde un evento más antiguo no debe
  sobrescribir uno más reciente, se introduciría una condición
  \texttt{ConditionExpression:\ lastUpdated\ \textless{}\ :ts}.
\end{itemize}

\section{7.7 Versionado de la API}\label{versionado-de-la-api}

Política de versionado adoptada:

\begin{itemize}
\tightlist
\item
  Versión en el path: \texttt{https://.../v1/spots}. (En el piloto se
  usa el stage \texttt{prod} de API Gateway sin versionado explícito; al
  pasar a productivo se introducirá \texttt{/v1/}).
\item
  Cambios retrocompatibles: adición de atributos opcionales en el JSON.
\item
  Cambios no retrocompatibles: nueva versión \texttt{/v2/}; las dos
  versiones conviven seis meses para migrar a los consumidores.
\end{itemize}

\section{7.8 Documentación legible y
autodescriptiva}\label{documentaciuxf3n-legible-y-autodescriptiva}

El fichero \texttt{openapi.yaml} se sirve también desde un sitio
estático en S3 + CloudFront (no implementado en el prototipo) y puede
explorarse con Swagger UI / Redoc. Esto cubre \textbf{RNF-18}
(autodescriptiva y estándar).

\section{7.9 Modelos de eventos y catálogo de
KPIs}\label{modelos-de-eventos-y-catuxe1logo-de-kpis}

Catálogo de KPIs operacionales que la solución entrega al ayuntamiento:

{\def\LTcaptype{none} % do not increment counter
\begin{longtable}[]{@{}
  >{\raggedright\arraybackslash}p{(\linewidth - 4\tabcolsep) * \real{0.1852}}
  >{\raggedright\arraybackslash}p{(\linewidth - 4\tabcolsep) * \real{0.5185}}
  >{\raggedright\arraybackslash}p{(\linewidth - 4\tabcolsep) * \real{0.2963}}@{}}
\toprule\noalign{}
\begin{minipage}[b]{\linewidth}\raggedright
KPI
\end{minipage} & \begin{minipage}[b]{\linewidth}\raggedright
Periodicidad
\end{minipage} & \begin{minipage}[b]{\linewidth}\raggedright
Origen
\end{minipage} \\
\midrule\noalign{}
\endhead
\bottomrule\noalign{}
\endlastfoot
\% ocupación por sub-zona & Tiempo real & Lambda aggregator \\
\% ocupación por calle & A demanda & \texttt{Scan} sobre
\texttt{parking-state} \\
Plazas libres absolutas por sub-zona & Tiempo real & Lambda
aggregator \\
Tiempo medio de ocupación por plaza & Diario & Trabajo futuro (Flink) \\
Rotación por plaza & Diario & Trabajo futuro (Flink) \\
Sensores caídos / batería baja & Cada hora & Trabajo futuro
(auditoría) \\
Saturación pico (max \% ocupación) & Diario & Trabajo futuro (S3 +
Athena) \\
\end{longtable}
}

\newpage

\chapter{08. Seguridad y privacidad}\label{seguridad-y-privacidad}

La seguridad en IoT no es una capa que se añada al final: hay que
diseñarla en todas las fases del ciclo (provisión, comunicación,
persistencia, exposición y operación). Este capítulo describe las
decisiones de seguridad adoptadas en el prototipo y las ampliaciones
planificadas para producción, distinguiendo siempre lo que ya está
\textbf{implementado} de lo que está \textbf{diseñado pero no
desplegado}.

\section{8.1 Modelo de amenazas
(resumen)}\label{modelo-de-amenazas-resumen}

{\def\LTcaptype{none} % do not increment counter
\begin{longtable}[]{@{}
  >{\raggedright\arraybackslash}p{(\linewidth - 6\tabcolsep) * \real{0.2162}}
  >{\raggedright\arraybackslash}p{(\linewidth - 6\tabcolsep) * \real{0.2432}}
  >{\raggedright\arraybackslash}p{(\linewidth - 6\tabcolsep) * \real{0.2162}}
  >{\raggedright\arraybackslash}p{(\linewidth - 6\tabcolsep) * \real{0.3243}}@{}}
\toprule\noalign{}
\begin{minipage}[b]{\linewidth}\raggedright
Activo
\end{minipage} & \begin{minipage}[b]{\linewidth}\raggedright
Amenaza
\end{minipage} & \begin{minipage}[b]{\linewidth}\raggedright
Vector
\end{minipage} & \begin{minipage}[b]{\linewidth}\raggedright
Mitigación
\end{minipage} \\
\midrule\noalign{}
\endhead
\bottomrule\noalign{}
\endlastfoot
Sensor en calzada & Manipulación física, robo & Acceso físico & Caja
blindada, anclaje, alarma de movimiento. \\
Identidad del sensor & Suplantación & Robo de certificado & Certificado
por dispositivo, rotación, revocación. \\
Canal sensor↔cloud & Eavesdropping, MITM & Red pública & mTLS (TLS 1.2
con certificado mutuo). \\
Cloud (datos) & Acceso no autorizado & Credenciales filtradas & IAM con
principio de mínimo privilegio; auditoría. \\
API pública & Abuso, scraping, DoS & Internet & Throttling, cuotas, API
keys, WAF. \\
Datos personales & Inferencia de hábitos & Trazas de matrículas & El
sistema no almacena matrículas ni datos del conductor. \\
Datos operacionales & Modificación maliciosa & Atacante con credenciales
& Auditoría inmutable, logs en bucket separado. \\
Firmware del sensor & Backdoor en actualizaciones & Despliegue OTA &
Firma criptográfica del firmware, validación en el dispositivo. \\
\end{longtable}
}

\section{8.2 Capa de dispositivo (provisión e
identidad)}\label{capa-de-dispositivo-provisiuxf3n-e-identidad}

\textbf{Implementado en el prototipo:}

\begin{itemize}
\tightlist
\item
  Cada \textbf{Thing} en AWS IoT Core representa una plaza identificable
  de forma única (\texttt{ALB-Z1-001}, \ldots).
\item
  Un único \textbf{certificado X.509} compartido por toda la flota
  piloto, atado a una \textbf{IoT Policy} restrictiva
  (\texttt{smart-parking-albacete-sensor-policy}).
\item
  La policy autoriza únicamente \texttt{Connect}, \texttt{Publish},
  \texttt{Subscribe}, \texttt{Receive} y solo sobre topics que comienzan
  por \texttt{parking/}.
\end{itemize}

\begin{Shaded}
\begin{Highlighting}[]
\FunctionTok{\{}
  \DataTypeTok{"Version"}\FunctionTok{:} \StringTok{"2012{-}10{-}17"}\FunctionTok{,}
  \DataTypeTok{"Statement"}\FunctionTok{:} \OtherTok{[}
    \FunctionTok{\{}\DataTypeTok{"Effect"}\FunctionTok{:} \StringTok{"Allow"}\FunctionTok{,} \DataTypeTok{"Action"}\FunctionTok{:} \StringTok{"iot:Connect"}\FunctionTok{,} \DataTypeTok{"Resource"}\FunctionTok{:} \StringTok{"*"}\FunctionTok{\}}\OtherTok{,}
    \FunctionTok{\{}\DataTypeTok{"Effect"}\FunctionTok{:} \StringTok{"Allow"}\FunctionTok{,} \DataTypeTok{"Action"}\FunctionTok{:} \StringTok{"iot:Publish"}\FunctionTok{,}   \DataTypeTok{"Resource"}\FunctionTok{:} \StringTok{"arn:aws:iot:*:*:topic/parking/*"}\FunctionTok{\}}\OtherTok{,}
    \FunctionTok{\{}\DataTypeTok{"Effect"}\FunctionTok{:} \StringTok{"Allow"}\FunctionTok{,} \DataTypeTok{"Action"}\FunctionTok{:} \StringTok{"iot:Subscribe"}\FunctionTok{,} \DataTypeTok{"Resource"}\FunctionTok{:} \StringTok{"arn:aws:iot:*:*:topicfilter/parking/*"}\FunctionTok{\}}\OtherTok{,}
    \FunctionTok{\{}\DataTypeTok{"Effect"}\FunctionTok{:} \StringTok{"Allow"}\FunctionTok{,} \DataTypeTok{"Action"}\FunctionTok{:} \StringTok{"iot:Receive"}\FunctionTok{,}   \DataTypeTok{"Resource"}\FunctionTok{:} \StringTok{"arn:aws:iot:*:*:topic/parking/*"}\FunctionTok{\}}
  \OtherTok{]}
\FunctionTok{\}}
\end{Highlighting}
\end{Shaded}

\textbf{Diseñado para producción (no en piloto):}

\begin{itemize}
\tightlist
\item
  \textbf{Un certificado por dispositivo}
  (\texttt{Just-in-Time\ Provisioning} o
  \texttt{Fleet\ Provisioning\ Templates}): permite revocaciones
  granulares y auditoría individual.
\item
  \textbf{Renovación periódica} de certificados (cada 12-24 meses)
  mediante el endpoint \texttt{iot:UpdateCertificate}.
\item
  \textbf{AWS IoT Device Defender} activado para auditar configuraciones
  (políticas demasiado permisivas, certificados a punto de expirar) y
  detectar comportamiento anómalo (volumen de mensajes inusual, conexión
  desde IP no habitual).
\end{itemize}

\section{8.3 Capa de transporte}\label{capa-de-transporte}

\begin{itemize}
\tightlist
\item
  \textbf{TLS 1.2} obligatorio en todas las conexiones MQTT (puerto
  8883). Cifrado AES-128/256 según negociación.
\item
  \textbf{mTLS} (autenticación mutua) con certificado del cliente
  firmado por la CA de AWS IoT Core.
\item
  Verificación del certificado raíz \textbf{AmazonRootCA1.pem} en el
  dispositivo (incluido en los certificados que se descargan en
  \texttt{simulator/certs/}).
\item
  API REST \textbf{siempre por HTTPS} (TLS 1.2+) con el certificado
  gestionado por API Gateway.
\end{itemize}

\section{8.4 Capa de procesamiento
(Lambdas)}\label{capa-de-procesamiento-lambdas}

\begin{itemize}
\tightlist
\item
  Las tres Lambdas se ejecutan bajo \texttt{LabRole}, un rol del lab de
  Academy con permisos sobre los servicios que el laboratorio expone. En
  producción se sustituiría por \textbf{un rol IAM por Lambda} con
  políticas a medida (principio de mínimo privilegio):

  \begin{itemize}
  \tightlist
  \item
    \texttt{lambda-ingest-role}: permisos \texttt{dynamodb:PutItem}
    sobre \texttt{parking-state} y \texttt{lambda:InvokeFunction} sobre
    el agregador.
  \item
    \texttt{lambda-aggregator-role}: \texttt{dynamodb:Scan} (con
    condición sobre \texttt{zoneId}) sobre \texttt{parking-state} y
    \texttt{dynamodb:PutItem} sobre \texttt{zone-kpis}.
  \item
    \texttt{lambda-api-role}: \texttt{dynamodb:GetItem} / \texttt{Query}
    / \texttt{Scan} sobre las dos tablas (solo lectura).
  \end{itemize}
\item
  Variables de entorno \textbf{sin credenciales en claro}; los nombres
  de tabla y de función auxiliar viajan como \texttt{STATE\_TABLE},
  \texttt{KPIS\_TABLE}, \texttt{AGGREGATOR\_FN}.
\item
  Trazabilidad: cada invocación genera logs en CloudWatch con request
  ID, código de estado y tiempo.
\end{itemize}

\section{8.5 Capa de persistencia}\label{capa-de-persistencia}

\begin{itemize}
\tightlist
\item
  DynamoDB está cifrado \textbf{en reposo por defecto} con claves
  gestionadas por AWS (AWS KMS).
\item
  En producción se utilizaría \textbf{CMK propia} (Customer Master Key)
  para satisfacer auditorías y políticas internas.
\item
  Backup: DynamoDB on-demand permite habilitar
  \texttt{Point-in-Time\ Recovery} (PITR) con 35 días de retención.
  Recomendado activarlo en producción.
\end{itemize}

\section{8.6 Capa de exposición (API)}\label{capa-de-exposiciuxf3n-api}

\textbf{Implementado:}

\begin{itemize}
\tightlist
\item
  HTTPS con TLS 1.2+ en API Gateway.
\item
  CORS controlado a nivel de Lambda (necesario para Streamlit local).
\end{itemize}

\textbf{Diseñado para producción:}

\begin{itemize}
\tightlist
\item
  \textbf{API Keys + Usage Plans} en API Gateway para terceros:
  throttling por minuto y cuotas mensuales (p.~ej. 10 000 peticiones/día
  por aplicación).
\item
  \textbf{Amazon Cognito} o \textbf{JWT custom} para diferenciar API
  pública (sin auth o con clave) vs API interna del operador
  (autenticada).
\item
  \textbf{AWS WAF} delante de API Gateway para mitigar abuso, OWASP Top
  10 y rate-limiting basado en IP.
\item
  \textbf{Custom domain} con certificado ACM para una URL estable y
  fácil de comunicar a integradores.
\end{itemize}

\section{8.7 Capa de operación}\label{capa-de-operaciuxf3n}

\begin{itemize}
\tightlist
\item
  \textbf{CloudWatch Logs}: retiene los logs de cada Lambda y de la IoT
  Rule con retención configurable.
\item
  \textbf{CloudTrail} (activo por defecto en la cuenta) audita todas las
  llamadas a la API de AWS (quién creó/borró/modificó qué).
\item
  \textbf{AWS Config} (recomendado en producción) para alertar de
  cambios en políticas IAM o en certificados IoT.
\item
  \textbf{Auditoría de manipulación}: los Things creados quedan
  registrados con timestamp; cualquier cambio futuro genera evento de
  CloudTrail.
\end{itemize}

\section{8.8 Cumplimiento RGPD y
privacidad}\label{cumplimiento-rgpd-y-privacidad}

El diseño es \textbf{privacy by design}:

\begin{itemize}
\tightlist
\item
  \textbf{No se almacenan datos personales} identificables. La unidad
  mínima de información es la plaza, no el conductor.
\item
  Las cámaras ANPR (capa complementaria) procesan la matrícula
  \textbf{localmente en el gateway} y solo envían a la cloud un
  \textbf{hash unidireccional} y la dirección de paso. Las imágenes
  originales se retienen un máximo de 72 horas localmente para
  reclamaciones, y después se borran (LOPDGDD y RGPD).
\item
  Las consultas de la API son \textbf{anónimas} para el consumidor; no
  se trazan identidades de usuario final.
\item
  \textbf{Política de retención} para los KPIs: 24 meses agregados; los
  logs operacionales rotan a S3 Glacier después de 90 días.
\end{itemize}

\section{8.9 Plan de respuesta a incidentes
(resumen)}\label{plan-de-respuesta-a-incidentes-resumen}

\begin{enumerate}
\def\labelenumi{\arabic{enumi}.}
\tightlist
\item
  \textbf{Detección}: alertas de CloudWatch o IoT Device Defender
  (intentos de conexión rechazados, picos de tráfico anómalos).
\item
  \textbf{Contención}: revocación inmediata del certificado afectado con
  \texttt{iot:UpdateCertificate(newStatus="REVOKED")}.
\item
  \textbf{Erradicación}: rotación de certificados afectados, sustitución
  física del sensor si hay sospecha de tampering.
\item
  \textbf{Recuperación}: restauración del estado desde DynamoDB PITR si
  fue manipulado.
\item
  \textbf{Post-mortem}: revisión de logs CloudTrail y CloudWatch;
  informe a la concejalía; mejora de la policy o de la lógica afectada.
\end{enumerate}

\section{8.10 Resumen de cumplimiento de requisitos no funcionales de
seguridad}\label{resumen-de-cumplimiento-de-requisitos-no-funcionales-de-seguridad}

{\def\LTcaptype{none} % do not increment counter
\begin{longtable}[]{@{}
  >{\raggedright\arraybackslash}p{(\linewidth - 2\tabcolsep) * \real{0.5789}}
  >{\raggedright\arraybackslash}p{(\linewidth - 2\tabcolsep) * \real{0.4211}}@{}}
\toprule\noalign{}
\begin{minipage}[b]{\linewidth}\raggedright
Requisito
\end{minipage} & \begin{minipage}[b]{\linewidth}\raggedright
Estado
\end{minipage} \\
\midrule\noalign{}
\endhead
\bottomrule\noalign{}
\endlastfoot
RNF-09 (cifrado en tránsito) & OK (MQTT/TLS, HTTPS) \\
RNF-10 (identidad por dispositivo) & Parcial en piloto (cert
compartido); diseñado por dispositivo en producción \\
RNF-11 (control de consumo API) & Diseñado (API Keys + WAF) \\
RNF-12 (sin PII) & OK \\
RNF-13 (auditoría y trazabilidad) & OK (CloudWatch, CloudTrail) \\
RNF-20 (RGPD) & OK (privacy by design) \\
\end{longtable}
}

\newpage

\chapter{09. Escalabilidad: del piloto a la
ciudad}\label{escalabilidad-del-piloto-a-la-ciudad}

El enunciado exige analizar dos escenarios bien diferenciados: un
\textbf{piloto} de centenares de plazas y un \textbf{escenario ciudad}
de miles. Este capítulo cuantifica volumen, capacidad de procesamiento,
límites del sistema y plan de migración entre ambos escenarios.

\section{9.1 Escenario piloto (≈ 500
plazas)}\label{escenario-piloto-500-plazas}

\subsection{Volumen de datos}\label{volumen-de-datos}

\begin{itemize}
\tightlist
\item
  500 plazas × \textasciitilde5 cambios/día = 2 500 eventos/día.
\item
  500 plazas × 288 heartbeats/día (cada 5 min) = 144 000 eventos/día.
\item
  Total ≈ \textbf{150 000 mensajes/día} (≈ 1,7 msg/s medio; pico
  \textasciitilde10 msg/s).
\item
  Tamaño medio del mensaje: \textasciitilde600 bytes →
  \textbf{\textasciitilde85 MB/día} de tráfico bruto.
\end{itemize}

\subsection{Capacidad de
procesamiento}\label{capacidad-de-procesamiento}

{\def\LTcaptype{none} % do not increment counter
\begin{longtable}[]{@{}
  >{\raggedright\arraybackslash}p{(\linewidth - 6\tabcolsep) * \real{0.1429}}
  >{\raggedright\arraybackslash}p{(\linewidth - 6\tabcolsep) * \real{0.3333}}
  >{\raggedright\arraybackslash}p{(\linewidth - 6\tabcolsep) * \real{0.3968}}
  >{\raggedright\arraybackslash}p{(\linewidth - 6\tabcolsep) * \real{0.1270}}@{}}
\toprule\noalign{}
\begin{minipage}[b]{\linewidth}\raggedright
Recurso
\end{minipage} & \begin{minipage}[b]{\linewidth}\raggedright
Límite del servicio
\end{minipage} & \begin{minipage}[b]{\linewidth}\raggedright
Uso esperado del piloto
\end{minipage} & \begin{minipage}[b]{\linewidth}\raggedright
Margen
\end{minipage} \\
\midrule\noalign{}
\endhead
\bottomrule\noalign{}
\endlastfoot
AWS IoT Core & 20 000 conn/s, 30 000 pub/s por cuenta & 1,7 msg/s & x10
000 \\
IoT Rule & ilimitado dentro de la cuenta & 1 regla & OK \\
Lambda concurrente & 1 000 por cuenta default & 1-2 simultáneas & OK \\
DynamoDB on-demand & 40 000 WCU absorbidos sin throttling & 1,7 WCU/s &
x10 000 \\
API Gateway REST & 10 000 req/s soft limit & \textless1 req/s & OK \\
\end{longtable}
}

\subsection{Limitaciones observadas en el
piloto}\label{limitaciones-observadas-en-el-piloto}

\begin{itemize}
\tightlist
\item
  El \texttt{Scan} filtrado por zona crece linealmente. Hasta
  \textasciitilde10 000 plazas no es problema. A partir de ahí conviene
  migrar a \texttt{Query} con GSI por \texttt{zoneId}.
\item
  La invocación asíncrona de la Lambda agregador puede generar múltiples
  ejecuciones por cambio de estado simultáneo, pero el resultado es
  idempotente (se reescribe el KPI con el último cálculo).
\end{itemize}

\section{9.2 Escenario ciudad (≈ 10 000
plazas)}\label{escenario-ciudad-10-000-plazas}

\subsection{Volumen de datos}\label{volumen-de-datos-1}

\begin{itemize}
\tightlist
\item
  10 000 plazas × \textasciitilde5 cambios/día = 50 000 eventos/día.
\item
  10 000 plazas × 288 heartbeats/día = 2 880 000 eventos/día.
\item
  Total ≈ \textbf{3 millones de mensajes/día} (≈ 35 msg/s medio; pico
  \textasciitilde200 msg/s en eventos deportivos).
\item
  Tráfico bruto: \textasciitilde1,7 GB/día.
\end{itemize}

\subsection{Cambios arquitectónicos
recomendados}\label{cambios-arquitectuxf3nicos-recomendados}

{\def\LTcaptype{none} % do not increment counter
\begin{longtable}[]{@{}
  >{\raggedright\arraybackslash}p{(\linewidth - 2\tabcolsep) * \real{0.6000}}
  >{\raggedright\arraybackslash}p{(\linewidth - 2\tabcolsep) * \real{0.4000}}@{}}
\toprule\noalign{}
\begin{minipage}[b]{\linewidth}\raggedright
Componente
\end{minipage} & \begin{minipage}[b]{\linewidth}\raggedright
Cambio
\end{minipage} \\
\midrule\noalign{}
\endhead
\bottomrule\noalign{}
\endlastfoot
\textbf{Sensores} & Onboarding masivo con AWS IoT Fleet Provisioning
Templates. Certificado por dispositivo (no compartido). \\
\textbf{IoT Core} & Sin cambios; sobra capacidad. Habilitar Device
Defender para auditoría continua. \\
\textbf{Procesamiento} & Sustituir la Lambda agregador por un job de
\textbf{Apache Flink en Amazon Kinesis Data Analytics} (o Amazon Managed
Service for Apache Flink). Ventanas tumbling de 1 min para KPIs por
zona. \\
\textbf{Buffer} & Intercalar \textbf{Amazon Kinesis Data Streams} entre
IoT Core y el procesador. Desacopla picos y permite consumidores
múltiples (Flink, archivado en S3, replicación a Orion). \\
\textbf{Estado actual} & DynamoDB on-demand sigue sirviendo, pero añadir
un GSI \texttt{zoneId-status-index} para responder consultas por zona
sin \texttt{Scan}. \\
\textbf{Histórico} & Migrar serie temporal a \textbf{Amazon Timestream}:
agregaciones temporales nativas, retención automática (memoria caliente
+ almacenamiento frío). \\
\textbf{Datalake} & S3 raw como capa bronze; AWS Glue + Athena para
analítica ad-hoc; opcionalmente Redshift/Snowflake para BI municipal. \\
\textbf{APIs} & API Gateway con cuotas y API keys; CloudFront para caché
de las respuestas más leídas (lista de plazas estática + estado
refrescado cada 5 s). \\
\textbf{Multi-AZ} & Confirmar uso de servicios gestionados con
redundancia multi-AZ activada (DynamoDB Global Tables si fuera
multi-región). \\
\end{longtable}
}

\subsection{Arquitectura objetivo para
ciudad}\label{arquitectura-objetivo-para-ciudad}

\begin{Shaded}
\begin{Highlighting}[]
\NormalTok{flowchart LR}
\NormalTok{    SENS[10 000 sensores] {-}{-}\textgreater{} IOTC[AWS IoT Core]}
\NormalTok{    IOTC {-}{-}\textgreater{} KDS[Kinesis Data Streams]}
\NormalTok{    KDS {-}{-}\textgreater{} FLINK[Managed Apache Flink]}
\NormalTok{    KDS {-}{-}\textgreater{} S3R[(S3 raw {-} bronze)]}
\NormalTok{    FLINK {-}{-}\textgreater{} DDB[(DynamoDB state + GSI)]}
\NormalTok{    FLINK {-}{-}\textgreater{} TS[(Timestream KPIs)]}
\NormalTok{    FLINK {-}{-}\textgreater{} ORION[FIWARE Orion {-} opcional]}
\NormalTok{    DDB {-}{-}\textgreater{} APIGW[API Gateway + WAF + Cuotas]}
\NormalTok{    TS {-}{-}\textgreater{} APIGW}
\NormalTok{    APIGW {-}{-}\textgreater{} CF[CloudFront cache]}
\NormalTok{    CF {-}{-}\textgreater{} CONS[Apps / Vehiculos / Operador]}
\NormalTok{    S3R {-}{-}\textgreater{} GLUE[Glue + Athena]}
\NormalTok{    GLUE {-}{-}\textgreater{} QS[QuickSight]}
\end{Highlighting}
\end{Shaded}

\subsection{Coste agregado estimado
(ciudad)}\label{coste-agregado-estimado-ciudad}

Detalle en el capítulo 10. Resumen: \textbf{600-900 USD/mes} en AWS para
los 10 000 sensores en operación normal, no incluyendo el CAPEX de los
sensores (\textasciitilde75 €/unidad).

\section{9.3 Estrategia de
particionado}\label{estrategia-de-particionado}

{\def\LTcaptype{none} % do not increment counter
\begin{longtable}[]{@{}
  >{\raggedright\arraybackslash}p{(\linewidth - 4\tabcolsep) * \real{0.3438}}
  >{\raggedright\arraybackslash}p{(\linewidth - 4\tabcolsep) * \real{0.3125}}
  >{\raggedright\arraybackslash}p{(\linewidth - 4\tabcolsep) * \real{0.3438}}@{}}
\toprule\noalign{}
\begin{minipage}[b]{\linewidth}\raggedright
Partición
\end{minipage} & \begin{minipage}[b]{\linewidth}\raggedright
Criterio
\end{minipage} & \begin{minipage}[b]{\linewidth}\raggedright
Beneficio
\end{minipage} \\
\midrule\noalign{}
\endhead
\bottomrule\noalign{}
\endlastfoot
Sub-zona operativa & \texttt{zoneId} & Localiza KPIs a la unidad de
gestión del ayuntamiento. \\
Geográfica gruesa & Distrito o código postal & Permite reportes
municipales por distrito. \\
Temporal & Por mes (en S3) & Hace eficiente la analítica retrospectiva
con Athena. \\
Por operador concesionario & Tag por concesión & Para zonas reguladas
con gestor externo. \\
\end{longtable}
}

\section{9.4 Estrategia de alta
disponibilidad}\label{estrategia-de-alta-disponibilidad}

\begin{itemize}
\tightlist
\item
  Todos los servicios gestionados (IoT Core, Lambda, DynamoDB, API
  Gateway, S3, Kinesis, Timestream) son \textbf{multi-AZ por defecto} en
  \texttt{us-east-1} / \texttt{eu-west-1}.
\item
  Para tolerancia a la caída de toda una región se planificaría una
  segunda región pasiva con \textbf{DynamoDB Global Tables} y
  replicación de Lambda + API Gateway desplegable por CI/CD (Terraform /
  CDK).
\item
  \textbf{Sensores}: caché local de hasta 24 horas; reenvío en orden al
  recuperarse la conectividad.
\end{itemize}

\section{9.5 Estrategia de despliegue continuo
(CI/CD)}\label{estrategia-de-despliegue-continuo-cicd}

Para el piloto, los scripts \texttt{infra/*.py} son suficientes y se
ejecutan a mano. Para producción se recomienda:

\begin{itemize}
\tightlist
\item
  \textbf{AWS CDK} (TypeScript o Python) o \textbf{Terraform} para
  infraestructura como código versionable.
\item
  \textbf{Pipeline de CI/CD} con GitHub Actions o CodePipeline: tests →
  build → despliegue blue/green con rollback automático.
\item
  \textbf{Tests de carga} trimestrales con Locust o Artillery contra la
  API y un simulador a escala (5 000+ sensores en local).
\end{itemize}

\section{9.6 Plan de pruebas a escala}\label{plan-de-pruebas-a-escala}

{\def\LTcaptype{none} % do not increment counter
\begin{longtable}[]{@{}
  >{\raggedright\arraybackslash}p{(\linewidth - 4\tabcolsep) * \real{0.2963}}
  >{\raggedright\arraybackslash}p{(\linewidth - 4\tabcolsep) * \real{0.3704}}
  >{\raggedright\arraybackslash}p{(\linewidth - 4\tabcolsep) * \real{0.3333}}@{}}
\toprule\noalign{}
\begin{minipage}[b]{\linewidth}\raggedright
Prueba
\end{minipage} & \begin{minipage}[b]{\linewidth}\raggedright
Objetivo
\end{minipage} & \begin{minipage}[b]{\linewidth}\raggedright
Métrica
\end{minipage} \\
\midrule\noalign{}
\endhead
\bottomrule\noalign{}
\endlastfoot
Carga sostenida & Validar 35 msg/s constantes durante 24 h & p99
latencia API \textless{} 2 s \\
Pico de evento deportivo & 1 000 eventos/min durante 15 min & Sin
throttling DynamoDB ni Lambda \\
Caída de IoT Core (simulada) & Comprobar buffering en sensor y reenvío &
Pérdida 0 eventos en cola de \textless{} 24 h \\
Caída de Lambda agregador & Comprobar continuidad de la ingesta & KPIs
eventualmente consistentes al recuperar \\
Fuga de credencial sensor & Revocación rápida del certificado &
\textless{} 5 min entre detección y revocación \\
\end{longtable}
}

\section{9.7 Riesgos específicos del
escalado}\label{riesgos-especuxedficos-del-escalado}

{\def\LTcaptype{none} % do not increment counter
\begin{longtable}[]{@{}
  >{\raggedright\arraybackslash}p{(\linewidth - 2\tabcolsep) * \real{0.4211}}
  >{\raggedright\arraybackslash}p{(\linewidth - 2\tabcolsep) * \real{0.5789}}@{}}
\toprule\noalign{}
\begin{minipage}[b]{\linewidth}\raggedright
Riesgo
\end{minipage} & \begin{minipage}[b]{\linewidth}\raggedright
Mitigación
\end{minipage} \\
\midrule\noalign{}
\endhead
\bottomrule\noalign{}
\endlastfoot
Coste cloud creciendo más rápido que el número de plazas & Revisar
arquitectura: pasar a Kinesis + Flink, batching, ventanas mayores. \\
Saturación de la API durante eventos & CloudFront caching + WAF
rate-limiting + scale-out automático de Lambda. \\
Sensores fuera de cobertura LPWAN & Auditoría periódica con
\texttt{lastUpdated}; envío de cuadrilla. \\
Modelos NGSI-v2 / Smart Data Models incompatibles entre versiones &
Mantener mapping en una Lambda dedicada; tests de contrato semanales. \\
Crecimiento de la tabla \texttt{zone-kpis} & Política de TTL (DynamoDB
TTL nativo) para purgar datos \textgreater{} N meses o migrar a
Timestream. \\
\end{longtable}
}

\newpage

\chapter{10. Análisis de costes}\label{anuxe1lisis-de-costes}

Este capítulo presenta una estimación realista del coste del proyecto en
dos escenarios (piloto y ciudad) descompuesto en \textbf{CAPEX}
(inversión inicial en hardware y obra civil) y \textbf{OPEX} (operación
recurrente, principalmente cloud y conectividad). Los precios son
aproximados y se sitúan en el orden de magnitud habitual del sector en
España a fecha 2026; en una propuesta real se cerrarían con los
proveedores específicos.

\section{10.1 Hipótesis y unidades}\label{hipuxf3tesis-y-unidades}

{\def\LTcaptype{none} % do not increment counter
\begin{longtable}[]{@{}
  >{\raggedright\arraybackslash}p{(\linewidth - 2\tabcolsep) * \real{0.3704}}
  >{\raggedright\arraybackslash}p{(\linewidth - 2\tabcolsep) * \real{0.6296}}@{}}
\toprule\noalign{}
\begin{minipage}[b]{\linewidth}\raggedright
Concepto
\end{minipage} & \begin{minipage}[b]{\linewidth}\raggedright
Valor adoptado
\end{minipage} \\
\midrule\noalign{}
\endhead
\bottomrule\noalign{}
\endlastfoot
Tipo de cambio EUR/USD & 1 EUR = 1,08 USD \\
Vida útil del sensor & 6 años \\
Tasa de descuento (NPV simple) & 5 \% \\
Horas-hombre instalación & 0,3 h/sensor (cuadrilla de 2 personas,
\textasciitilde35 sensores/noche) \\
Coste hora-cuadrilla & 60 €/h \\
Coste medio SIM NB-IoT M2M (\textgreater10 000 unidades) & 0,60 €/mes \\
Coste medio gateway edge & 700 € unidad + 100 € instalación \\
Coste cámara ANPR con compute edge & 900 € \\
\end{longtable}
}

\section{10.2 CAPEX -- piloto (500
plazas)}\label{capex-piloto-500-plazas}

{\def\LTcaptype{none} % do not increment counter
\begin{longtable}[]{@{}
  >{\raggedright\arraybackslash}p{(\linewidth - 6\tabcolsep) * \real{0.1923}}
  >{\raggedright\arraybackslash}p{(\linewidth - 6\tabcolsep) * \real{0.1923}}
  >{\raggedright\arraybackslash}p{(\linewidth - 6\tabcolsep) * \real{0.4038}}
  >{\raggedright\arraybackslash}p{(\linewidth - 6\tabcolsep) * \real{0.2115}}@{}}
\toprule\noalign{}
\begin{minipage}[b]{\linewidth}\raggedright
Concepto
\end{minipage} & \begin{minipage}[b]{\linewidth}\raggedright
Unidades
\end{minipage} & \begin{minipage}[b]{\linewidth}\raggedright
Coste unitario (€)
\end{minipage} & \begin{minipage}[b]{\linewidth}\raggedright
Total (€)
\end{minipage} \\
\midrule\noalign{}
\endhead
\bottomrule\noalign{}
\endlastfoot
Sensores magnéticos AMR & 500 & 75 & 37 500 \\
Instalación sensores (incluye corte de tráfico, asfaltado) & 500 & 40 &
20 000 \\
Gateways edge (3 zonas + redundancia) & 4 & 800 & 3 200 \\
Cámaras ANPR puntuales & 8 & 900 & 7 200 \\
Mástiles / soportes municipales & 8 & 250 & 2 000 \\
Nodos ambientales (1 cada \textasciitilde150 plazas) & 4 & 350 & 1
400 \\
Sistema central de respaldo & 1 & 1 500 & 1 500 \\
Subtotal hardware & & & 72 800 \\
Ingeniería, integración, gestión proyecto (20 \%) & & & 14 560 \\
\textbf{CAPEX piloto} & & & \textbf{87 360 €} \\
\end{longtable}
}

\section{10.3 OPEX mensual -- piloto}\label{opex-mensual-piloto}

{\def\LTcaptype{none} % do not increment counter
\begin{longtable}[]{@{}
  >{\raggedright\arraybackslash}p{(\linewidth - 6\tabcolsep) * \real{0.2273}}
  >{\raggedright\arraybackslash}p{(\linewidth - 6\tabcolsep) * \real{0.2273}}
  >{\raggedright\arraybackslash}p{(\linewidth - 6\tabcolsep) * \real{0.2273}}
  >{\raggedright\arraybackslash}p{(\linewidth - 6\tabcolsep) * \real{0.3182}}@{}}
\toprule\noalign{}
\begin{minipage}[b]{\linewidth}\raggedright
Concepto
\end{minipage} & \begin{minipage}[b]{\linewidth}\raggedright
Cantidad
\end{minipage} & \begin{minipage}[b]{\linewidth}\raggedright
Unitario
\end{minipage} & \begin{minipage}[b]{\linewidth}\raggedright
Mensual (€)
\end{minipage} \\
\midrule\noalign{}
\endhead
\bottomrule\noalign{}
\endlastfoot
SIM NB-IoT por sensor & 500 & 0,60 €/mes & 300 \\
Backhaul gateway/cámara (fibra municipal) & 4 + 8 & incluido en convenio
& 0 \\
AWS IoT Core mensajes (publicación y reglas) & 4,5 M msg/mes & 1 USD / 1
M & 4,2 € \\
AWS Lambda (≈ 5 M invocaciones, 256 MB, 200 ms) & 5 M inv &
\textasciitilde0,30 USD/M req + \textasciitilde0,20 USD GB-s & 0,7 € \\
Amazon DynamoDB (on-demand, \textasciitilde5 M write + 1 M read) & &
1,25 USD/M WRU + 0,25 USD/M RRU & 7,5 € \\
Amazon API Gateway (REST) & 1 M req/mes & 3,5 USD/M req & 3,3 € \\
Amazon S3 (raw events, 3 GB/mes acumulando) & 3 GB & 0,023 USD/GB & 0,07
€ \\
Amazon CloudWatch (logs \textasciitilde5 GB/mes) & 5 GB & 0,50 USD/GB &
2,3 € \\
Operación y mantenimiento (10 h/mes × 35 €/h) & 10 h & 35 €/h & 350 \\
\textbf{OPEX piloto mensual} & & & \textbf{≈ 670 €} \\
\end{longtable}
}

Coste cloud puro (sin OPEX humano ni conectividad): \textbf{≈ 18 €/mes},
ampliamente por debajo de RNF-17.

\section{10.4 CAPEX -- ciudad (10 000
plazas)}\label{capex-ciudad-10-000-plazas}

{\def\LTcaptype{none} % do not increment counter
\begin{longtable}[]{@{}
  >{\raggedright\arraybackslash}p{(\linewidth - 6\tabcolsep) * \real{0.1923}}
  >{\raggedright\arraybackslash}p{(\linewidth - 6\tabcolsep) * \real{0.1923}}
  >{\raggedright\arraybackslash}p{(\linewidth - 6\tabcolsep) * \real{0.4038}}
  >{\raggedright\arraybackslash}p{(\linewidth - 6\tabcolsep) * \real{0.2115}}@{}}
\toprule\noalign{}
\begin{minipage}[b]{\linewidth}\raggedright
Concepto
\end{minipage} & \begin{minipage}[b]{\linewidth}\raggedright
Unidades
\end{minipage} & \begin{minipage}[b]{\linewidth}\raggedright
Coste unitario (€)
\end{minipage} & \begin{minipage}[b]{\linewidth}\raggedright
Total (€)
\end{minipage} \\
\midrule\noalign{}
\endhead
\bottomrule\noalign{}
\endlastfoot
Sensores magnéticos AMR & 10 000 & 75 (con descuento por volumen) & 750
000 \\
Instalación sensores & 10 000 & 35 & 350 000 \\
Gateways edge (40 zonas con redundancia) & 80 & 800 & 64 000 \\
Cámaras ANPR (puntos clave + zonas reguladas) & 50 & 900 & 45 000 \\
Nodos ambientales & 70 & 350 & 24 500 \\
Centro de operaciones (servidor de respaldo, monitor) & 1 & 8 000 & 8
000 \\
Subtotal hardware & & & 1 241 500 \\
Ingeniería + gestión + formación (15 \%) & & & 186 225 \\
\textbf{CAPEX ciudad} & & & \textbf{≈ 1 427 725 €} \\
\end{longtable}
}

\section{10.5 OPEX mensual -- ciudad}\label{opex-mensual-ciudad}

{\def\LTcaptype{none} % do not increment counter
\begin{longtable}[]{@{}
  >{\raggedright\arraybackslash}p{(\linewidth - 6\tabcolsep) * \real{0.2273}}
  >{\raggedright\arraybackslash}p{(\linewidth - 6\tabcolsep) * \real{0.2273}}
  >{\raggedright\arraybackslash}p{(\linewidth - 6\tabcolsep) * \real{0.2273}}
  >{\raggedright\arraybackslash}p{(\linewidth - 6\tabcolsep) * \real{0.3182}}@{}}
\toprule\noalign{}
\begin{minipage}[b]{\linewidth}\raggedright
Concepto
\end{minipage} & \begin{minipage}[b]{\linewidth}\raggedright
Cantidad
\end{minipage} & \begin{minipage}[b]{\linewidth}\raggedright
Unitario
\end{minipage} & \begin{minipage}[b]{\linewidth}\raggedright
Mensual (€)
\end{minipage} \\
\midrule\noalign{}
\endhead
\bottomrule\noalign{}
\endlastfoot
SIM NB-IoT & 10 000 & 0,55 €/mes (volumen) & 5 500 \\
Mantenimiento físico (cuadrillas, repuestos) & --- & --- & 6 000 \\
Conectividad gateways (fibra municipal) & --- & --- & 0 \\
AWS IoT Core mensajes & 90 M msg/mes & 1 USD/M & 83 € \\
AWS Kinesis Data Streams (3 shards, 200 PUT/s pico) & --- &
\textasciitilde75 USD/mes & 70 € \\
Managed Apache Flink (1 KPU) & --- & \textasciitilde110 USD/mes & 102
€ \\
Lambda + DynamoDB + API + S3 + CloudWatch & --- & --- & ≈ 250 € \\
Amazon Timestream (KPIs históricos) & --- & --- & ≈ 90 € \\
WAF + CloudFront & --- & --- & ≈ 30 € \\
Personal operaciones (1 FTE) & --- & --- & 3 500 \\
\textbf{OPEX ciudad mensual} & & & \textbf{≈ 15 625 €} \\
\end{longtable}
}

Coste AWS puro (sin RR. HH. ni SIMs ni operación física): \textbf{≈ 625
€/mes para 10 000 plazas}, es decir, \textbf{\textasciitilde6 céntimos
por plaza y mes} en infraestructura cloud.

\section{10.6 Coste de propiedad total (TCO) a 5
años}\label{coste-de-propiedad-total-tco-a-5-auxf1os}

{\def\LTcaptype{none} % do not increment counter
\begin{longtable}[]{@{}
  >{\raggedright\arraybackslash}p{(\linewidth - 8\tabcolsep) * \real{0.1467}}
  >{\raggedright\arraybackslash}p{(\linewidth - 8\tabcolsep) * \real{0.1467}}
  >{\raggedright\arraybackslash}p{(\linewidth - 8\tabcolsep) * \real{0.2267}}
  >{\raggedright\arraybackslash}p{(\linewidth - 8\tabcolsep) * \real{0.2133}}
  >{\raggedright\arraybackslash}p{(\linewidth - 8\tabcolsep) * \real{0.2667}}@{}}
\toprule\noalign{}
\begin{minipage}[b]{\linewidth}\raggedright
Escenario
\end{minipage} & \begin{minipage}[b]{\linewidth}\raggedright
CAPEX (€)
\end{minipage} & \begin{minipage}[b]{\linewidth}\raggedright
OPEX 5 años (€)
\end{minipage} & \begin{minipage}[b]{\linewidth}\raggedright
TCO total (€)
\end{minipage} & \begin{minipage}[b]{\linewidth}\raggedright
TCO por plaza-año
\end{minipage} \\
\midrule\noalign{}
\endhead
\bottomrule\noalign{}
\endlastfoot
Piloto (500 plazas) & 87 360 & 670 × 60 = 40 200 & 127 560 & \textbf{51
€/plaza/año} \\
Ciudad (10 000 plazas) & 1 427 725 & 15 625 × 60 = 937 500 & 2 365 225 &
\textbf{47 €/plaza/año} \\
\end{longtable}
}

Comparativa cualitativa:

\begin{itemize}
\tightlist
\item
  Una plaza regulada en zona azul típica española factura por encima de
  1 €/día (\textasciitilde365 €/año).
\item
  Solo con incrementar la rotación en un 5 \% o capturar un 2 \%
  adicional en sanciones evitadas, la solución se autofinancia.
\end{itemize}

\section{10.7 Notas adicionales}\label{notas-adicionales}

\begin{itemize}
\tightlist
\item
  Las tarifas de AWS usadas son las publicadas en
  \texttt{aws.amazon.com/pricing} para \texttt{us-east-1}. En
  \texttt{eu-west-1} (región natural para Albacete) los precios son
  \textasciitilde10 \% superiores; los importes se mantienen en el mismo
  orden de magnitud.
\item
  No se incluye el coste de licencias FIWARE (open source) ni de
  cualquier despliegue propio adicional.
\item
  La estimación NO incluye IVA.
\item
  La estimación de coste de operación humana (≈ 1 FTE en ciudad) puede
  absorberse por la propia plantilla del ayuntamiento si la operación se
  hace en el centro municipal de gestión de movilidad.
\end{itemize}

\section{10.8 Conclusiones del análisis de
coste}\label{conclusiones-del-anuxe1lisis-de-coste}

\begin{enumerate}
\def\labelenumi{\arabic{enumi}.}
\tightlist
\item
  La fracción cloud del coste es \textbf{marginal frente al hardware y
  la operación física}.
\item
  El coste cloud se mantiene aproximadamente lineal con el número de
  plazas; el verdadero coste crece con la flota física.
\item
  NB-IoT como red elegida es decisiva en el OPEX: cualquier alternativa
  con cuota mensual \textgreater2 €/SIM duplicaría el coste operativo.
\item
  La arquitectura serverless evita inversión en infraestructura cloud
  propia (sin EC2/RDS) y permite un piloto que cabe en el presupuesto
  típico de un proyecto de innovación de un ayuntamiento mediano.
\end{enumerate}

\newpage

\chapter{11. Prototipo funcional: ejecución, verificación y
capturas}\label{prototipo-funcional-ejecuciuxf3n-verificaciuxf3n-y-capturas}

Este capítulo describe paso a paso cómo se ha construido, desplegado y
verificado el prototipo entregable. Incluye la cronología de comandos
efectivamente ejecutados, los resultados observados y la lista de
capturas de pantalla que acompañan el documento.

\section{11.1 Resumen del prototipo}\label{resumen-del-prototipo}

El prototipo demuestra el flujo extremo a extremo:

\begin{verbatim}
Simulador Python (40 sensores MQTT/TLS)
  -> AWS IoT Core (Things, certificados, policy, Topic Rule)
    -> Lambda ingest -> DynamoDB parking-state
      -> Lambda aggregator -> DynamoDB zone-kpis
    -> API Gateway REST (+ formato GeoJSON)
      -> Streamlit dashboard (mapa, KPIs, serie temporal)
      -> Cualquier tercero por HTTPS
\end{verbatim}

Todo está implementado y desplegado realmente en una cuenta AWS Academy
Learner Lab (\texttt{Account\ 583916379944}, \texttt{us-east-1}).

\section{11.2 Estructura del prototipo en
disco}\label{estructura-del-prototipo-en-disco}

\begin{verbatim}
prototipo/
├── README.md
├── requirements.txt
├── .env.example
├── infra/
│   ├── common.py
│   ├── 01_setup_iot_core.py
│   ├── 02_setup_dynamodb.py
│   ├── 03_setup_lambda.py
│   ├── 04_setup_api_gateway.py
│   ├── 99_teardown.py
│   ├── parking_zone_seed.json
│   └── infra_state.json            (generado en runtime; no en git)
├── simulator/
│   ├── parking_sensor.py
│   ├── fleet_runner.py
│   └── certs/                      (generado al desplegar IoT Core)
├── lambdas/
│   ├── ingest/handler.py
│   ├── aggregator/handler.py
│   └── api/handler.py
├── api/openapi.yaml
└── dashboard/streamlit_app.py
\end{verbatim}

\section{11.3 Requisitos previos}\label{requisitos-previos}

\begin{itemize}
\tightlist
\item
  Python 3.11 o superior (probado en 3.13.9).
\item
  AWS Academy Learner Lab activo con credenciales en
  \texttt{\%USERPROFILE\%\textbackslash{}.aws\textbackslash{}credentials}.
\item
  Conexión a Internet.
\item
  Sistema operativo Windows / Linux / macOS (probado en Windows 11).
\end{itemize}

\section{11.4 Cómo lanzar el prototipo paso a
paso}\label{cuxf3mo-lanzar-el-prototipo-paso-a-paso}

\subsection{11.4.1 Instalación de
dependencias}\label{instalaciuxf3n-de-dependencias}

\begin{Shaded}
\begin{Highlighting}[]
\FunctionTok{cd} \StringTok{"D:\textbackslash{}DISCO DURO PORTABLE\textbackslash{}...\textbackslash{}proyectofinal\textbackslash{}prototipo"}
\NormalTok{python }\OperatorTok{{-}}\NormalTok{m pip install }\OperatorTok{{-}}\NormalTok{r requirements}\OperatorTok{.}\FunctionTok{txt}
\FunctionTok{Copy{-}Item} \OperatorTok{.}\FunctionTok{env}\OperatorTok{.}\FunctionTok{example} \OperatorTok{.}\FunctionTok{env}
\end{Highlighting}
\end{Shaded}

\subsection{11.4.2 Validación rápida de
credenciales}\label{validaciuxf3n-ruxe1pida-de-credenciales}

\begin{Shaded}
\begin{Highlighting}[]
\NormalTok{python }\OperatorTok{{-}}\NormalTok{c }\StringTok{"import boto3; print(boto3.client(\textquotesingle{}sts\textquotesingle{}, region\_name=\textquotesingle{}us{-}east{-}1\textquotesingle{}).get\_caller\_identity())"}
\end{Highlighting}
\end{Shaded}

Debe devolver un objeto con \texttt{Account}, \texttt{Arn} (terminado en
\texttt{:assumed-role/voclabs/...}).

\subsection{11.4.3 Despliegue de la infraestructura
AWS}\label{despliegue-de-la-infraestructura-aws}

Ejecutar en orden:

\begin{Shaded}
\begin{Highlighting}[]
\NormalTok{python infra\textbackslash{}01\_setup\_iot\_core}\OperatorTok{.}\FunctionTok{py}
\NormalTok{python infra\textbackslash{}02\_setup\_dynamodb}\OperatorTok{.}\FunctionTok{py}
\NormalTok{python infra\textbackslash{}03\_setup\_lambda}\OperatorTok{.}\FunctionTok{py}
\NormalTok{python infra\textbackslash{}04\_setup\_api\_gateway}\OperatorTok{.}\FunctionTok{py}
\end{Highlighting}
\end{Shaded}

Tras \texttt{01\_*} se generan los ficheros
\texttt{simulator/certs/\{device.cert.pem,\ device.private.key,\ AmazonRootCA1.pem\}}.
Tras \texttt{04\_*} el archivo \texttt{infra/infra\_state.json} contiene
la URL pública de la API (\texttt{apiBaseUrl}).

\subsection{11.4.4 Lanzamiento del simulador de
telemetría}\label{lanzamiento-del-simulador-de-telemetruxeda}

\begin{Shaded}
\begin{Highlighting}[]
\NormalTok{python simulator\textbackslash{}fleet\_runner}\OperatorTok{.}\FunctionTok{py} \OperatorTok{{-}{-}}\NormalTok{num{-}spots }\DecValTok{40} \OperatorTok{{-}{-}}\NormalTok{duration }\DecValTok{180} \OperatorTok{{-}{-}}\NormalTok{heartbeat }\DecValTok{20} \OperatorTok{{-}{-}}\NormalTok{tick }\DecValTok{2}
\end{Highlighting}
\end{Shaded}

Parámetros:

\begin{itemize}
\tightlist
\item
  \texttt{-\/-num-spots}: número de plazas a simular (máx. 40 según el
  seed).
\item
  \texttt{-\/-duration}: segundos a ejecutar (0 = indefinido).
\item
  \texttt{-\/-heartbeat}: segundos entre heartbeats sin cambio de
  estado.
\item
  \texttt{-\/-tick}: período del bucle interno del simulador.
\end{itemize}

\subsection{11.4.5 Lanzamiento del
dashboard}\label{lanzamiento-del-dashboard}

En otra consola:

\begin{Shaded}
\begin{Highlighting}[]
\NormalTok{python }\OperatorTok{{-}}\NormalTok{m streamlit run dashboard\textbackslash{}streamlit\_app}\OperatorTok{.}\FunctionTok{py}
\end{Highlighting}
\end{Shaded}

Se abre el navegador en \texttt{http://localhost:8501}. La aplicación
toma la URL de la API de \texttt{infra/infra\_state.json}. El refresco
automático está configurable en la barra lateral.

\subsection{11.4.6 Teardown obligatorio}\label{teardown-obligatorio}

Para no dejar recursos colgando (y proteger las horas de lab restantes):

\begin{Shaded}
\begin{Highlighting}[]
\NormalTok{python infra\textbackslash{}99\_teardown}\OperatorTok{.}\FunctionTok{py}
\end{Highlighting}
\end{Shaded}

\section{11.5 Cómo utilizar el
sistema}\label{cuxf3mo-utilizar-el-sistema}

\subsection{11.5.1 Como operador municipal
(dashboard)}\label{como-operador-municipal-dashboard}

\begin{enumerate}
\def\labelenumi{\arabic{enumi}.}
\tightlist
\item
  Abrir \texttt{http://localhost:8501}.
\item
  La cabecera muestra los KPIs globales (total / libres / ocupadas / sin
  datos / \% ocupación).
\item
  El mapa central muestra cada plaza con color (verde = libre, rojo =
  ocupada, gris = sin datos).
\item
  La sección ``KPIs por sub-zona'' muestra una tabla y un gráfico de
  barras por sub-zona.
\item
  En ``Evolución temporal'' se selecciona una sub-zona y se ve la serie
  temporal del \% de ocupación.
\item
  La tabla ``Detalle de plazas'' permite filtrar e inspeccionar.
\end{enumerate}

\subsection{11.5.2 Como sistema externo (vehículo autónomo,
app)}\label{como-sistema-externo-vehuxedculo-autuxf3nomo-app}

\begin{Shaded}
\begin{Highlighting}[]
\CommentTok{\# Estado global}
\ExtensionTok{curl}\NormalTok{ https://}\OperatorTok{\textless{}}\NormalTok{api{-}id}\OperatorTok{\textgreater{}}\NormalTok{.execute{-}api.us{-}east{-}1.amazonaws.com/prod/spots}

\CommentTok{\# Plazas libres en una sub{-}zona}
\ExtensionTok{curl} \StringTok{\textquotesingle{}https://\textless{}api{-}id\textgreater{}.execute{-}api.us{-}east{-}1.amazonaws.com/prod/spots?zone=Z2{-}DEPORTIVO\textquotesingle{}}

\CommentTok{\# GeoJSON para un mapa cliente}
\ExtensionTok{curl} \StringTok{\textquotesingle{}https://\textless{}api{-}id\textgreater{}.execute{-}api.us{-}east{-}1.amazonaws.com/prod/spots?format=geojson\textquotesingle{}}

\CommentTok{\# Detalle de una plaza}
\ExtensionTok{curl}\NormalTok{ https://}\OperatorTok{\textless{}}\NormalTok{api{-}id}\OperatorTok{\textgreater{}}\NormalTok{.execute{-}api.us{-}east{-}1.amazonaws.com/prod/spots/ALB{-}Z1{-}001}

\CommentTok{\# Serie temporal KPIs}
\ExtensionTok{curl}\NormalTok{ https://}\OperatorTok{\textless{}}\NormalTok{api{-}id}\OperatorTok{\textgreater{}}\NormalTok{.execute{-}api.us{-}east{-}1.amazonaws.com/prod/zones/Z1{-}CAMPUS/kpis}\PreprocessorTok{?}\NormalTok{limit=20}
\end{Highlighting}
\end{Shaded}

\section{11.6 Cómo visualizar los
resultados}\label{cuxf3mo-visualizar-los-resultados}

{\def\LTcaptype{none} % do not increment counter
\begin{longtable}[]{@{}
  >{\raggedright\arraybackslash}p{(\linewidth - 4\tabcolsep) * \real{0.3500}}
  >{\raggedright\arraybackslash}p{(\linewidth - 4\tabcolsep) * \real{0.3500}}
  >{\raggedright\arraybackslash}p{(\linewidth - 4\tabcolsep) * \real{0.3000}}@{}}
\toprule\noalign{}
\begin{minipage}[b]{\linewidth}\raggedright
Vista
\end{minipage} & \begin{minipage}[b]{\linewidth}\raggedright
Dónde
\end{minipage} & \begin{minipage}[b]{\linewidth}\raggedright
Cómo
\end{minipage} \\
\midrule\noalign{}
\endhead
\bottomrule\noalign{}
\endlastfoot
Estado por plaza & Dashboard Streamlit + \texttt{GET\ /spots} & Tabla y
mapa \\
KPIs por zona en vivo & Dashboard + \texttt{GET\ /zones} & Tabla +
gráfico de barras \\
Evolución temporal & Dashboard + \texttt{GET\ /zones/\{id\}/kpis} &
Plotly line \\
Mensajes MQTT crudos & AWS Console \textgreater{} IoT \textgreater{}
Test \textgreater{} MQTT test client, suscribiéndose a
\texttt{parking/\#} & Útil para defensa \\
Items en DynamoDB & AWS Console \textgreater{} DynamoDB \textgreater{}
Tables \textgreater{} \texttt{smart-parking-albacete-state}
\textgreater{} Explore & Ver datos como se persisten \\
Logs de Lambda & AWS Console \textgreater{} CloudWatch \textgreater{}
Log groups \textgreater{}
\texttt{/aws/lambda/smart-parking-albacete-ingest} & Diagnóstico \\
\end{longtable}
}

\section{11.7 Cómo explicar el prototipo en la
defensa}\label{cuxf3mo-explicar-el-prototipo-en-la-defensa}

Discurso recomendado (≈ 5 minutos):

\begin{enumerate}
\def\labelenumi{\arabic{enumi}.}
\tightlist
\item
  \textbf{Contexto} (30 s): smart parking en zona universitaria de
  Albacete; cliente ficticio TECO S.L.; problema real de tráfico
  ineficiente buscando aparcamiento.
\item
  \textbf{Decisiones clave} (1 min): sensor magnético + NB-IoT + AWS
  serverless. Por qué cada elección (referencias al capítulo 3 y 4 de la
  memoria).
\item
  \textbf{Arquitectura} (1 min): diagrama mermaid del capítulo 5;
  remarcar que es serverless gestionada y multi-AZ.
\item
  \textbf{Demostración en vivo} (2 min):

  \begin{itemize}
  \tightlist
  \item
    Mostrar el dashboard ya cargado con el mapa.
  \item
    Lanzar el simulador en una consola y mostrar cómo, al cabo de unos
    segundos, las plazas cambian de estado en el dashboard y la serie
    temporal sube.
  \item
    Mostrar la API con un \texttt{curl} al GeoJSON (la respuesta puede
    pintarse en \texttt{geojson.io} si hay conexión).
  \item
    Mostrar en la consola AWS \textgreater{} IoT \textgreater{} MQTT
    test client suscrito a \texttt{parking/\#} que recibe los mensajes
    en tiempo real.
  \end{itemize}
\item
  \textbf{Escalado y coste} (30 s): de 500 a 10 000 plazas con los
  cambios del capítulo 9; coste cloud despreciable (\textasciitilde6
  céntimos/plaza/mes).
\item
  \textbf{Pregunta esperada --- ``¿Por qué no FIWARE?''} (15 s): se ha
  diseñado la integración híbrida; en piloto sobra; en ciudad se
  introduce sin tocar el resto.
\end{enumerate}

\section{11.8 Verificación end-to-end (resultado real
obtenido)}\label{verificaciuxf3n-end-to-end-resultado-real-obtenido}

{\def\LTcaptype{none} % do not increment counter
\begin{longtable}[]{@{}
  >{\raggedright\arraybackslash}p{(\linewidth - 4\tabcolsep) * \real{0.1622}}
  >{\raggedright\arraybackslash}p{(\linewidth - 4\tabcolsep) * \real{0.2432}}
  >{\raggedright\arraybackslash}p{(\linewidth - 4\tabcolsep) * \real{0.5946}}@{}}
\toprule\noalign{}
\begin{minipage}[b]{\linewidth}\raggedright
Paso
\end{minipage} & \begin{minipage}[b]{\linewidth}\raggedright
Comando
\end{minipage} & \begin{minipage}[b]{\linewidth}\raggedright
Resultado observado
\end{minipage} \\
\midrule\noalign{}
\endhead
\bottomrule\noalign{}
\endlastfoot
Validar credenciales & \texttt{python\ -c\ "import\ boto3;\ ..."} &
\texttt{Account=583916379944,\ Arn=...voclabs/...} \\
Crear IoT Core & \texttt{python\ infra/01\_setup\_iot\_core.py} & 40
Things, certificado, policy, attach OK \\
Crear DynamoDB & \texttt{python\ infra/02\_setup\_dynamodb.py} & Dos
tablas en estado ACTIVE \\
Crear Lambdas & \texttt{python\ infra/03\_setup\_lambda.py} & 3 Lambdas
creadas, IoT Rule activa \\
Crear API Gateway & \texttt{python\ infra/04\_setup\_api\_gateway.py} &
URL:
\texttt{https://85fbp0svzc.execute-api.us-east-1.amazonaws.com/prod} \\
Ejecutar simulador (60 s) &
\texttt{python\ simulator/fleet\_runner.py\ -\/-duration\ 60\ ...} & 40
sensores conectados, eventos publicados \\
\texttt{GET\ /spots} & \texttt{curl\ .../prod/spots} & 200 OK, 39 items
(la 40ª es sensor en fallo intencional) \\
\texttt{GET\ /zones} & \texttt{curl\ .../prod/zones} & 4 sub-zonas con
KPIs (Z1 40 \%, Z2 60 \%, Z3 40 \%, Z4 22 \%) \\
\texttt{GET\ /spots?format=geojson} & \texttt{curl\ ...\&format=geojson}
& FeatureCollection con 10 features para Z1 \\
\texttt{GET\ /zones/Z1-CAMPUS/kpis} &
\texttt{curl\ .../zones/Z1-CAMPUS/kpis?limit=5} & 5 entradas con
\texttt{windowEnd} y serie temporal \\
\end{longtable}
}

\section{11.9 Capturas de pantalla a incluir en la memoria (a generar
por el
autor)}\label{capturas-de-pantalla-a-incluir-en-la-memoria-a-generar-por-el-autor}

\begin{quote}
\textbf{Importante}: el agente no puede tomar capturas de pantalla del
sistema operativo del autor. Los siguientes assets deben generarse
manualmente antes de la entrega final y guardarse en
\texttt{memoria/diagramas/}. Se sugiere nomenclatura para que el
documento las localice automáticamente:
\end{quote}

\begin{enumerate}
\def\labelenumi{\arabic{enumi}.}
\tightlist
\item
  \textbf{\texttt{captura\_aws\_iot\_things.png}} -- Consola AWS
  \textgreater{} IoT Core \textgreater{} Manage \textgreater{} Things,
  mostrando los 40 Things \texttt{ALB-Z1-001} \ldots{}
  \texttt{ALB-Z4-010} listados.
\item
  \textbf{\texttt{captura\_aws\_iot\_mqtt\_test.png}} -- Consola AWS
  \textgreater{} IoT Core \textgreater{} Test \textgreater{} MQTT test
  client, suscrito a \texttt{parking/\#}, mostrando mensajes recientes
  con el JSON publicado por el simulador.
\item
  \textbf{\texttt{captura\_aws\_iot\_rule.png}} -- Consola AWS
  \textgreater{} IoT Core \textgreater{} Message routing \textgreater{}
  Rules \textgreater{} \texttt{smart\_parking\_albacete\_ingest\_rule},
  mostrando el SQL y la acción Lambda asociada.
\item
  \textbf{\texttt{captura\_aws\_dynamodb\_state.png}} -- Consola AWS
  \textgreater{} DynamoDB \textgreater{} Tables \textgreater{}
  \texttt{smart-parking-albacete-state} \textgreater{} Explore,
  mostrando una decena de items con sus \texttt{spotId},
  \texttt{status}, \texttt{lat}, \texttt{lon}.
\item
  \textbf{\texttt{captura\_aws\_dynamodb\_kpis.png}} -- Ídem para
  \texttt{smart-parking-albacete-zone-kpis}, mostrando varias entradas
  con \texttt{windowEnd} distintos.
\item
  \textbf{\texttt{captura\_aws\_lambda\_ingest.png}} -- Consola AWS
  \textgreater{} Lambda \textgreater{}
  \texttt{smart-parking-albacete-ingest}, mostrando la configuración
  (runtime Python 3.12, rol LabRole) y un test de invocación.
\item
  \textbf{\texttt{captura\_aws\_cloudwatch\_ingest.png}} -- Consola AWS
  \textgreater{} CloudWatch \textgreater{} Log groups \textgreater{}
  \texttt{/aws/lambda/smart-parking-albacete-ingest}, mostrando un log
  stream con varias invocaciones.
\item
  \textbf{\texttt{captura\_aws\_apigateway.png}} -- Consola AWS
  \textgreater{} API Gateway \textgreater{} APIs \textgreater{}
  \texttt{smart-parking-albacete-api}, mostrando los recursos y el stage
  \texttt{prod} desplegado.
\item
  \textbf{\texttt{captura\_dashboard\_principal.png}} -- Streamlit con
  KPIs cabecera, mapa de plazas y leyenda.
\item
  \textbf{\texttt{captura\_dashboard\_kpis\_zona.png}} -- Sección ``KPIs
  por sub-zona'' con tabla y gráfico de barras.
\item
  \textbf{\texttt{captura\_dashboard\_serie\_temporal.png}} -- Sección
  de evolución temporal mostrando una serie con varios puntos.
\item
  \textbf{\texttt{captura\_curl\_geojson.png}} -- Consola PowerShell con
  la respuesta GeoJSON formateada.
\item
  \textbf{\texttt{captura\_curl\_zones.png}} -- Consola PowerShell con
  la respuesta de \texttt{GET\ /zones}.
\end{enumerate}

Los diagramas Mermaid de los capítulos 5, 6 y 9 se renderizan
automáticamente al exportar a PDF a través de pandoc (siempre que la
cadena de plantillas tenga soporte). Si no, basta con renderizarlos en
\texttt{mermaid.live} y guardar el PNG con el mismo nombre que el
bloque.

\section{11.10 Mapa ``tal cual'' del sistema en
operación}\label{mapa-tal-cual-del-sistema-en-operaciuxf3n}

Para que el lector se haga una idea inmediata del aspecto del prototipo:

\begin{itemize}
\tightlist
\item
  40 puntos coloreados sobre el plano de la zona universitaria de
  Albacete.
\item
  4 grupos por sub-zona.
\item
  Barra lateral con configuración y enlace a la API.
\item
  Refresco visible cada 5-10 segundos.
\item
  Latencia entre cambio de estado (en el simulador) y refresco en
  pantalla: \textbf{2-3 segundos} en operación normal.
\end{itemize}

\newpage

\chapter{12. Limitaciones y trabajo
futuro}\label{limitaciones-y-trabajo-futuro}

Ningún piloto es la versión final. Este capítulo enumera honestamente
las limitaciones de la solución entregada y propone una hoja de ruta
para los siguientes pasos. Esta sinceridad es importante de cara a la
defensa: muestra capacidad crítica y demuestra que el autor distingue lo
que está demostrado de lo que está pendiente.

\section{12.1 Limitaciones derivadas del entorno
académico}\label{limitaciones-derivadas-del-entorno-acaduxe9mico}

{\def\LTcaptype{none} % do not increment counter
\begin{longtable}[]{@{}
  >{\raggedright\arraybackslash}p{(\linewidth - 4\tabcolsep) * \real{0.3000}}
  >{\raggedright\arraybackslash}p{(\linewidth - 4\tabcolsep) * \real{0.2250}}
  >{\raggedright\arraybackslash}p{(\linewidth - 4\tabcolsep) * \real{0.4750}}@{}}
\toprule\noalign{}
\begin{minipage}[b]{\linewidth}\raggedright
Limitación
\end{minipage} & \begin{minipage}[b]{\linewidth}\raggedright
Impacto
\end{minipage} & \begin{minipage}[b]{\linewidth}\raggedright
Mitigación / Plan
\end{minipage} \\
\midrule\noalign{}
\endhead
\bottomrule\noalign{}
\endlastfoot
Credenciales AWS Academy caducan a las 3 h & El despliegue debe
completarse y demostrarse en una sesión & Scripts idempotentes
(\texttt{infra/*.py}) y \texttt{99\_teardown.py}. \\
Solo \texttt{LabRole} para Lambda & No se pueden afinar permisos por
servicio & Documentado; en producción se crea un rol por Lambda. \\
Servicios no disponibles en el lab & Amazon Timestream y Kinesis pueden
no estar habilitados & Diseño documentado; no impacta el prototipo
actual. \\
Cuenta compartida con compañeros & Posible colisión de nombres & Prefijo
\texttt{smart-parking-albacete-} único. \\
\end{longtable}
}

\section{12.2 Limitaciones técnicas del
prototipo}\label{limitaciones-tuxe9cnicas-del-prototipo}

{\def\LTcaptype{none} % do not increment counter
\begin{longtable}[]{@{}
  >{\raggedright\arraybackslash}p{(\linewidth - 4\tabcolsep) * \real{0.3243}}
  >{\raggedright\arraybackslash}p{(\linewidth - 4\tabcolsep) * \real{0.2432}}
  >{\raggedright\arraybackslash}p{(\linewidth - 4\tabcolsep) * \real{0.4324}}@{}}
\toprule\noalign{}
\begin{minipage}[b]{\linewidth}\raggedright
Limitación
\end{minipage} & \begin{minipage}[b]{\linewidth}\raggedright
Impacto
\end{minipage} & \begin{minipage}[b]{\linewidth}\raggedright
Plan de futuro
\end{minipage} \\
\midrule\noalign{}
\endhead
\bottomrule\noalign{}
\endlastfoot
Certificado X.509 compartido por toda la flota & Revocación granular
imposible & AWS IoT Fleet Provisioning Templates por dispositivo. \\
\texttt{Scan} filtrado en DynamoDB & Coste creciente a partir de
\textasciitilde10 000 plazas & Crear GSI
\texttt{zoneId-status-index}. \\
Agregador en Lambda síncrona & Latencia añadida en eventos
correlacionados & Sustituir por Flink en escenario ciudad. \\
Sin alertas automáticas de sensor caído & Operador debe revisar
manualmente & Job EventBridge cron + Lambda que audita
\texttt{lastUpdated}. \\
Sin retentiva configurable en KPIs & Crece la tabla con el tiempo &
Habilitar DynamoDB TTL o migrar a Timestream. \\
Sin autenticación en la API & Cualquiera puede consultar & API Keys +
WAF + Cognito (diseñado, no implementado en piloto). \\
Dashboard local & No multi-usuario, sin permisos & Streamlit Cloud / ECS
+ Cognito o sustituir por QuickSight. \\
Una sola región (\texttt{us-east-1}) & Sin DR multi-región & Diseño con
DynamoDB Global Tables; redespliegue por CDK. \\
Simulador en lugar de sensores reales & Comportamiento idealizado &
Validar con sensores piloto (fase 1 del plan capítulo 4). \\
\end{longtable}
}

\section{12.3 Limitaciones funcionales}\label{limitaciones-funcionales}

{\def\LTcaptype{none} % do not increment counter
\begin{longtable}[]{@{}
  >{\raggedright\arraybackslash}p{(\linewidth - 4\tabcolsep) * \real{0.3243}}
  >{\raggedright\arraybackslash}p{(\linewidth - 4\tabcolsep) * \real{0.2432}}
  >{\raggedright\arraybackslash}p{(\linewidth - 4\tabcolsep) * \real{0.4324}}@{}}
\toprule\noalign{}
\begin{minipage}[b]{\linewidth}\raggedright
Limitación
\end{minipage} & \begin{minipage}[b]{\linewidth}\raggedright
Por qué
\end{minipage} & \begin{minipage}[b]{\linewidth}\raggedright
Trabajo futuro
\end{minipage} \\
\midrule\noalign{}
\endhead
\bottomrule\noalign{}
\endlastfoot
Sin estimación de duración de la ocupación & El agregador no calcula
\texttt{dwell\ time} & Session windows en Flink. \\
Sin predicción a futuro & El piloto solo refleja el estado actual &
Modelo ML batch en SageMaker; integrar predicciones en la API. \\
Sin gestión de plazas reservadas (PMR, carga eléctrica, residentes) &
Modelo simplificado & Añadir atributo \texttt{restrictedTo} y reglas de
visibilidad por consumidor. \\
Sin integración con sistema de cobro & Fuera de alcance & Integración
con plataforma de tarificación municipal. \\
Sin panel del ciudadano & El dashboard es de operador & App ciudadana
con tiles libres por zona y guiado por GPS. \\
Sin paneles de mensajería variable & Fuera de alcance & Integración con
DMS municipales por API. \\
\end{longtable}
}

\section{12.4 Roadmap propuesto (12
meses)}\label{roadmap-propuesto-12-meses}

{\def\LTcaptype{none} % do not increment counter
\begin{longtable}[]{@{}
  >{\raggedright\arraybackslash}p{(\linewidth - 2\tabcolsep) * \real{0.4000}}
  >{\raggedright\arraybackslash}p{(\linewidth - 2\tabcolsep) * \real{0.6000}}@{}}
\toprule\noalign{}
\begin{minipage}[b]{\linewidth}\raggedright
Mes
\end{minipage} & \begin{minipage}[b]{\linewidth}\raggedright
Hito
\end{minipage} \\
\midrule\noalign{}
\endhead
\bottomrule\noalign{}
\endlastfoot
1 & Piloto técnico con 20 sensores reales y validación del flujo
end-to-end con tráfico real. \\
2-3 & Despliegue de los 500 sensores piloto; calibración de patrones
reales por zona. \\
4 & Auditoría de seguridad (pentest) y endurecimiento (Cognito, WAF,
Device Defender). \\
5 & Integración FIWARE Orion para interoperabilidad con la plataforma
municipal. \\
6 & Migración del histórico a Timestream + cuadros de mando en
QuickSight. \\
7 & App móvil ciudadana (PWA) con localización de plazas libres más
cercanas. \\
8 & Introducción de Flink para KPIs por ventanas y detección de
patrones. \\
9 & Integración con paneles de mensajería variable municipales. \\
10-12 & Extensión a 10 000 plazas; entrada en producción ``ciudad''. \\
\end{longtable}
}

\section{12.5 Investigación y mejoras
prospectivas}\label{investigaciuxf3n-y-mejoras-prospectivas}

\begin{itemize}
\tightlist
\item
  \textbf{Modelos de predicción de ocupación}: redes neuronales
  recurrentes (LSTM) o transformers entrenados con histórico anual y
  variables exógenas (calendario universitario, calendario deportivo,
  festividades, meteorología).
\item
  \textbf{Routing dinámico para vehículos autónomos}: integración del
  API REST con plataformas C-V2X (RSU + MEC) para guiado plaza a plaza.
\item
  \textbf{Carbon awareness}: correlación de ocupación con calidad del
  aire (nodos ambientales) para informar políticas de bajas emisiones.
\item
  \textbf{Plazas reservables}: módulo de reserva temporal para vehículos
  eléctricos en carga o vehículos de emergencia.
\item
  \textbf{Open data}: publicación del histórico anonimizado en el portal
  de datos abiertos del ayuntamiento.
\end{itemize}

\section{12.6 Riesgos abiertos}\label{riesgos-abiertos}

{\def\LTcaptype{none} % do not increment counter
\begin{longtable}[]{@{}
  >{\raggedright\arraybackslash}p{(\linewidth - 2\tabcolsep) * \real{0.2500}}
  >{\raggedright\arraybackslash}p{(\linewidth - 2\tabcolsep) * \real{0.7500}}@{}}
\toprule\noalign{}
\begin{minipage}[b]{\linewidth}\raggedright
Riesgo
\end{minipage} & \begin{minipage}[b]{\linewidth}\raggedright
Mitigación recomendada
\end{minipage} \\
\midrule\noalign{}
\endhead
\bottomrule\noalign{}
\endlastfoot
Cambio del operador NB-IoT con subida de precios & Cláusulas
contractuales multianuales + alternativa LoRaWAN preparada. \\
Pérdida de soporte de un SDK / librería usado & Sin dependencias
propietarias críticas; código modular fácil de migrar. \\
Resistencia ciudadana a las cámaras ANPR & Comunicación clara, política
de retención breve, hash de matrículas. \\
Variabilidad meteorológica afecta a las baterías & Selección de sensores
certificados para -20 a +60 °C; auditoría anual. \\
Vandalismo en una zona concreta & Movilizar cuadrilla de mantenimiento;
reposición rápida (sensor barato). \\
\end{longtable}
}

\newpage

\chapter{13. Preparación de la
defensa}\label{preparaciuxf3n-de-la-defensa}

La guía docente otorga \textbf{20 \% de la nota (TR2)} a la defensa oral
y \textbf{5 \% (PR2)} a la entrevista de autoría. Este capítulo recopila
las preguntas previsibles, las respuestas razonadas y los puntos de
control que el alumno debe dominar antes de la defensa.

\section{13.1 Guion sugerido (10
minutos)}\label{guion-sugerido-10-minutos}

{\def\LTcaptype{none} % do not increment counter
\begin{longtable}[]{@{}
  >{\raggedright\arraybackslash}p{(\linewidth - 4\tabcolsep) * \real{0.2424}}
  >{\raggedright\arraybackslash}p{(\linewidth - 4\tabcolsep) * \real{0.3333}}
  >{\raggedright\arraybackslash}p{(\linewidth - 4\tabcolsep) * \real{0.4242}}@{}}
\toprule\noalign{}
\begin{minipage}[b]{\linewidth}\raggedright
Minuto
\end{minipage} & \begin{minipage}[b]{\linewidth}\raggedright
Contenido
\end{minipage} & \begin{minipage}[b]{\linewidth}\raggedright
Apoyo visual
\end{minipage} \\
\midrule\noalign{}
\endhead
\bottomrule\noalign{}
\endlastfoot
0:00 - 1:00 & Presentación: nombre, asignatura, cliente ficticio,
problema & Slide portada con BBOX y foto satélite zona \\
1:00 - 2:30 & Requisitos funcionales y no funcionales más relevantes &
Slide con tabla resumen capítulo 2 \\
2:30 - 4:00 & Decisión telemetría + conectividad & Slides con
comparativa de los capítulos 3 y 4 \\
4:00 - 6:00 & Arquitectura cloud AWS, flujo extremo a extremo & Diagrama
mermaid del capítulo 5 \\
6:00 - 8:00 & \textbf{Demostración en vivo del prototipo} & Dashboard +
consola AWS IoT + curl \\
8:00 - 9:00 & Escalabilidad y coste & Tablas capítulos 9 y 10 \\
9:00 - 10:00 & Limitaciones, trabajo futuro y conclusiones & Capítulo
12 \\
\end{longtable}
}

\section{13.2 Preguntas previsibles y
respuestas}\label{preguntas-previsibles-y-respuestas}

\subsection{13.2.1 ¿Por qué elegiste sensores magnéticos en lugar de
cámaras?}\label{por-quuxe9-elegiste-sensores-magnuxe9ticos-en-lugar-de-cuxe1maras}

\begin{itemize}
\tightlist
\item
  Privacidad (no captura matrículas).
\item
  Consumo (vida útil 5-7 años con pila D).
\item
  Coste (\textasciitilde75 € frente a 500-1200 € de una cámara que cubre
  5-10 plazas).
\item
  Robustez climática (IP68 bajo asfalto).
\item
  Las cámaras se reservan para los accesos como conteo agregado y
  posible control de zonas de bajas emisiones.
\end{itemize}

\subsection{13.2.2 ¿Qué pasa si falla un sensor o un
gateway?}\label{quuxe9-pasa-si-falla-un-sensor-o-un-gateway}

\begin{itemize}
\tightlist
\item
  Pérdida de señal del sensor: el agregador no recibe heartbeat.
  Mecanismo previsto (trabajo futuro): EventBridge cron + Lambda que
  audita \texttt{lastUpdated} y emite alerta a SNS.
\item
  Pérdida del gateway edge: en NB-IoT directo no hay gateway propio.
  Para sensores LoRaWAN se mitiga con redundancia (un gateway adicional
  por zona).
\item
  Pérdida de conectividad temporal: el sensor mantiene caché local y
  reenvía al recuperar (en el simulador no se ha implementado, sí se
  documenta).
\end{itemize}

\subsection{13.2.3 ¿Cómo evitas reportar falsos
libres/ocupados?}\label{cuxf3mo-evitas-reportar-falsos-libresocupados}

\begin{itemize}
\tightlist
\item
  Debounce en edge: el cambio sólo se considera tras 2-3 s estables.
\item
  Umbral de \texttt{confidence} en el sensor (campo en el payload).
\item
  Calibración inicial in situ del magnetómetro.
\item
  Filtrado adicional en Lambda ingest (descartar payloads malformados).
\end{itemize}

\subsection{13.2.4 ¿Cómo escalas la arquitectura de 500 a 10 000
plazas?}\label{cuxf3mo-escalas-la-arquitectura-de-500-a-10-000-plazas}

Resumen del capítulo 9:

\begin{itemize}
\tightlist
\item
  Sensores: certificado por dispositivo via Fleet Provisioning
  Templates.
\item
  Procesamiento: Kinesis Data Streams + Managed Apache Flink en lugar de
  Lambda agregador.
\item
  Estado: añadir GSI \texttt{zoneId-status-index} en DynamoDB.
\item
  Histórico: migrar \texttt{zone-kpis} a Amazon Timestream.
\item
  API: WAF + CloudFront caché + cuotas.
\item
  Coste cloud: \textasciitilde625 €/mes (6 céntimos por plaza/mes).
\end{itemize}

\subsection{13.2.5 ¿Qué latencia se espera y de dónde
sale?}\label{quuxe9-latencia-se-espera-y-de-duxf3nde-sale}

\begin{itemize}
\tightlist
\item
  NB-IoT: 1-3 s típicos (CP/IDLE -\textgreater{} conexión radio + envío
  IP).
\item
  IoT Core → IoT Rule → Lambda → DynamoDB: 50-300 ms.
\item
  DynamoDB → API Gateway → cliente: 100-500 ms.
\item
  \textbf{Total}: \textless{} 5 s extremo a extremo. Medido en el
  prototipo: \textbf{2-3 s} entre publicación del simulador y refresco
  en el dashboard.
\end{itemize}

\subsection{13.2.6 ¿Qué datos guardas como estado y cuáles como
histórico?}\label{quuxe9-datos-guardas-como-estado-y-cuuxe1les-como-histuxf3rico}

\begin{itemize}
\tightlist
\item
  \textbf{Estado actual} en \texttt{parking-state}: el último valor de
  cada plaza, sobrescrito por cada evento.
\item
  \textbf{Histórico operacional} en \texttt{zone-kpis}: una fila por
  cada cálculo del agregador, indexada por
  \texttt{(zoneId,\ windowEnd)}. Permite reconstruir cómo cambió la
  ocupación en el tiempo.
\item
  \textbf{Raw events} (diseño, no implementado): cada mensaje crudo a S3
  (data lake bronze) para analítica futura con Athena.
\end{itemize}

\subsection{13.2.7 ¿Cómo proteges las
APIs?}\label{cuxf3mo-proteges-las-apis}

\begin{itemize}
\tightlist
\item
  HTTPS obligatorio (TLS 1.2+).
\item
  CORS controlado en Lambda.
\item
  API Gateway permite throttling y API keys (no activados en piloto,
  diseñados para producción).
\item
  WAF y Cognito como ampliaciones en producción.
\item
  IAM principle of least privilege: en piloto se usa \texttt{LabRole},
  en producción un rol por Lambda.
\end{itemize}

\subsection{13.2.8 ¿Qué parte del prototipo es real y qué parte está
simulada?}\label{quuxe9-parte-del-prototipo-es-real-y-quuxe9-parte-estuxe1-simulada}

Real (corriendo en AWS):

\begin{itemize}
\tightlist
\item
  AWS IoT Core con 40 Things, certificado, policy, Topic Rule.
\item
  3 Lambdas (\texttt{ingest}, \texttt{aggregator}, \texttt{api}).
\item
  2 tablas DynamoDB (\texttt{parking-state}, \texttt{zone-kpis}).
\item
  API Gateway REST con stage \texttt{prod}.
\item
  Dashboard Streamlit consumiendo la API real.
\end{itemize}

Simulado:

\begin{itemize}
\tightlist
\item
  Los sensores físicos (sustituidos por un cliente paho-mqtt que habla
  MQTT/TLS contra IoT Core, exactamente igual que un sensor real haría).
\item
  Las cámaras ANPR (solo se documentan).
\item
  El gateway edge (su lógica de debounce está dentro del simulador).
\end{itemize}

\subsection{13.2.9 ¿Por qué decidiste no implementar FIWARE /
Flink?}\label{por-quuxe9-decidiste-no-implementar-fiware-flink}

\begin{itemize}
\tightlist
\item
  AWS Academy Lab dura 3 h: gastarlas en levantar Orion + Flink no
  aporta valor académico, ya que el flujo conceptual queda demostrado
  con la implementación AWS.
\item
  En el piloto, el coste/beneficio de Flink es negativo (Lambda alcanza
  sobradamente).
\item
  Se ha dejado el diseño detallado de la integración para el escenario
  ciudad, donde sí aporta valor (capítulo 6 y 9).
\end{itemize}

\subsection{13.2.10 ¿Qué limitaciones tiene la
solución?}\label{quuxe9-limitaciones-tiene-la-soluciuxf3n}

Resumen del capítulo 12:

\begin{itemize}
\tightlist
\item
  Certificado X.509 compartido (cambiable a por-dispositivo).
\item
  Sin alertas automáticas de sensores caídos (Lambda cron + SNS futuro).
\item
  Sin autenticación de la API pública (API Keys + WAF futuro).
\item
  Sin predicción ML (SageMaker futuro).
\item
  Sin app ciudadana (PWA futura).
\end{itemize}

\section{13.3 Resultados de aprendizaje y dónde se
cubren}\label{resultados-de-aprendizaje-y-duxf3nde-se-cubren}

{\def\LTcaptype{none} % do not increment counter
\begin{longtable}[]{@{}
  >{\raggedright\arraybackslash}p{(\linewidth - 2\tabcolsep) * \real{0.2667}}
  >{\raggedright\arraybackslash}p{(\linewidth - 2\tabcolsep) * \real{0.7333}}@{}}
\toprule\noalign{}
\begin{minipage}[b]{\linewidth}\raggedright
RA
\end{minipage} & \begin{minipage}[b]{\linewidth}\raggedright
Cobertura
\end{minipage} \\
\midrule\noalign{}
\endhead
\bottomrule\noalign{}
\endlastfoot
\textbf{CN02} (arquitecturas de tratamiento masivo) & Capítulos 5, 7,
9. \\
\textbf{HA03} (orquestación ETL / data lakes) & Capítulo 7 (modelo de
datos) + diseño S3 raw bronze (capítulo 9). \\
\textbf{CP02} (IoT, edge, streams) & Capítulos 3 (telemetría), 4
(conectividad), 5 (cloud) y 6 (Flink). \\
\end{longtable}
}

\section{13.4 Checklist final antes de la
entrega}\label{checklist-final-antes-de-la-entrega}

\begin{itemize}
\tightlist
\item[$\square$]
  La memoria (\texttt{MEMORIA\_TECNICA.md} + \texttt{.tex} +
  \texttt{.pdf}) está en \texttt{memoria/}.
\item[$\square$]
  El prototipo es ejecutable de cero siguiendo \texttt{README\_GUIA.md}
  y \texttt{prototipo/README.md}.
\item[$\square$]
  Las capturas de pantalla del capítulo 11 están en
  \texttt{memoria/diagramas/}.
\item[$\square$]
  El \texttt{99\_teardown.py} se ha ejecutado tras la demo (no quedan
  recursos en AWS).
\item[$\square$]
  El ZIP final no contiene credenciales (\texttt{.aws/},
  \texttt{infra/infra\_state.json}, \texttt{simulator/certs/}).
\item[$\square$]
  El nombre del autor consta como \texttt{alonso.marcos@alu.uclm.es}.
\end{itemize}

\section{13.5 Recordatorio operativo de
defensa}\label{recordatorio-operativo-de-defensa}

\begin{itemize}
\tightlist
\item
  Llegar 15 min antes con la sesión de AWS Academy ya iniciada y el
  despliegue corriendo.
\item
  Tener navegador con tres pestañas listas: Streamlit, AWS IoT MQTT test
  client, AWS DynamoDB Explore.
\item
  Tener una pestaña con la memoria en PDF para señalar tablas/diagramas
  si surge la duda.
\item
  Practicar la demo en seco al menos dos veces para que el simulador y
  el dashboard estén sincronizados.
\item
  Estar preparado para responder en castellano y, si el tribunal lo
  pide, mostrar dónde en el código se implementa cada decisión.
\end{itemize}

\end{document}
